% Môn: Vật lý
% Lớp: 12
% Chương: Chương IV. Vật lí hạt nhân
% Bài: L12C4 Bài 25. Bài tập về vật lí hạt nhân
% Số câu: 189
% App đọc file này trực tiếp — sửa rồi Commit trên GitHub, bấm Đồng bộ GitHub .tex

% L12C4 Bài 25. Bài tập về vật lí hạt nhân

% ===== Câu 1 =====
% ID: L12C4B25-01-DS
% Mức: TH
\dangbt{Bài toán vận dụng cao kết hợp nhiều nội dung của Chương IV}
\begin{ex}
% ID: L12C4B25-01-DS
% Mức: TH
Xét bài toán vận dụng cao kết hợp nhiều nội dung của chương iv (tình huống 1). Hãy xác định tính đúng, sai của bốn phát biểu sau.
\choiceTF
{\True Bài vận dụng cao cần tách giai đoạn, dùng đúng bảo toàn, liên hệ khối lượng–năng lượng và định luật phóng xạ, đồng thời kiểm tra đơn vị và ý nghĩa kết quả.}
{Có thể dùng một công thức duy nhất cho mọi bài toán hạt nhân mà không xét điều kiện.}
{\True Khi giải cần đọc đúng dữ kiện, dùng hệ đơn vị thống nhất và kiểm tra điều kiện áp dụng.}
{Có thể bỏ qua dữ kiện, đơn vị và ý nghĩa vật lí của kết quả.}
\loigiai{
\begin{itemchoice}
\item (Đúng) Bài vận dụng cao cần tách giai đoạn, dùng đúng bảo toàn, liên hệ khối lượng–năng lượng và định luật phóng xạ, đồng thời kiểm tra đơn vị và ý nghĩa kết quả. (Vì): Bài vận dụng cao cần tách giai đoạn, dùng đúng bảo toàn, liên hệ khối lượng–năng lượng và định luật phóng xạ, đồng thời kiểm tra đơn vị và ý nghĩa kết quả. b) Sai: Có thể dùng một công thức duy nhất cho mọi bài toán hạt nhân mà không xét điều kiện. c) Đúng: Khi giải cần đọc đúng dữ kiện, dùng hệ đơn vị thống nhất và kiểm tra điều kiện áp dụng. d) Sai: Có thể bỏ qua dữ kiện, đơn vị và ý nghĩa vật lí của kết quả. Kiến thức cốt lõi: Bài vận dụng cao cần tách giai đoạn, dùng đúng bảo toàn, liên hệ khối lượng–năng lượng và định luật phóng xạ, đồng thời kiểm tra đơn vị và ý nghĩa kết quả
\item (Sai) Có thể dùng một công thức duy nhất cho mọi bài toán hạt nhân mà không xét điều kiện.
\item (Đúng) Khi giải cần đọc đúng dữ kiện, dùng hệ đơn vị thống nhất và kiểm tra điều kiện áp dụng.
\item (Sai) Có thể bỏ qua dữ kiện, đơn vị và ý nghĩa vật lí của kết quả.
\end{itemchoice}
}
\end{ex}

% ===== Câu 2 =====
% ID: L12C4B25-01-TL
% Mức: TH
\begin{bt}
% ID: L12C4B25-01-TL
% Mức: TH
Trình bày kiến thức cốt lõi của bài toán vận dụng cao kết hợp nhiều nội dung của chương iv và nêu điều kiện áp dụng.
\loigiai{
Cần nêu được: Bài vận dụng cao cần tách giai đoạn, dùng đúng bảo toàn, liên hệ khối lượng–năng lượng và định luật phóng xạ, đồng thời kiểm tra đơn vị và ý nghĩa kết quả. Khi vận dụng phải xác định đúng dữ kiện, điều kiện áp dụng, hệ đơn vị và kiểm tra tính hợp lí của kết quả. Nhận định “Có thể dùng một công thức duy nhất cho mọi bài toán hạt nhân mà không xét điều kiện.” sai vì trái với kiến thức trên.
}
\end{bt}

% ===== Câu 3 =====
% ID: L12C4B25-01-TLN
% Mức: TH
\begin{ex}
% ID: L12C4B25-01-TLN
% Mức: TH
Trong bốn phát biểu sau về bài toán vận dụng cao kết hợp nhiều nội dung của chương iv, có bao nhiêu phát biểu đúng? Chỉ ghi một số nguyên. (1) Bài vận dụng cao cần tách giai đoạn, dùng đúng bảo toàn, liên hệ khối lượng–năng lượng và định luật phóng xạ, đồng thời kiểm tra đơn vị và ý nghĩa kết quả. (2) Có thể dùng một công thức duy nhất cho mọi bài toán hạt nhân mà không xét điều kiện. (3) Cần kiểm tra điều kiện áp dụng trước khi tính toán. (4) Có thể bỏ qua đơn vị của kết quả.
\shortans{2}
\loigiai{
Đối chiếu từng phát biểu với kiến thức cốt lõi: Bài vận dụng cao cần tách giai đoạn, dùng đúng bảo toàn, liên hệ khối lượng–năng lượng và định luật phóng xạ, đồng thời kiểm tra đơn vị và ý nghĩa kết quả. Có 2 phát biểu đúng.
}
\end{ex}

% ===== Câu 4 =====
% ID: L12C4B25-01-TN
% Mức: NB
\begin{ex}
% ID: L12C4B25-01-TN
% Mức: NB
Phát biểu nào sau đây đúng về bài toán vận dụng cao kết hợp nhiều nội dung của chương iv?
\choice
{\True Bài vận dụng cao cần tách giai đoạn, dùng đúng bảo toàn, liên hệ khối lượng–năng lượng và định luật phóng xạ, đồng thời kiểm tra đơn vị và ý nghĩa kết quả.}
{Có thể dùng một công thức duy nhất cho mọi bài toán hạt nhân mà không xét điều kiện.}
{Có thể bỏ qua điều kiện áp dụng và đơn vị trong mọi trường hợp.}
{Kết quả của bài toán không phụ thuộc dữ kiện đã cho.}
\loigiai{
Bài vận dụng cao cần tách giai đoạn, dùng đúng bảo toàn, liên hệ khối lượng–năng lượng và định luật phóng xạ, đồng thời kiểm tra đơn vị và ý nghĩa kết quả. Vì vậy phương án $A$ đúng; các phương án còn lại trái với kiến thức hoặc điều kiện áp dụng.
}
\end{ex}

% ===== Câu 5 =====
% ID: L12C4B25-02-DS
% Mức: TH
\begin{ex}
% ID: L12C4B25-02-DS
% Mức: TH
Xét bài toán vận dụng cao kết hợp nhiều nội dung của chương iv (tình huống 2). Hãy xác định tính đúng, sai của bốn phát biểu sau.
\choiceTF
{Có thể dùng một công thức duy nhất cho mọi bài toán hạt nhân mà không xét điều kiện.}
{\True Bài vận dụng cao cần tách giai đoạn, dùng đúng bảo toàn, liên hệ khối lượng–năng lượng và định luật phóng xạ, đồng thời kiểm tra đơn vị và ý nghĩa kết quả.}
{Có thể bỏ qua dữ kiện, đơn vị và ý nghĩa vật lí của kết quả.}
{\True Khi giải cần đọc đúng dữ kiện, dùng hệ đơn vị thống nhất và kiểm tra điều kiện áp dụng.}
\loigiai{
\begin{itemchoice}
\item (Sai) Có thể dùng một công thức duy nhất cho mọi bài toán hạt nhân mà không xét điều kiện. (Vì): Có thể dùng một công thức duy nhất cho mọi bài toán hạt nhân mà không xét điều kiện. b) Đúng: Bài vận dụng cao cần tách giai đoạn, dùng đúng bảo toàn, liên hệ khối lượng–năng lượng và định luật phóng xạ, đồng thời kiểm tra đơn vị và ý nghĩa kết quả. c) Sai: Có thể bỏ qua dữ kiện, đơn vị và ý nghĩa vật lí của kết quả. d) Đúng: Khi giải cần đọc đúng dữ kiện, dùng hệ đơn vị thống nhất và kiểm tra điều kiện áp dụng. Kiến thức cốt lõi: Bài vận dụng cao cần tách giai đoạn, dùng đúng bảo toàn, liên hệ khối lượng–năng lượng và định luật phóng xạ, đồng thời kiểm tra đơn vị và ý nghĩa kết quả
\item (Đúng) Bài vận dụng cao cần tách giai đoạn, dùng đúng bảo toàn, liên hệ khối lượng–năng lượng và định luật phóng xạ, đồng thời kiểm tra đơn vị và ý nghĩa kết quả.
\item (Sai) Có thể bỏ qua dữ kiện, đơn vị và ý nghĩa vật lí của kết quả.
\item (Đúng) Khi giải cần đọc đúng dữ kiện, dùng hệ đơn vị thống nhất và kiểm tra điều kiện áp dụng.
\end{itemchoice}
}
\end{ex}

% ===== Câu 6 =====
% ID: L12C4B25-02-TL
% Mức: TH
\begin{bt}
% ID: L12C4B25-02-TL
% Mức: TH
Giải thích vì sao nhận định sau đúng: “Bài vận dụng cao cần tách giai đoạn, dùng đúng bảo toàn, liên hệ khối lượng–năng lượng và định luật phóng xạ, đồng thời kiểm tra đơn vị và ý nghĩa kết quả.”
\loigiai{
Cần nêu được: Bài vận dụng cao cần tách giai đoạn, dùng đúng bảo toàn, liên hệ khối lượng–năng lượng và định luật phóng xạ, đồng thời kiểm tra đơn vị và ý nghĩa kết quả. Khi vận dụng phải xác định đúng dữ kiện, điều kiện áp dụng, hệ đơn vị và kiểm tra tính hợp lí của kết quả. Nhận định “Có thể dùng một công thức duy nhất cho mọi bài toán hạt nhân mà không xét điều kiện.” sai vì trái với kiến thức trên.
}
\end{bt}

% ===== Câu 7 =====
% ID: L12C4B25-02-TLN
% Mức: TH
\begin{ex}
% ID: L12C4B25-02-TLN
% Mức: TH
Trong bốn phát biểu sau về bài toán vận dụng cao kết hợp nhiều nội dung của chương iv, có bao nhiêu phát biểu đúng? Chỉ ghi một số nguyên. (1) Có thể dùng một công thức duy nhất cho mọi bài toán hạt nhân mà không xét điều kiện. (2) Bài vận dụng cao cần tách giai đoạn, dùng đúng bảo toàn, liên hệ khối lượng–năng lượng và định luật phóng xạ, đồng thời kiểm tra đơn vị và ý nghĩa kết quả. (3) Kết quả phải phù hợp ý nghĩa vật lí. (4) Mọi đại lượng đều có cùng bản chất.
\shortans{2}
\loigiai{
Đối chiếu từng phát biểu với kiến thức cốt lõi: Bài vận dụng cao cần tách giai đoạn, dùng đúng bảo toàn, liên hệ khối lượng–năng lượng và định luật phóng xạ, đồng thời kiểm tra đơn vị và ý nghĩa kết quả. Có 2 phát biểu đúng.
}
\end{ex}

% ===== Câu 8 =====
% ID: L12C4B25-02-TN
% Mức: NB
\begin{ex}
% ID: L12C4B25-02-TN
% Mức: NB
Nhận định nào phù hợp với kiến thức Vật lí về bài toán vận dụng cao kết hợp nhiều nội dung của chương iv?
\choice
{Có thể dùng một công thức duy nhất cho mọi bài toán hạt nhân mà không xét điều kiện.}
{\True Bài vận dụng cao cần tách giai đoạn, dùng đúng bảo toàn, liên hệ khối lượng–năng lượng và định luật phóng xạ, đồng thời kiểm tra đơn vị và ý nghĩa kết quả.}
{Có thể bỏ qua điều kiện áp dụng và đơn vị trong mọi trường hợp.}
{Kết quả của bài toán không phụ thuộc dữ kiện đã cho.}
\loigiai{
Bài vận dụng cao cần tách giai đoạn, dùng đúng bảo toàn, liên hệ khối lượng–năng lượng và định luật phóng xạ, đồng thời kiểm tra đơn vị và ý nghĩa kết quả. Vì vậy phương án $B$ đúng; các phương án còn lại trái với kiến thức hoặc điều kiện áp dụng.
}
\end{ex}

% ===== Câu 9 =====
% ID: L12C4B25-03-DS
% Mức: VD
\begin{ex}
% ID: L12C4B25-03-DS
% Mức: VD
Xét bài toán vận dụng cao kết hợp nhiều nội dung của chương iv (tình huống 3). Hãy xác định tính đúng, sai của bốn phát biểu sau.
\choiceTF
{\True Bài vận dụng cao cần tách giai đoạn, dùng đúng bảo toàn, liên hệ khối lượng–năng lượng và định luật phóng xạ, đồng thời kiểm tra đơn vị và ý nghĩa kết quả.}
{Có thể dùng một công thức duy nhất cho mọi bài toán hạt nhân mà không xét điều kiện.}
{\True Khi giải cần đọc đúng dữ kiện, dùng hệ đơn vị thống nhất và kiểm tra điều kiện áp dụng.}
{Có thể bỏ qua dữ kiện, đơn vị và ý nghĩa vật lí của kết quả.}
\loigiai{
\begin{itemchoice}
\item (Đúng) Bài vận dụng cao cần tách giai đoạn, dùng đúng bảo toàn, liên hệ khối lượng–năng lượng và định luật phóng xạ, đồng thời kiểm tra đơn vị và ý nghĩa kết quả. (Vì): Bài vận dụng cao cần tách giai đoạn, dùng đúng bảo toàn, liên hệ khối lượng–năng lượng và định luật phóng xạ, đồng thời kiểm tra đơn vị và ý nghĩa kết quả. b) Sai: Có thể dùng một công thức duy nhất cho mọi bài toán hạt nhân mà không xét điều kiện. c) Đúng: Khi giải cần đọc đúng dữ kiện, dùng hệ đơn vị thống nhất và kiểm tra điều kiện áp dụng. d) Sai: Có thể bỏ qua dữ kiện, đơn vị và ý nghĩa vật lí của kết quả. Kiến thức cốt lõi: Bài vận dụng cao cần tách giai đoạn, dùng đúng bảo toàn, liên hệ khối lượng–năng lượng và định luật phóng xạ, đồng thời kiểm tra đơn vị và ý nghĩa kết quả
\item (Sai) Có thể dùng một công thức duy nhất cho mọi bài toán hạt nhân mà không xét điều kiện.
\item (Đúng) Khi giải cần đọc đúng dữ kiện, dùng hệ đơn vị thống nhất và kiểm tra điều kiện áp dụng.
\item (Sai) Có thể bỏ qua dữ kiện, đơn vị và ý nghĩa vật lí của kết quả.
\end{itemchoice}
}
\end{ex}

% ===== Câu 10 =====
% ID: L12C4B25-03-TL
% Mức: VD
\begin{bt}
% ID: L12C4B25-03-TL
% Mức: VD
Phân tích sai lầm trong nhận định: “Có thể dùng một công thức duy nhất cho mọi bài toán hạt nhân mà không xét điều kiện.”
\loigiai{
Cần nêu được: Bài vận dụng cao cần tách giai đoạn, dùng đúng bảo toàn, liên hệ khối lượng–năng lượng và định luật phóng xạ, đồng thời kiểm tra đơn vị và ý nghĩa kết quả. Khi vận dụng phải xác định đúng dữ kiện, điều kiện áp dụng, hệ đơn vị và kiểm tra tính hợp lí của kết quả. Nhận định “Có thể dùng một công thức duy nhất cho mọi bài toán hạt nhân mà không xét điều kiện.” sai vì trái với kiến thức trên.
}
\end{bt}

% ===== Câu 11 =====
% ID: L12C4B25-03-TLN
% Mức: TH
\begin{ex}
% ID: L12C4B25-03-TLN
% Mức: TH
Trong bốn phát biểu sau về bài toán vận dụng cao kết hợp nhiều nội dung của chương iv, có bao nhiêu phát biểu đúng? Chỉ ghi một số nguyên. (1) Cần đổi các đại lượng về hệ đơn vị thống nhất. (2) Bài vận dụng cao cần tách giai đoạn, dùng đúng bảo toàn, liên hệ khối lượng–năng lượng và định luật phóng xạ, đồng thời kiểm tra đơn vị và ý nghĩa kết quả. (3) Có thể dùng một công thức duy nhất cho mọi bài toán hạt nhân mà không xét điều kiện. (4) Không cần kiểm tra dữ kiện.
\shortans{2}
\loigiai{
Đối chiếu từng phát biểu với kiến thức cốt lõi: Bài vận dụng cao cần tách giai đoạn, dùng đúng bảo toàn, liên hệ khối lượng–năng lượng và định luật phóng xạ, đồng thời kiểm tra đơn vị và ý nghĩa kết quả. Có 2 phát biểu đúng.
}
\end{ex}

% ===== Câu 12 =====
% ID: L12C4B25-03-TN
% Mức: NB
\begin{ex}
% ID: L12C4B25-03-TN
% Mức: NB
Chọn kết luận chính xác về bài toán vận dụng cao kết hợp nhiều nội dung của chương iv?
\choice
{Có thể dùng một công thức duy nhất cho mọi bài toán hạt nhân mà không xét điều kiện.}
{Có thể bỏ qua điều kiện áp dụng và đơn vị trong mọi trường hợp.}
{\True Bài vận dụng cao cần tách giai đoạn, dùng đúng bảo toàn, liên hệ khối lượng–năng lượng và định luật phóng xạ, đồng thời kiểm tra đơn vị và ý nghĩa kết quả.}
{Kết quả của bài toán không phụ thuộc dữ kiện đã cho.}
\loigiai{
Bài vận dụng cao cần tách giai đoạn, dùng đúng bảo toàn, liên hệ khối lượng–năng lượng và định luật phóng xạ, đồng thời kiểm tra đơn vị và ý nghĩa kết quả. Vì vậy phương án $C$ đúng; các phương án còn lại trái với kiến thức hoặc điều kiện áp dụng.
}
\end{ex}

% ===== Câu 13 =====
% ID: L12C4B25-04-DS
% Mức: VD
\begin{ex}
% ID: L12C4B25-04-DS
% Mức: VD
Xét bài toán vận dụng cao kết hợp nhiều nội dung của chương iv (tình huống 4). Hãy xác định tính đúng, sai của bốn phát biểu sau.
\choiceTF
{Có thể dùng một công thức duy nhất cho mọi bài toán hạt nhân mà không xét điều kiện.}
{\True Bài vận dụng cao cần tách giai đoạn, dùng đúng bảo toàn, liên hệ khối lượng–năng lượng và định luật phóng xạ, đồng thời kiểm tra đơn vị và ý nghĩa kết quả.}
{Có thể bỏ qua dữ kiện, đơn vị và ý nghĩa vật lí của kết quả.}
{\True Khi giải cần đọc đúng dữ kiện, dùng hệ đơn vị thống nhất và kiểm tra điều kiện áp dụng.}
\loigiai{
\begin{itemchoice}
\item (Sai) Có thể dùng một công thức duy nhất cho mọi bài toán hạt nhân mà không xét điều kiện. (Vì): Có thể dùng một công thức duy nhất cho mọi bài toán hạt nhân mà không xét điều kiện. b) Đúng: Bài vận dụng cao cần tách giai đoạn, dùng đúng bảo toàn, liên hệ khối lượng–năng lượng và định luật phóng xạ, đồng thời kiểm tra đơn vị và ý nghĩa kết quả. c) Sai: Có thể bỏ qua dữ kiện, đơn vị và ý nghĩa vật lí của kết quả. d) Đúng: Khi giải cần đọc đúng dữ kiện, dùng hệ đơn vị thống nhất và kiểm tra điều kiện áp dụng. Kiến thức cốt lõi: Bài vận dụng cao cần tách giai đoạn, dùng đúng bảo toàn, liên hệ khối lượng–năng lượng và định luật phóng xạ, đồng thời kiểm tra đơn vị và ý nghĩa kết quả
\item (Đúng) Bài vận dụng cao cần tách giai đoạn, dùng đúng bảo toàn, liên hệ khối lượng–năng lượng và định luật phóng xạ, đồng thời kiểm tra đơn vị và ý nghĩa kết quả.
\item (Sai) Có thể bỏ qua dữ kiện, đơn vị và ý nghĩa vật lí của kết quả.
\item (Đúng) Khi giải cần đọc đúng dữ kiện, dùng hệ đơn vị thống nhất và kiểm tra điều kiện áp dụng.
\end{itemchoice}
}
\end{ex}

% ===== Câu 14 =====
% ID: L12C4B25-04-TL
% Mức: VD
\begin{bt}
% ID: L12C4B25-04-TL
% Mức: VD
Đề xuất các bước giải một bài toán thuộc dạng bài toán vận dụng cao kết hợp nhiều nội dung của chương iv.
\loigiai{
Cần nêu được: Bài vận dụng cao cần tách giai đoạn, dùng đúng bảo toàn, liên hệ khối lượng–năng lượng và định luật phóng xạ, đồng thời kiểm tra đơn vị và ý nghĩa kết quả. Khi vận dụng phải xác định đúng dữ kiện, điều kiện áp dụng, hệ đơn vị và kiểm tra tính hợp lí của kết quả. Nhận định “Có thể dùng một công thức duy nhất cho mọi bài toán hạt nhân mà không xét điều kiện.” sai vì trái với kiến thức trên.
}
\end{bt}

% ===== Câu 15 =====
% ID: L12C4B25-04-TLN
% Mức: VD
\begin{ex}
% ID: L12C4B25-04-TLN
% Mức: VD
Trong bốn phát biểu sau về bài toán vận dụng cao kết hợp nhiều nội dung của chương iv, có bao nhiêu phát biểu đúng? Chỉ ghi một số nguyên. (1) Bài vận dụng cao cần tách giai đoạn, dùng đúng bảo toàn, liên hệ khối lượng–năng lượng và định luật phóng xạ, đồng thời kiểm tra đơn vị và ý nghĩa kết quả. (2) Cần đối chiếu kết quả với điều kiện bài toán. (3) Có thể dùng một công thức duy nhất cho mọi bài toán hạt nhân mà không xét điều kiện. (4) Mọi trường hợp đều cho cùng một kết quả.
\shortans{2}
\loigiai{
Đối chiếu từng phát biểu với kiến thức cốt lõi: Bài vận dụng cao cần tách giai đoạn, dùng đúng bảo toàn, liên hệ khối lượng–năng lượng và định luật phóng xạ, đồng thời kiểm tra đơn vị và ý nghĩa kết quả. Có 2 phát biểu đúng.
}
\end{ex}

% ===== Câu 16 =====
% ID: L12C4B25-04-TN
% Mức: TH
\begin{ex}
% ID: L12C4B25-04-TN
% Mức: TH
Nội dung nào mô tả đúng nhất về bài toán vận dụng cao kết hợp nhiều nội dung của chương iv?
\choice
{Có thể dùng một công thức duy nhất cho mọi bài toán hạt nhân mà không xét điều kiện.}
{Có thể bỏ qua điều kiện áp dụng và đơn vị trong mọi trường hợp.}
{Kết quả của bài toán không phụ thuộc dữ kiện đã cho.}
{\True Bài vận dụng cao cần tách giai đoạn, dùng đúng bảo toàn, liên hệ khối lượng–năng lượng và định luật phóng xạ, đồng thời kiểm tra đơn vị và ý nghĩa kết quả.}
\loigiai{
Bài vận dụng cao cần tách giai đoạn, dùng đúng bảo toàn, liên hệ khối lượng–năng lượng và định luật phóng xạ, đồng thời kiểm tra đơn vị và ý nghĩa kết quả. Vì vậy phương án $D$ đúng; các phương án còn lại trái với kiến thức hoặc điều kiện áp dụng.
}
\end{ex}

% ===== Câu 17 =====
% ID: L12C4B25-05-DS
% Mức: VD
\begin{ex}
% ID: L12C4B25-05-DS
% Mức: VD
Xét bài toán vận dụng cao kết hợp nhiều nội dung của chương iv (tình huống 5). Hãy xác định tính đúng, sai của bốn phát biểu sau.
\choiceTF
{\True Bài vận dụng cao cần tách giai đoạn, dùng đúng bảo toàn, liên hệ khối lượng–năng lượng và định luật phóng xạ, đồng thời kiểm tra đơn vị và ý nghĩa kết quả.}
{Có thể bỏ qua dữ kiện, đơn vị và ý nghĩa vật lí của kết quả.}
{Có thể dùng một công thức duy nhất cho mọi bài toán hạt nhân mà không xét điều kiện.}
{Có thể bỏ qua dữ kiện, đơn vị và ý nghĩa vật lí của kết quả.}
\loigiai{
\begin{itemchoice}
\item (Đúng) Bài vận dụng cao cần tách giai đoạn, dùng đúng bảo toàn, liên hệ khối lượng–năng lượng và định luật phóng xạ, đồng thời kiểm tra đơn vị và ý nghĩa kết quả. (Vì): Bài vận dụng cao cần tách giai đoạn, dùng đúng bảo toàn, liên hệ khối lượng–năng lượng và định luật phóng xạ, đồng thời kiểm tra đơn vị và ý nghĩa kết quả. b) Sai: Có thể bỏ qua dữ kiện, đơn vị và ý nghĩa vật lí của kết quả. c) Sai: Có thể dùng một công thức duy nhất cho mọi bài toán hạt nhân mà không xét điều kiện. d) Sai: Có thể bỏ qua dữ kiện, đơn vị và ý nghĩa vật lí của kết quả. Kiến thức cốt lõi: Bài vận dụng cao cần tách giai đoạn, dùng đúng bảo toàn, liên hệ khối lượng–năng lượng và định luật phóng xạ, đồng thời kiểm tra đơn vị và ý nghĩa kết quả
\item (Sai) Có thể bỏ qua dữ kiện, đơn vị và ý nghĩa vật lí của kết quả.
\item (Sai) Có thể dùng một công thức duy nhất cho mọi bài toán hạt nhân mà không xét điều kiện.
\item (Sai) Có thể bỏ qua dữ kiện, đơn vị và ý nghĩa vật lí của kết quả.
\end{itemchoice}
}
\end{ex}

% ===== Câu 18 =====
% ID: L12C4B25-05-TL
% Mức: VD
\begin{bt}
% ID: L12C4B25-05-TL
% Mức: VD
Nêu một tình huống thực tiễn liên quan đến bài toán vận dụng cao kết hợp nhiều nội dung của chương iv và giải thích bằng kiến thức Vật lí.
\loigiai{
Cần nêu được: Bài vận dụng cao cần tách giai đoạn, dùng đúng bảo toàn, liên hệ khối lượng–năng lượng và định luật phóng xạ, đồng thời kiểm tra đơn vị và ý nghĩa kết quả. Khi vận dụng phải xác định đúng dữ kiện, điều kiện áp dụng, hệ đơn vị và kiểm tra tính hợp lí của kết quả. Nhận định “Có thể dùng một công thức duy nhất cho mọi bài toán hạt nhân mà không xét điều kiện.” sai vì trái với kiến thức trên.
}
\end{bt}

% ===== Câu 19 =====
% ID: L12C4B25-05-TLN
% Mức: VD
\begin{ex}
% ID: L12C4B25-05-TLN
% Mức: VD
Trong bốn phát biểu sau về bài toán vận dụng cao kết hợp nhiều nội dung của chương iv, có bao nhiêu phát biểu đúng? Chỉ ghi một số nguyên. (1) Có thể dùng một công thức duy nhất cho mọi bài toán hạt nhân mà không xét điều kiện. (2) Cần đọc đúng dữ kiện và giả thiết. (3) Bài vận dụng cao cần tách giai đoạn, dùng đúng bảo toàn, liên hệ khối lượng–năng lượng và định luật phóng xạ, đồng thời kiểm tra đơn vị và ý nghĩa kết quả. (4) Có thể thay số trước khi chọn công thức.
\shortans{2}
\loigiai{
Đối chiếu từng phát biểu với kiến thức cốt lõi: Bài vận dụng cao cần tách giai đoạn, dùng đúng bảo toàn, liên hệ khối lượng–năng lượng và định luật phóng xạ, đồng thời kiểm tra đơn vị và ý nghĩa kết quả. Có 2 phát biểu đúng.
}
\end{ex}

% ===== Câu 20 =====
% ID: L12C4B25-05-TN
% Mức: TH
\begin{ex}
% ID: L12C4B25-05-TN
% Mức: TH
Khi phân tích, phát biểu nào đúng về bài toán vận dụng cao kết hợp nhiều nội dung của chương iv?
\choice
{\True Bài vận dụng cao cần tách giai đoạn, dùng đúng bảo toàn, liên hệ khối lượng–năng lượng và định luật phóng xạ, đồng thời kiểm tra đơn vị và ý nghĩa kết quả.}
{Có thể dùng một công thức duy nhất cho mọi bài toán hạt nhân mà không xét điều kiện.}
{Có thể bỏ qua điều kiện áp dụng và đơn vị trong mọi trường hợp.}
{Kết quả của bài toán không phụ thuộc dữ kiện đã cho.}
\loigiai{
Bài vận dụng cao cần tách giai đoạn, dùng đúng bảo toàn, liên hệ khối lượng–năng lượng và định luật phóng xạ, đồng thời kiểm tra đơn vị và ý nghĩa kết quả. Vì vậy phương án $A$ đúng; các phương án còn lại trái với kiến thức hoặc điều kiện áp dụng.
}
\end{ex}

% ===== Câu 21 =====
% ID: L12C4B25-06-DS
% Mức: TH
\dangbt{Tổng hợp cấu trúc hạt nhân và đồng vị}
\begin{ex}
% ID: L12C4B25-06-DS
% Mức: TH
Xét tổng hợp cấu trúc hạt nhân và đồng vị (tình huống 1). Hãy xác định tính đúng, sai của bốn phát biểu sau.
\choiceTF
{\True Bài tổng hợp cấu trúc cần xác định đúng $Z,N,A$, quan hệ đồng vị và trung bình có trọng số khi biết tỉ lệ đồng vị.}
{Các đồng vị của cùng nguyên tố có số proton khác nhau.}
{\True Khi giải cần đọc đúng dữ kiện, dùng hệ đơn vị thống nhất và kiểm tra điều kiện áp dụng.}
{Có thể bỏ qua dữ kiện, đơn vị và ý nghĩa vật lí của kết quả.}
\loigiai{
\begin{itemchoice}
\item (Đúng) Bài tổng hợp cấu trúc cần xác định đúng $Z,N,A$, quan hệ đồng vị và trung bình có trọng số khi biết tỉ lệ đồng vị. (Vì): Bài tổng hợp cấu trúc cần xác định đúng $Z,N,A$, quan hệ đồng vị và trung bình có trọng số khi biết tỉ lệ đồng vị. b) Sai: Các đồng vị của cùng nguyên tố có số proton khác nhau. c) Đúng: Khi giải cần đọc đúng dữ kiện, dùng hệ đơn vị thống nhất và kiểm tra điều kiện áp dụng. d) Sai: Có thể bỏ qua dữ kiện, đơn vị và ý nghĩa vật lí của kết quả. Kiến thức cốt lõi: Bài tổng hợp cấu trúc cần xác định đúng $Z,N,A$, quan hệ đồng vị và trung bình có trọng số khi biết tỉ lệ đồng vị
\item (Sai) Các đồng vị của cùng nguyên tố có số proton khác nhau.
\item (Đúng) Khi giải cần đọc đúng dữ kiện, dùng hệ đơn vị thống nhất và kiểm tra điều kiện áp dụng.
\item (Sai) Có thể bỏ qua dữ kiện, đơn vị và ý nghĩa vật lí của kết quả.
\end{itemchoice}
}
\end{ex}

% ===== Câu 22 =====
% ID: L12C4B25-06-TL
% Mức: VDC
\dangbt{Bài toán vận dụng cao kết hợp nhiều nội dung của Chương IV}
\begin{bt}
% ID: L12C4B25-06-TL
% Mức: VDC
So sánh nhận định đúng “Bài vận dụng cao cần tách giai đoạn, dùng đúng bảo toàn, liên hệ khối lượng–năng lượng và định luật phóng xạ, đồng thời kiểm tra đơn vị và ý nghĩa kết quả.” với nhận định sai “Có thể dùng một công thức duy nhất cho mọi bài toán hạt nhân mà không xét điều kiện.”; chỉ rõ căn cứ Vật lí.
\loigiai{
Cần nêu được: Bài vận dụng cao cần tách giai đoạn, dùng đúng bảo toàn, liên hệ khối lượng–năng lượng và định luật phóng xạ, đồng thời kiểm tra đơn vị và ý nghĩa kết quả. Khi vận dụng phải xác định đúng dữ kiện, điều kiện áp dụng, hệ đơn vị và kiểm tra tính hợp lí của kết quả. Nhận định “Có thể dùng một công thức duy nhất cho mọi bài toán hạt nhân mà không xét điều kiện.” sai vì trái với kiến thức trên.
}
\end{bt}

% ===== Câu 23 =====
% ID: L12C4B25-06-TLN
% Mức: VDC
\begin{ex}
% ID: L12C4B25-06-TLN
% Mức: VDC
Trong bốn phát biểu sau về bài toán vận dụng cao kết hợp nhiều nội dung của chương iv, có bao nhiêu phát biểu đúng? Chỉ ghi một số nguyên. (1) Kết quả cần có ý nghĩa vật lí. (2) Có thể dùng một công thức duy nhất cho mọi bài toán hạt nhân mà không xét điều kiện. (3) Phải dùng đúng định luật cho tình huống. (4) Bài vận dụng cao cần tách giai đoạn, dùng đúng bảo toàn, liên hệ khối lượng–năng lượng và định luật phóng xạ, đồng thời kiểm tra đơn vị và ý nghĩa kết quả.
\shortans{3}
\loigiai{
Đối chiếu từng phát biểu với kiến thức cốt lõi: Bài vận dụng cao cần tách giai đoạn, dùng đúng bảo toàn, liên hệ khối lượng–năng lượng và định luật phóng xạ, đồng thời kiểm tra đơn vị và ý nghĩa kết quả. Có 3 phát biểu đúng.
}
\end{ex}

% ===== Câu 24 =====
% ID: L12C4B25-06-TN
% Mức: TH
\begin{ex}
% ID: L12C4B25-06-TN
% Mức: TH
Kết luận nào của học sinh là đúng về bài toán vận dụng cao kết hợp nhiều nội dung của chương iv?
\choice
{Có thể dùng một công thức duy nhất cho mọi bài toán hạt nhân mà không xét điều kiện.}
{\True Bài vận dụng cao cần tách giai đoạn, dùng đúng bảo toàn, liên hệ khối lượng–năng lượng và định luật phóng xạ, đồng thời kiểm tra đơn vị và ý nghĩa kết quả.}
{Có thể bỏ qua điều kiện áp dụng và đơn vị trong mọi trường hợp.}
{Kết quả của bài toán không phụ thuộc dữ kiện đã cho.}
\loigiai{
Bài vận dụng cao cần tách giai đoạn, dùng đúng bảo toàn, liên hệ khối lượng–năng lượng và định luật phóng xạ, đồng thời kiểm tra đơn vị và ý nghĩa kết quả. Vì vậy phương án $B$ đúng; các phương án còn lại trái với kiến thức hoặc điều kiện áp dụng.
}
\end{ex}

% ===== Câu 25 =====
% ID: L12C4B25-07-DS
% Mức: TH
\dangbt{Tổng hợp cấu trúc hạt nhân và đồng vị}
\begin{ex}
% ID: L12C4B25-07-DS
% Mức: TH
Xét tổng hợp cấu trúc hạt nhân và đồng vị (tình huống 2). Hãy xác định tính đúng, sai của bốn phát biểu sau.
\choiceTF
{Các đồng vị của cùng nguyên tố có số proton khác nhau.}
{\True Bài tổng hợp cấu trúc cần xác định đúng $Z,N,A$, quan hệ đồng vị và trung bình có trọng số khi biết tỉ lệ đồng vị.}
{Có thể bỏ qua dữ kiện, đơn vị và ý nghĩa vật lí của kết quả.}
{\True Khi giải cần đọc đúng dữ kiện, dùng hệ đơn vị thống nhất và kiểm tra điều kiện áp dụng.}
\loigiai{
\begin{itemchoice}
\item (Sai) Các đồng vị của cùng nguyên tố có số proton khác nhau. (Vì): Các đồng vị của cùng nguyên tố có số proton khác nhau. b) Đúng: Bài tổng hợp cấu trúc cần xác định đúng $Z,N,A$, quan hệ đồng vị và trung bình có trọng số khi biết tỉ lệ đồng vị. c) Sai: Có thể bỏ qua dữ kiện, đơn vị và ý nghĩa vật lí của kết quả. d) Đúng: Khi giải cần đọc đúng dữ kiện, dùng hệ đơn vị thống nhất và kiểm tra điều kiện áp dụng. Kiến thức cốt lõi: Bài tổng hợp cấu trúc cần xác định đúng $Z,N,A$, quan hệ đồng vị và trung bình có trọng số khi biết tỉ lệ đồng vị
\item (Đúng) Bài tổng hợp cấu trúc cần xác định đúng $Z,N,A$, quan hệ đồng vị và trung bình có trọng số khi biết tỉ lệ đồng vị.
\item (Sai) Có thể bỏ qua dữ kiện, đơn vị và ý nghĩa vật lí của kết quả.
\item (Đúng) Khi giải cần đọc đúng dữ kiện, dùng hệ đơn vị thống nhất và kiểm tra điều kiện áp dụng.
\end{itemchoice}
}
\end{ex}

% ===== Câu 26 =====
% ID: L12C4B25-07-TL
% Mức: TH
\begin{bt}
% ID: L12C4B25-07-TL
% Mức: TH
Trình bày kiến thức cốt lõi của tổng hợp cấu trúc hạt nhân và đồng vị và nêu điều kiện áp dụng.
\loigiai{
Cần nêu được: Bài tổng hợp cấu trúc cần xác định đúng $Z,N,A$, quan hệ đồng vị và trung bình có trọng số khi biết tỉ lệ đồng vị. Khi vận dụng phải xác định đúng dữ kiện, điều kiện áp dụng, hệ đơn vị và kiểm tra tính hợp lí của kết quả. Nhận định “Các đồng vị của cùng nguyên tố có số proton khác nhau.” sai vì trái với kiến thức trên.
}
\end{bt}

% ===== Câu 27 =====
% ID: L12C4B25-07-TLN
% Mức: TH
\begin{ex}
% ID: L12C4B25-07-TLN
% Mức: TH
Trong bốn phát biểu sau về tổng hợp cấu trúc hạt nhân và đồng vị, có bao nhiêu phát biểu đúng? Chỉ ghi một số nguyên. (1) Bài tổng hợp cấu trúc cần xác định đúng $Z,N,A$, quan hệ đồng vị và trung bình có trọng số khi biết tỉ lệ đồng vị. (2) Các đồng vị của cùng nguyên tố có số proton khác nhau. (3) Cần kiểm tra điều kiện áp dụng trước khi tính toán. (4) Có thể bỏ qua đơn vị của kết quả.
\shortans{2}
\loigiai{
Đối chiếu từng phát biểu với kiến thức cốt lõi: Bài tổng hợp cấu trúc cần xác định đúng $Z,N,A$, quan hệ đồng vị và trung bình có trọng số khi biết tỉ lệ đồng vị. Có 2 phát biểu đúng.
}
\end{ex}

% ===== Câu 28 =====
% ID: L12C4B25-07-TN
% Mức: TH
\dangbt{Bài toán vận dụng cao kết hợp nhiều nội dung của Chương IV}
\begin{ex}
% ID: L12C4B25-07-TN
% Mức: TH
Chọn phát biểu không trái với kiến thức về bài toán vận dụng cao kết hợp nhiều nội dung của chương iv?
\choice
{Có thể dùng một công thức duy nhất cho mọi bài toán hạt nhân mà không xét điều kiện.}
{Có thể bỏ qua điều kiện áp dụng và đơn vị trong mọi trường hợp.}
{\True Bài vận dụng cao cần tách giai đoạn, dùng đúng bảo toàn, liên hệ khối lượng–năng lượng và định luật phóng xạ, đồng thời kiểm tra đơn vị và ý nghĩa kết quả.}
{Kết quả của bài toán không phụ thuộc dữ kiện đã cho.}
\loigiai{
Bài vận dụng cao cần tách giai đoạn, dùng đúng bảo toàn, liên hệ khối lượng–năng lượng và định luật phóng xạ, đồng thời kiểm tra đơn vị và ý nghĩa kết quả. Vì vậy phương án $C$ đúng; các phương án còn lại trái với kiến thức hoặc điều kiện áp dụng.
}
\end{ex}

% ===== Câu 29 =====
% ID: L12C4B25-08-DS
% Mức: VD
\dangbt{Tổng hợp cấu trúc hạt nhân và đồng vị}
\begin{ex}
% ID: L12C4B25-08-DS
% Mức: VD
Xét tổng hợp cấu trúc hạt nhân và đồng vị (tình huống 3). Hãy xác định tính đúng, sai của bốn phát biểu sau.
\choiceTF
{\True Bài tổng hợp cấu trúc cần xác định đúng $Z,N,A$, quan hệ đồng vị và trung bình có trọng số khi biết tỉ lệ đồng vị.}
{Các đồng vị của cùng nguyên tố có số proton khác nhau.}
{\True Khi giải cần đọc đúng dữ kiện, dùng hệ đơn vị thống nhất và kiểm tra điều kiện áp dụng.}
{Có thể bỏ qua dữ kiện, đơn vị và ý nghĩa vật lí của kết quả.}
\loigiai{
\begin{itemchoice}
\item (Đúng) Bài tổng hợp cấu trúc cần xác định đúng $Z,N,A$, quan hệ đồng vị và trung bình có trọng số khi biết tỉ lệ đồng vị. (Vì): Bài tổng hợp cấu trúc cần xác định đúng $Z,N,A$, quan hệ đồng vị và trung bình có trọng số khi biết tỉ lệ đồng vị. b) Sai: Các đồng vị của cùng nguyên tố có số proton khác nhau. c) Đúng: Khi giải cần đọc đúng dữ kiện, dùng hệ đơn vị thống nhất và kiểm tra điều kiện áp dụng. d) Sai: Có thể bỏ qua dữ kiện, đơn vị và ý nghĩa vật lí của kết quả. Kiến thức cốt lõi: Bài tổng hợp cấu trúc cần xác định đúng $Z,N,A$, quan hệ đồng vị và trung bình có trọng số khi biết tỉ lệ đồng vị
\item (Sai) Các đồng vị của cùng nguyên tố có số proton khác nhau.
\item (Đúng) Khi giải cần đọc đúng dữ kiện, dùng hệ đơn vị thống nhất và kiểm tra điều kiện áp dụng.
\item (Sai) Có thể bỏ qua dữ kiện, đơn vị và ý nghĩa vật lí của kết quả.
\end{itemchoice}
}
\end{ex}

% ===== Câu 30 =====
% ID: L12C4B25-08-TL
% Mức: TH
\begin{bt}
% ID: L12C4B25-08-TL
% Mức: TH
Giải thích vì sao nhận định sau đúng: “Bài tổng hợp cấu trúc cần xác định đúng $Z,N,A$, quan hệ đồng vị và trung bình có trọng số khi biết tỉ lệ đồng vị.”
\loigiai{
Cần nêu được: Bài tổng hợp cấu trúc cần xác định đúng $Z,N,A$, quan hệ đồng vị và trung bình có trọng số khi biết tỉ lệ đồng vị. Khi vận dụng phải xác định đúng dữ kiện, điều kiện áp dụng, hệ đơn vị và kiểm tra tính hợp lí của kết quả. Nhận định “Các đồng vị của cùng nguyên tố có số proton khác nhau.” sai vì trái với kiến thức trên.
}
\end{bt}

% ===== Câu 31 =====
% ID: L12C4B25-08-TLN
% Mức: TH
\begin{ex}
% ID: L12C4B25-08-TLN
% Mức: TH
Trong bốn phát biểu sau về tổng hợp cấu trúc hạt nhân và đồng vị, có bao nhiêu phát biểu đúng? Chỉ ghi một số nguyên. (1) Các đồng vị của cùng nguyên tố có số proton khác nhau. (2) Bài tổng hợp cấu trúc cần xác định đúng $Z,N,A$, quan hệ đồng vị và trung bình có trọng số khi biết tỉ lệ đồng vị. (3) Kết quả phải phù hợp ý nghĩa vật lí. (4) Mọi đại lượng đều có cùng bản chất.
\shortans{2}
\loigiai{
Đối chiếu từng phát biểu với kiến thức cốt lõi: Bài tổng hợp cấu trúc cần xác định đúng $Z,N,A$, quan hệ đồng vị và trung bình có trọng số khi biết tỉ lệ đồng vị. Có 2 phát biểu đúng.
}
\end{ex}

% ===== Câu 32 =====
% ID: L12C4B25-08-TN
% Mức: VD
\dangbt{Bài toán vận dụng cao kết hợp nhiều nội dung của Chương IV}
\begin{ex}
% ID: L12C4B25-08-TN
% Mức: VD
Cơ sở Vật lí đúng để giải bài là gì về bài toán vận dụng cao kết hợp nhiều nội dung của chương iv?
\choice
{Có thể dùng một công thức duy nhất cho mọi bài toán hạt nhân mà không xét điều kiện.}
{Có thể bỏ qua điều kiện áp dụng và đơn vị trong mọi trường hợp.}
{Kết quả của bài toán không phụ thuộc dữ kiện đã cho.}
{\True Bài vận dụng cao cần tách giai đoạn, dùng đúng bảo toàn, liên hệ khối lượng–năng lượng và định luật phóng xạ, đồng thời kiểm tra đơn vị và ý nghĩa kết quả.}
\loigiai{
Bài vận dụng cao cần tách giai đoạn, dùng đúng bảo toàn, liên hệ khối lượng–năng lượng và định luật phóng xạ, đồng thời kiểm tra đơn vị và ý nghĩa kết quả. Vì vậy phương án $D$ đúng; các phương án còn lại trái với kiến thức hoặc điều kiện áp dụng.
}
\end{ex}

% ===== Câu 33 =====
% ID: L12C4B25-09-DS
% Mức: VD
\dangbt{Tổng hợp cấu trúc hạt nhân và đồng vị}
\begin{ex}
% ID: L12C4B25-09-DS
% Mức: VD
Xét tổng hợp cấu trúc hạt nhân và đồng vị (tình huống 4). Hãy xác định tính đúng, sai của bốn phát biểu sau.
\choiceTF
{Các đồng vị của cùng nguyên tố có số proton khác nhau.}
{\True Bài tổng hợp cấu trúc cần xác định đúng $Z,N,A$, quan hệ đồng vị và trung bình có trọng số khi biết tỉ lệ đồng vị.}
{Có thể bỏ qua dữ kiện, đơn vị và ý nghĩa vật lí của kết quả.}
{\True Khi giải cần đọc đúng dữ kiện, dùng hệ đơn vị thống nhất và kiểm tra điều kiện áp dụng.}
\loigiai{
\begin{itemchoice}
\item (Sai) Các đồng vị của cùng nguyên tố có số proton khác nhau. (Vì): Các đồng vị của cùng nguyên tố có số proton khác nhau. b) Đúng: Bài tổng hợp cấu trúc cần xác định đúng $Z,N,A$, quan hệ đồng vị và trung bình có trọng số khi biết tỉ lệ đồng vị. c) Sai: Có thể bỏ qua dữ kiện, đơn vị và ý nghĩa vật lí của kết quả. d) Đúng: Khi giải cần đọc đúng dữ kiện, dùng hệ đơn vị thống nhất và kiểm tra điều kiện áp dụng. Kiến thức cốt lõi: Bài tổng hợp cấu trúc cần xác định đúng $Z,N,A$, quan hệ đồng vị và trung bình có trọng số khi biết tỉ lệ đồng vị
\item (Đúng) Bài tổng hợp cấu trúc cần xác định đúng $Z,N,A$, quan hệ đồng vị và trung bình có trọng số khi biết tỉ lệ đồng vị.
\item (Sai) Có thể bỏ qua dữ kiện, đơn vị và ý nghĩa vật lí của kết quả.
\item (Đúng) Khi giải cần đọc đúng dữ kiện, dùng hệ đơn vị thống nhất và kiểm tra điều kiện áp dụng.
\end{itemchoice}
}
\end{ex}

% ===== Câu 34 =====
% ID: L12C4B25-09-TL
% Mức: VD
\begin{bt}
% ID: L12C4B25-09-TL
% Mức: VD
Phân tích sai lầm trong nhận định: “Các đồng vị của cùng nguyên tố có số proton khác nhau.”
\loigiai{
Cần nêu được: Bài tổng hợp cấu trúc cần xác định đúng $Z,N,A$, quan hệ đồng vị và trung bình có trọng số khi biết tỉ lệ đồng vị. Khi vận dụng phải xác định đúng dữ kiện, điều kiện áp dụng, hệ đơn vị và kiểm tra tính hợp lí của kết quả. Nhận định “Các đồng vị của cùng nguyên tố có số proton khác nhau.” sai vì trái với kiến thức trên.
}
\end{bt}

% ===== Câu 35 =====
% ID: L12C4B25-09-TLN
% Mức: TH
\begin{ex}
% ID: L12C4B25-09-TLN
% Mức: TH
Trong bốn phát biểu sau về tổng hợp cấu trúc hạt nhân và đồng vị, có bao nhiêu phát biểu đúng? Chỉ ghi một số nguyên. (1) Cần đổi các đại lượng về hệ đơn vị thống nhất. (2) Bài tổng hợp cấu trúc cần xác định đúng $Z,N,A$, quan hệ đồng vị và trung bình có trọng số khi biết tỉ lệ đồng vị. (3) Các đồng vị của cùng nguyên tố có số proton khác nhau. (4) Không cần kiểm tra dữ kiện.
\shortans{2}
\loigiai{
Đối chiếu từng phát biểu với kiến thức cốt lõi: Bài tổng hợp cấu trúc cần xác định đúng $Z,N,A$, quan hệ đồng vị và trung bình có trọng số khi biết tỉ lệ đồng vị. Có 2 phát biểu đúng.
}
\end{ex}

% ===== Câu 36 =====
% ID: L12C4B25-09-TN
% Mức: VD
\dangbt{Bài toán vận dụng cao kết hợp nhiều nội dung của Chương IV}
\begin{ex}
% ID: L12C4B25-09-TN
% Mức: VD
Trong một bài thực hành, nhận xét nào đúng về bài toán vận dụng cao kết hợp nhiều nội dung của chương iv?
\choice
{\True Bài vận dụng cao cần tách giai đoạn, dùng đúng bảo toàn, liên hệ khối lượng–năng lượng và định luật phóng xạ, đồng thời kiểm tra đơn vị và ý nghĩa kết quả.}
{Có thể dùng một công thức duy nhất cho mọi bài toán hạt nhân mà không xét điều kiện.}
{Có thể bỏ qua điều kiện áp dụng và đơn vị trong mọi trường hợp.}
{Kết quả của bài toán không phụ thuộc dữ kiện đã cho.}
\loigiai{
Bài vận dụng cao cần tách giai đoạn, dùng đúng bảo toàn, liên hệ khối lượng–năng lượng và định luật phóng xạ, đồng thời kiểm tra đơn vị và ý nghĩa kết quả. Vì vậy phương án $A$ đúng; các phương án còn lại trái với kiến thức hoặc điều kiện áp dụng.
}
\end{ex}

% ===== Câu 37 =====
% ID: L12C4B25-10-DS
% Mức: VD
\dangbt{Tổng hợp cấu trúc hạt nhân và đồng vị}
\begin{ex}
% ID: L12C4B25-10-DS
% Mức: VD
Xét tổng hợp cấu trúc hạt nhân và đồng vị (tình huống 5). Hãy xác định tính đúng, sai của bốn phát biểu sau.
\choiceTF
{\True Bài tổng hợp cấu trúc cần xác định đúng $Z,N,A$, quan hệ đồng vị và trung bình có trọng số khi biết tỉ lệ đồng vị.}
{Có thể bỏ qua dữ kiện, đơn vị và ý nghĩa vật lí của kết quả.}
{Các đồng vị của cùng nguyên tố có số proton khác nhau.}
{Có thể bỏ qua dữ kiện, đơn vị và ý nghĩa vật lí của kết quả.}
\loigiai{
\begin{itemchoice}
\item (Đúng) Bài tổng hợp cấu trúc cần xác định đúng $Z,N,A$, quan hệ đồng vị và trung bình có trọng số khi biết tỉ lệ đồng vị. (Vì): Bài tổng hợp cấu trúc cần xác định đúng $Z,N,A$, quan hệ đồng vị và trung bình có trọng số khi biết tỉ lệ đồng vị. b) Sai: Có thể bỏ qua dữ kiện, đơn vị và ý nghĩa vật lí của kết quả. c) Sai: Các đồng vị của cùng nguyên tố có số proton khác nhau. d) Sai: Có thể bỏ qua dữ kiện, đơn vị và ý nghĩa vật lí của kết quả. Kiến thức cốt lõi: Bài tổng hợp cấu trúc cần xác định đúng $Z,N,A$, quan hệ đồng vị và trung bình có trọng số khi biết tỉ lệ đồng vị
\item (Sai) Có thể bỏ qua dữ kiện, đơn vị và ý nghĩa vật lí của kết quả.
\item (Sai) Các đồng vị của cùng nguyên tố có số proton khác nhau.
\item (Sai) Có thể bỏ qua dữ kiện, đơn vị và ý nghĩa vật lí của kết quả.
\end{itemchoice}
}
\end{ex}

% ===== Câu 38 =====
% ID: L12C4B25-10-TL
% Mức: VD
\begin{bt}
% ID: L12C4B25-10-TL
% Mức: VD
Đề xuất các bước giải một bài toán thuộc dạng tổng hợp cấu trúc hạt nhân và đồng vị.
\loigiai{
Cần nêu được: Bài tổng hợp cấu trúc cần xác định đúng $Z,N,A$, quan hệ đồng vị và trung bình có trọng số khi biết tỉ lệ đồng vị. Khi vận dụng phải xác định đúng dữ kiện, điều kiện áp dụng, hệ đơn vị và kiểm tra tính hợp lí của kết quả. Nhận định “Các đồng vị của cùng nguyên tố có số proton khác nhau.” sai vì trái với kiến thức trên.
}
\end{bt}

% ===== Câu 39 =====
% ID: L12C4B25-10-TLN
% Mức: VD
\begin{ex}
% ID: L12C4B25-10-TLN
% Mức: VD
Trong bốn phát biểu sau về tổng hợp cấu trúc hạt nhân và đồng vị, có bao nhiêu phát biểu đúng? Chỉ ghi một số nguyên. (1) Bài tổng hợp cấu trúc cần xác định đúng $Z,N,A$, quan hệ đồng vị và trung bình có trọng số khi biết tỉ lệ đồng vị. (2) Cần đối chiếu kết quả với điều kiện bài toán. (3) Các đồng vị của cùng nguyên tố có số proton khác nhau. (4) Mọi trường hợp đều cho cùng một kết quả.
\shortans{2}
\loigiai{
Đối chiếu từng phát biểu với kiến thức cốt lõi: Bài tổng hợp cấu trúc cần xác định đúng $Z,N,A$, quan hệ đồng vị và trung bình có trọng số khi biết tỉ lệ đồng vị. Có 2 phát biểu đúng.
}
\end{ex}

% ===== Câu 40 =====
% ID: L12C4B25-10-TN
% Mức: VD
\dangbt{Bài toán vận dụng cao kết hợp nhiều nội dung của Chương IV}
\begin{ex}
% ID: L12C4B25-10-TN
% Mức: VD
Kết luận đúng khi vận dụng là về bài toán vận dụng cao kết hợp nhiều nội dung của chương iv?
\choice
{Có thể dùng một công thức duy nhất cho mọi bài toán hạt nhân mà không xét điều kiện.}
{\True Bài vận dụng cao cần tách giai đoạn, dùng đúng bảo toàn, liên hệ khối lượng–năng lượng và định luật phóng xạ, đồng thời kiểm tra đơn vị và ý nghĩa kết quả.}
{Có thể bỏ qua điều kiện áp dụng và đơn vị trong mọi trường hợp.}
{Kết quả của bài toán không phụ thuộc dữ kiện đã cho.}
\loigiai{
Bài vận dụng cao cần tách giai đoạn, dùng đúng bảo toàn, liên hệ khối lượng–năng lượng và định luật phóng xạ, đồng thời kiểm tra đơn vị và ý nghĩa kết quả. Vì vậy phương án $B$ đúng; các phương án còn lại trái với kiến thức hoặc điều kiện áp dụng.
}
\end{ex}

% ===== Câu 41 =====
% ID: L12C4B25-11-DS
% Mức: TH
\dangbt{Tổng hợp độ hụt khối, năng lượng liên kết và độ bền vững}
\begin{ex}
% ID: L12C4B25-11-DS
% Mức: TH
Xét tổng hợp độ hụt khối, năng lượng liên kết và độ bền vững (tình huống 1). Hãy xác định tính đúng, sai của bốn phát biểu sau.
\choiceTF
{\True Phải tính độ hụt khối đúng dấu, đổi sang năng lượng liên kết và dùng năng lượng liên kết riêng khi so sánh độ bền.}
{Chỉ cần so sánh năng lượng liên kết toàn phần để kết luận độ bền cho mọi hạt nhân.}
{\True Khi giải cần đọc đúng dữ kiện, dùng hệ đơn vị thống nhất và kiểm tra điều kiện áp dụng.}
{Có thể bỏ qua dữ kiện, đơn vị và ý nghĩa vật lí của kết quả.}
\loigiai{
\begin{itemchoice}
\item (Đúng) Phải tính độ hụt khối đúng dấu, đổi sang năng lượng liên kết và dùng năng lượng liên kết riêng khi so sánh độ bền. (Vì): Phải tính độ hụt khối đúng dấu, đổi sang năng lượng liên kết và dùng năng lượng liên kết riêng khi so sánh độ bền. b) Sai: Chỉ cần so sánh năng lượng liên kết toàn phần để kết luận độ bền cho mọi hạt nhân. c) Đúng: Khi giải cần đọc đúng dữ kiện, dùng hệ đơn vị thống nhất và kiểm tra điều kiện áp dụng. d) Sai: Có thể bỏ qua dữ kiện, đơn vị và ý nghĩa vật lí của kết quả. Kiến thức cốt lõi: Phải tính độ hụt khối đúng dấu, đổi sang năng lượng liên kết và dùng năng lượng liên kết riêng khi so sánh độ bền
\item (Sai) Chỉ cần so sánh năng lượng liên kết toàn phần để kết luận độ bền cho mọi hạt nhân.
\item (Đúng) Khi giải cần đọc đúng dữ kiện, dùng hệ đơn vị thống nhất và kiểm tra điều kiện áp dụng.
\item (Sai) Có thể bỏ qua dữ kiện, đơn vị và ý nghĩa vật lí của kết quả.
\end{itemchoice}
}
\end{ex}

% ===== Câu 42 =====
% ID: L12C4B25-11-TL
% Mức: VD
\dangbt{Tổng hợp cấu trúc hạt nhân và đồng vị}
\begin{bt}
% ID: L12C4B25-11-TL
% Mức: VD
Nêu một tình huống thực tiễn liên quan đến tổng hợp cấu trúc hạt nhân và đồng vị và giải thích bằng kiến thức Vật lí.
\loigiai{
Cần nêu được: Bài tổng hợp cấu trúc cần xác định đúng $Z,N,A$, quan hệ đồng vị và trung bình có trọng số khi biết tỉ lệ đồng vị. Khi vận dụng phải xác định đúng dữ kiện, điều kiện áp dụng, hệ đơn vị và kiểm tra tính hợp lí của kết quả. Nhận định “Các đồng vị của cùng nguyên tố có số proton khác nhau.” sai vì trái với kiến thức trên.
}
\end{bt}

% ===== Câu 43 =====
% ID: L12C4B25-11-TLN
% Mức: VD
\begin{ex}
% ID: L12C4B25-11-TLN
% Mức: VD
Trong bốn phát biểu sau về tổng hợp cấu trúc hạt nhân và đồng vị, có bao nhiêu phát biểu đúng? Chỉ ghi một số nguyên. (1) Các đồng vị của cùng nguyên tố có số proton khác nhau. (2) Cần đọc đúng dữ kiện và giả thiết. (3) Bài tổng hợp cấu trúc cần xác định đúng $Z,N,A$, quan hệ đồng vị và trung bình có trọng số khi biết tỉ lệ đồng vị. (4) Có thể thay số trước khi chọn công thức.
\shortans{2}
\loigiai{
Đối chiếu từng phát biểu với kiến thức cốt lõi: Bài tổng hợp cấu trúc cần xác định đúng $Z,N,A$, quan hệ đồng vị và trung bình có trọng số khi biết tỉ lệ đồng vị. Có 2 phát biểu đúng.
}
\end{ex}

% ===== Câu 44 =====
% ID: L12C4B25-11-TN
% Mức: NB
\begin{ex}
% ID: L12C4B25-11-TN
% Mức: NB
Phát biểu nào sau đây đúng về tổng hợp cấu trúc hạt nhân và đồng vị?
\choice
{\True Bài tổng hợp cấu trúc cần xác định đúng $Z,N,A$, quan hệ đồng vị và trung bình có trọng số khi biết tỉ lệ đồng vị.}
{Các đồng vị của cùng nguyên tố có số proton khác nhau.}
{Có thể bỏ qua điều kiện áp dụng và đơn vị trong mọi trường hợp.}
{Kết quả của bài toán không phụ thuộc dữ kiện đã cho.}
\loigiai{
Bài tổng hợp cấu trúc cần xác định đúng $Z,N,A$, quan hệ đồng vị và trung bình có trọng số khi biết tỉ lệ đồng vị. Vì vậy phương án $A$ đúng; các phương án còn lại trái với kiến thức hoặc điều kiện áp dụng.
}
\end{ex}

% ===== Câu 45 =====
% ID: L12C4B25-12-DS
% Mức: TH
\dangbt{Tổng hợp độ hụt khối, năng lượng liên kết và độ bền vững}
\begin{ex}
% ID: L12C4B25-12-DS
% Mức: TH
Xét tổng hợp độ hụt khối, năng lượng liên kết và độ bền vững (tình huống 2). Hãy xác định tính đúng, sai của bốn phát biểu sau.
\choiceTF
{Chỉ cần so sánh năng lượng liên kết toàn phần để kết luận độ bền cho mọi hạt nhân.}
{\True Phải tính độ hụt khối đúng dấu, đổi sang năng lượng liên kết và dùng năng lượng liên kết riêng khi so sánh độ bền.}
{Có thể bỏ qua dữ kiện, đơn vị và ý nghĩa vật lí của kết quả.}
{\True Khi giải cần đọc đúng dữ kiện, dùng hệ đơn vị thống nhất và kiểm tra điều kiện áp dụng.}
\loigiai{
\begin{itemchoice}
\item (Sai) Chỉ cần so sánh năng lượng liên kết toàn phần để kết luận độ bền cho mọi hạt nhân. (Vì): Chỉ cần so sánh năng lượng liên kết toàn phần để kết luận độ bền cho mọi hạt nhân. b) Đúng: Phải tính độ hụt khối đúng dấu, đổi sang năng lượng liên kết và dùng năng lượng liên kết riêng khi so sánh độ bền. c) Sai: Có thể bỏ qua dữ kiện, đơn vị và ý nghĩa vật lí của kết quả. d) Đúng: Khi giải cần đọc đúng dữ kiện, dùng hệ đơn vị thống nhất và kiểm tra điều kiện áp dụng. Kiến thức cốt lõi: Phải tính độ hụt khối đúng dấu, đổi sang năng lượng liên kết và dùng năng lượng liên kết riêng khi so sánh độ bền
\item (Đúng) Phải tính độ hụt khối đúng dấu, đổi sang năng lượng liên kết và dùng năng lượng liên kết riêng khi so sánh độ bền.
\item (Sai) Có thể bỏ qua dữ kiện, đơn vị và ý nghĩa vật lí của kết quả.
\item (Đúng) Khi giải cần đọc đúng dữ kiện, dùng hệ đơn vị thống nhất và kiểm tra điều kiện áp dụng.
\end{itemchoice}
}
\end{ex}

% ===== Câu 46 =====
% ID: L12C4B25-12-TL
% Mức: VDC
\dangbt{Tổng hợp cấu trúc hạt nhân và đồng vị}
\begin{bt}
% ID: L12C4B25-12-TL
% Mức: VDC
So sánh nhận định đúng “Bài tổng hợp cấu trúc cần xác định đúng $Z,N,A$, quan hệ đồng vị và trung bình có trọng số khi biết tỉ lệ đồng vị.” với nhận định sai “Các đồng vị của cùng nguyên tố có số proton khác nhau.”; chỉ rõ căn cứ Vật lí.
\loigiai{
Cần nêu được: Bài tổng hợp cấu trúc cần xác định đúng $Z,N,A$, quan hệ đồng vị và trung bình có trọng số khi biết tỉ lệ đồng vị. Khi vận dụng phải xác định đúng dữ kiện, điều kiện áp dụng, hệ đơn vị và kiểm tra tính hợp lí của kết quả. Nhận định “Các đồng vị của cùng nguyên tố có số proton khác nhau.” sai vì trái với kiến thức trên.
}
\end{bt}

% ===== Câu 47 =====
% ID: L12C4B25-12-TLN
% Mức: VDC
\begin{ex}
% ID: L12C4B25-12-TLN
% Mức: VDC
Trong bốn phát biểu sau về tổng hợp cấu trúc hạt nhân và đồng vị, có bao nhiêu phát biểu đúng? Chỉ ghi một số nguyên. (1) Kết quả cần có ý nghĩa vật lí. (2) Các đồng vị của cùng nguyên tố có số proton khác nhau. (3) Phải dùng đúng định luật cho tình huống. (4) Bài tổng hợp cấu trúc cần xác định đúng $Z,N,A$, quan hệ đồng vị và trung bình có trọng số khi biết tỉ lệ đồng vị.
\shortans{3}
\loigiai{
Đối chiếu từng phát biểu với kiến thức cốt lõi: Bài tổng hợp cấu trúc cần xác định đúng $Z,N,A$, quan hệ đồng vị và trung bình có trọng số khi biết tỉ lệ đồng vị. Có 3 phát biểu đúng.
}
\end{ex}

% ===== Câu 48 =====
% ID: L12C4B25-12-TN
% Mức: NB
\begin{ex}
% ID: L12C4B25-12-TN
% Mức: NB
Nhận định nào phù hợp với kiến thức Vật lí về tổng hợp cấu trúc hạt nhân và đồng vị?
\choice
{Các đồng vị của cùng nguyên tố có số proton khác nhau.}
{\True Bài tổng hợp cấu trúc cần xác định đúng $Z,N,A$, quan hệ đồng vị và trung bình có trọng số khi biết tỉ lệ đồng vị.}
{Có thể bỏ qua điều kiện áp dụng và đơn vị trong mọi trường hợp.}
{Kết quả của bài toán không phụ thuộc dữ kiện đã cho.}
\loigiai{
Bài tổng hợp cấu trúc cần xác định đúng $Z,N,A$, quan hệ đồng vị và trung bình có trọng số khi biết tỉ lệ đồng vị. Vì vậy phương án $B$ đúng; các phương án còn lại trái với kiến thức hoặc điều kiện áp dụng.
}
\end{ex}

% ===== Câu 49 =====
% ID: L12C4B25-13-DS
% Mức: VD
\dangbt{Tổng hợp độ hụt khối, năng lượng liên kết và độ bền vững}
\begin{ex}
% ID: L12C4B25-13-DS
% Mức: VD
Xét tổng hợp độ hụt khối, năng lượng liên kết và độ bền vững (tình huống 3). Hãy xác định tính đúng, sai của bốn phát biểu sau.
\choiceTF
{\True Phải tính độ hụt khối đúng dấu, đổi sang năng lượng liên kết và dùng năng lượng liên kết riêng khi so sánh độ bền.}
{Chỉ cần so sánh năng lượng liên kết toàn phần để kết luận độ bền cho mọi hạt nhân.}
{\True Khi giải cần đọc đúng dữ kiện, dùng hệ đơn vị thống nhất và kiểm tra điều kiện áp dụng.}
{Có thể bỏ qua dữ kiện, đơn vị và ý nghĩa vật lí của kết quả.}
\loigiai{
\begin{itemchoice}
\item (Đúng) Phải tính độ hụt khối đúng dấu, đổi sang năng lượng liên kết và dùng năng lượng liên kết riêng khi so sánh độ bền. (Vì): Phải tính độ hụt khối đúng dấu, đổi sang năng lượng liên kết và dùng năng lượng liên kết riêng khi so sánh độ bền. b) Sai: Chỉ cần so sánh năng lượng liên kết toàn phần để kết luận độ bền cho mọi hạt nhân. c) Đúng: Khi giải cần đọc đúng dữ kiện, dùng hệ đơn vị thống nhất và kiểm tra điều kiện áp dụng. d) Sai: Có thể bỏ qua dữ kiện, đơn vị và ý nghĩa vật lí của kết quả. Kiến thức cốt lõi: Phải tính độ hụt khối đúng dấu, đổi sang năng lượng liên kết và dùng năng lượng liên kết riêng khi so sánh độ bền
\item (Sai) Chỉ cần so sánh năng lượng liên kết toàn phần để kết luận độ bền cho mọi hạt nhân.
\item (Đúng) Khi giải cần đọc đúng dữ kiện, dùng hệ đơn vị thống nhất và kiểm tra điều kiện áp dụng.
\item (Sai) Có thể bỏ qua dữ kiện, đơn vị và ý nghĩa vật lí của kết quả.
\end{itemchoice}
}
\end{ex}

% ===== Câu 50 =====
% ID: L12C4B25-13-TL
% Mức: TH
\begin{bt}
% ID: L12C4B25-13-TL
% Mức: TH
Trình bày kiến thức cốt lõi của tổng hợp độ hụt khối, năng lượng liên kết và độ bền vững và nêu điều kiện áp dụng.
\loigiai{
Cần nêu được: Phải tính độ hụt khối đúng dấu, đổi sang năng lượng liên kết và dùng năng lượng liên kết riêng khi so sánh độ bền. Khi vận dụng phải xác định đúng dữ kiện, điều kiện áp dụng, hệ đơn vị và kiểm tra tính hợp lí của kết quả. Nhận định “Chỉ cần so sánh năng lượng liên kết toàn phần để kết luận độ bền cho mọi hạt nhân.” sai vì trái với kiến thức trên.
}
\end{bt}

% ===== Câu 51 =====
% ID: L12C4B25-13-TLN
% Mức: TH
\begin{ex}
% ID: L12C4B25-13-TLN
% Mức: TH
Trong bốn phát biểu sau về tổng hợp độ hụt khối, năng lượng liên kết và độ bền vững, có bao nhiêu phát biểu đúng? Chỉ ghi một số nguyên. (1) Phải tính độ hụt khối đúng dấu, đổi sang năng lượng liên kết và dùng năng lượng liên kết riêng khi so sánh độ bền. (2) Chỉ cần so sánh năng lượng liên kết toàn phần để kết luận độ bền cho mọi hạt nhân. (3) Cần kiểm tra điều kiện áp dụng trước khi tính toán. (4) Có thể bỏ qua đơn vị của kết quả.
\shortans{2}
\loigiai{
Đối chiếu từng phát biểu với kiến thức cốt lõi: Phải tính độ hụt khối đúng dấu, đổi sang năng lượng liên kết và dùng năng lượng liên kết riêng khi so sánh độ bền. Có 2 phát biểu đúng.
}
\end{ex}

% ===== Câu 52 =====
% ID: L12C4B25-13-TN
% Mức: NB
\dangbt{Tổng hợp cấu trúc hạt nhân và đồng vị}
\begin{ex}
% ID: L12C4B25-13-TN
% Mức: NB
Chọn kết luận chính xác về tổng hợp cấu trúc hạt nhân và đồng vị?
\choice
{Các đồng vị của cùng nguyên tố có số proton khác nhau.}
{Có thể bỏ qua điều kiện áp dụng và đơn vị trong mọi trường hợp.}
{\True Bài tổng hợp cấu trúc cần xác định đúng $Z,N,A$, quan hệ đồng vị và trung bình có trọng số khi biết tỉ lệ đồng vị.}
{Kết quả của bài toán không phụ thuộc dữ kiện đã cho.}
\loigiai{
Bài tổng hợp cấu trúc cần xác định đúng $Z,N,A$, quan hệ đồng vị và trung bình có trọng số khi biết tỉ lệ đồng vị. Vì vậy phương án $C$ đúng; các phương án còn lại trái với kiến thức hoặc điều kiện áp dụng.
}
\end{ex}

% ===== Câu 53 =====
% ID: L12C4B25-14-DS
% Mức: VD
\dangbt{Tổng hợp độ hụt khối, năng lượng liên kết và độ bền vững}
\begin{ex}
% ID: L12C4B25-14-DS
% Mức: VD
Xét tổng hợp độ hụt khối, năng lượng liên kết và độ bền vững (tình huống 4). Hãy xác định tính đúng, sai của bốn phát biểu sau.
\choiceTF
{Chỉ cần so sánh năng lượng liên kết toàn phần để kết luận độ bền cho mọi hạt nhân.}
{\True Phải tính độ hụt khối đúng dấu, đổi sang năng lượng liên kết và dùng năng lượng liên kết riêng khi so sánh độ bền.}
{Có thể bỏ qua dữ kiện, đơn vị và ý nghĩa vật lí của kết quả.}
{\True Khi giải cần đọc đúng dữ kiện, dùng hệ đơn vị thống nhất và kiểm tra điều kiện áp dụng.}
\loigiai{
\begin{itemchoice}
\item (Sai) Chỉ cần so sánh năng lượng liên kết toàn phần để kết luận độ bền cho mọi hạt nhân. (Vì): Chỉ cần so sánh năng lượng liên kết toàn phần để kết luận độ bền cho mọi hạt nhân. b) Đúng: Phải tính độ hụt khối đúng dấu, đổi sang năng lượng liên kết và dùng năng lượng liên kết riêng khi so sánh độ bền. c) Sai: Có thể bỏ qua dữ kiện, đơn vị và ý nghĩa vật lí của kết quả. d) Đúng: Khi giải cần đọc đúng dữ kiện, dùng hệ đơn vị thống nhất và kiểm tra điều kiện áp dụng. Kiến thức cốt lõi: Phải tính độ hụt khối đúng dấu, đổi sang năng lượng liên kết và dùng năng lượng liên kết riêng khi so sánh độ bền
\item (Đúng) Phải tính độ hụt khối đúng dấu, đổi sang năng lượng liên kết và dùng năng lượng liên kết riêng khi so sánh độ bền.
\item (Sai) Có thể bỏ qua dữ kiện, đơn vị và ý nghĩa vật lí của kết quả.
\item (Đúng) Khi giải cần đọc đúng dữ kiện, dùng hệ đơn vị thống nhất và kiểm tra điều kiện áp dụng.
\end{itemchoice}
}
\end{ex}

% ===== Câu 54 =====
% ID: L12C4B25-14-TL
% Mức: TH
\begin{bt}
% ID: L12C4B25-14-TL
% Mức: TH
Giải thích vì sao nhận định sau đúng: “Phải tính độ hụt khối đúng dấu, đổi sang năng lượng liên kết và dùng năng lượng liên kết riêng khi so sánh độ bền.”
\loigiai{
Cần nêu được: Phải tính độ hụt khối đúng dấu, đổi sang năng lượng liên kết và dùng năng lượng liên kết riêng khi so sánh độ bền. Khi vận dụng phải xác định đúng dữ kiện, điều kiện áp dụng, hệ đơn vị và kiểm tra tính hợp lí của kết quả. Nhận định “Chỉ cần so sánh năng lượng liên kết toàn phần để kết luận độ bền cho mọi hạt nhân.” sai vì trái với kiến thức trên.
}
\end{bt}

% ===== Câu 55 =====
% ID: L12C4B25-14-TLN
% Mức: TH
\begin{ex}
% ID: L12C4B25-14-TLN
% Mức: TH
Trong bốn phát biểu sau về tổng hợp độ hụt khối, năng lượng liên kết và độ bền vững, có bao nhiêu phát biểu đúng? Chỉ ghi một số nguyên. (1) Chỉ cần so sánh năng lượng liên kết toàn phần để kết luận độ bền cho mọi hạt nhân. (2) Phải tính độ hụt khối đúng dấu, đổi sang năng lượng liên kết và dùng năng lượng liên kết riêng khi so sánh độ bền. (3) Kết quả phải phù hợp ý nghĩa vật lí. (4) Mọi đại lượng đều có cùng bản chất.
\shortans{2}
\loigiai{
Đối chiếu từng phát biểu với kiến thức cốt lõi: Phải tính độ hụt khối đúng dấu, đổi sang năng lượng liên kết và dùng năng lượng liên kết riêng khi so sánh độ bền. Có 2 phát biểu đúng.
}
\end{ex}

% ===== Câu 56 =====
% ID: L12C4B25-14-TN
% Mức: TH
\dangbt{Tổng hợp cấu trúc hạt nhân và đồng vị}
\begin{ex}
% ID: L12C4B25-14-TN
% Mức: TH
Nội dung nào mô tả đúng nhất về tổng hợp cấu trúc hạt nhân và đồng vị?
\choice
{Các đồng vị của cùng nguyên tố có số proton khác nhau.}
{Có thể bỏ qua điều kiện áp dụng và đơn vị trong mọi trường hợp.}
{Kết quả của bài toán không phụ thuộc dữ kiện đã cho.}
{\True Bài tổng hợp cấu trúc cần xác định đúng $Z,N,A$, quan hệ đồng vị và trung bình có trọng số khi biết tỉ lệ đồng vị.}
\loigiai{
Bài tổng hợp cấu trúc cần xác định đúng $Z,N,A$, quan hệ đồng vị và trung bình có trọng số khi biết tỉ lệ đồng vị. Vì vậy phương án $D$ đúng; các phương án còn lại trái với kiến thức hoặc điều kiện áp dụng.
}
\end{ex}

% ===== Câu 57 =====
% ID: L12C4B25-15-DS
% Mức: VD
\dangbt{Tổng hợp độ hụt khối, năng lượng liên kết và độ bền vững}
\begin{ex}
% ID: L12C4B25-15-DS
% Mức: VD
Xét tổng hợp độ hụt khối, năng lượng liên kết và độ bền vững (tình huống 5). Hãy xác định tính đúng, sai của bốn phát biểu sau.
\choiceTF
{\True Phải tính độ hụt khối đúng dấu, đổi sang năng lượng liên kết và dùng năng lượng liên kết riêng khi so sánh độ bền.}
{Có thể bỏ qua dữ kiện, đơn vị và ý nghĩa vật lí của kết quả.}
{Chỉ cần so sánh năng lượng liên kết toàn phần để kết luận độ bền cho mọi hạt nhân.}
{Có thể bỏ qua dữ kiện, đơn vị và ý nghĩa vật lí của kết quả.}
\loigiai{
\begin{itemchoice}
\item (Đúng) Phải tính độ hụt khối đúng dấu, đổi sang năng lượng liên kết và dùng năng lượng liên kết riêng khi so sánh độ bền. (Vì): Phải tính độ hụt khối đúng dấu, đổi sang năng lượng liên kết và dùng năng lượng liên kết riêng khi so sánh độ bền. b) Sai: Có thể bỏ qua dữ kiện, đơn vị và ý nghĩa vật lí của kết quả. c) Sai: Chỉ cần so sánh năng lượng liên kết toàn phần để kết luận độ bền cho mọi hạt nhân. d) Sai: Có thể bỏ qua dữ kiện, đơn vị và ý nghĩa vật lí của kết quả. Kiến thức cốt lõi: Phải tính độ hụt khối đúng dấu, đổi sang năng lượng liên kết và dùng năng lượng liên kết riêng khi so sánh độ bền
\item (Sai) Có thể bỏ qua dữ kiện, đơn vị và ý nghĩa vật lí của kết quả.
\item (Sai) Chỉ cần so sánh năng lượng liên kết toàn phần để kết luận độ bền cho mọi hạt nhân.
\item (Sai) Có thể bỏ qua dữ kiện, đơn vị và ý nghĩa vật lí của kết quả.
\end{itemchoice}
}
\end{ex}

% ===== Câu 58 =====
% ID: L12C4B25-15-TL
% Mức: VD
\begin{bt}
% ID: L12C4B25-15-TL
% Mức: VD
Phân tích sai lầm trong nhận định: “Chỉ cần so sánh năng lượng liên kết toàn phần để kết luận độ bền cho mọi hạt nhân.”
\loigiai{
Cần nêu được: Phải tính độ hụt khối đúng dấu, đổi sang năng lượng liên kết và dùng năng lượng liên kết riêng khi so sánh độ bền. Khi vận dụng phải xác định đúng dữ kiện, điều kiện áp dụng, hệ đơn vị và kiểm tra tính hợp lí của kết quả. Nhận định “Chỉ cần so sánh năng lượng liên kết toàn phần để kết luận độ bền cho mọi hạt nhân.” sai vì trái với kiến thức trên.
}
\end{bt}

% ===== Câu 59 =====
% ID: L12C4B25-15-TLN
% Mức: TH
\begin{ex}
% ID: L12C4B25-15-TLN
% Mức: TH
Trong bốn phát biểu sau về tổng hợp độ hụt khối, năng lượng liên kết và độ bền vững, có bao nhiêu phát biểu đúng? Chỉ ghi một số nguyên. (1) Cần đổi các đại lượng về hệ đơn vị thống nhất. (2) Phải tính độ hụt khối đúng dấu, đổi sang năng lượng liên kết và dùng năng lượng liên kết riêng khi so sánh độ bền. (3) Chỉ cần so sánh năng lượng liên kết toàn phần để kết luận độ bền cho mọi hạt nhân. (4) Không cần kiểm tra dữ kiện.
\shortans{2}
\loigiai{
Đối chiếu từng phát biểu với kiến thức cốt lõi: Phải tính độ hụt khối đúng dấu, đổi sang năng lượng liên kết và dùng năng lượng liên kết riêng khi so sánh độ bền. Có 2 phát biểu đúng.
}
\end{ex}

% ===== Câu 60 =====
% ID: L12C4B25-15-TN
% Mức: TH
\dangbt{Tổng hợp cấu trúc hạt nhân và đồng vị}
\begin{ex}
% ID: L12C4B25-15-TN
% Mức: TH
Khi phân tích, phát biểu nào đúng về tổng hợp cấu trúc hạt nhân và đồng vị?
\choice
{\True Bài tổng hợp cấu trúc cần xác định đúng $Z,N,A$, quan hệ đồng vị và trung bình có trọng số khi biết tỉ lệ đồng vị.}
{Các đồng vị của cùng nguyên tố có số proton khác nhau.}
{Có thể bỏ qua điều kiện áp dụng và đơn vị trong mọi trường hợp.}
{Kết quả của bài toán không phụ thuộc dữ kiện đã cho.}
\loigiai{
Bài tổng hợp cấu trúc cần xác định đúng $Z,N,A$, quan hệ đồng vị và trung bình có trọng số khi biết tỉ lệ đồng vị. Vì vậy phương án $A$ đúng; các phương án còn lại trái với kiến thức hoặc điều kiện áp dụng.
}
\end{ex}

% ===== Câu 61 =====
% ID: L12C4B25-16-DS
% Mức: TH
\dangbt{Tổng hợp phương trình và năng lượng phản ứng hạt nhân}
\begin{ex}
% ID: L12C4B25-16-DS
% Mức: TH
Xét tổng hợp phương trình và năng lượng phản ứng hạt nhân (tình huống 1). Hãy xác định tính đúng, sai của bốn phát biểu sau.
\choiceTF
{\True Cần bảo toàn $A,Z$, xác định hạt còn thiếu rồi tính $Q$ từ chênh lệch khối lượng trước và sau phản ứng.}
{Phương trình hạt nhân chỉ cần cân bằng số hạt ở hai vế.}
{\True Khi giải cần đọc đúng dữ kiện, dùng hệ đơn vị thống nhất và kiểm tra điều kiện áp dụng.}
{Có thể bỏ qua dữ kiện, đơn vị và ý nghĩa vật lí của kết quả.}
\loigiai{
\begin{itemchoice}
\item (Đúng) Cần bảo toàn $A,Z$, xác định hạt còn thiếu rồi tính $Q$ từ chênh lệch khối lượng trước và sau phản ứng. (Vì): Cần bảo toàn $A,Z$, xác định hạt còn thiếu rồi tính $Q$ từ chênh lệch khối lượng trước và sau phản ứng. b) Sai: Phương trình hạt nhân chỉ cần cân bằng số hạt ở hai vế. c) Đúng: Khi giải cần đọc đúng dữ kiện, dùng hệ đơn vị thống nhất và kiểm tra điều kiện áp dụng. d) Sai: Có thể bỏ qua dữ kiện, đơn vị và ý nghĩa vật lí của kết quả. Kiến thức cốt lõi: Cần bảo toàn $A,Z$, xác định hạt còn thiếu rồi tính $Q$ từ chênh lệch khối lượng trước và sau phản ứng
\item (Sai) Phương trình hạt nhân chỉ cần cân bằng số hạt ở hai vế.
\item (Đúng) Khi giải cần đọc đúng dữ kiện, dùng hệ đơn vị thống nhất và kiểm tra điều kiện áp dụng.
\item (Sai) Có thể bỏ qua dữ kiện, đơn vị và ý nghĩa vật lí của kết quả.
\end{itemchoice}
}
\end{ex}

% ===== Câu 62 =====
% ID: L12C4B25-16-TL
% Mức: VD
\dangbt{Tổng hợp độ hụt khối, năng lượng liên kết và độ bền vững}
\begin{bt}
% ID: L12C4B25-16-TL
% Mức: VD
Đề xuất các bước giải một bài toán thuộc dạng tổng hợp độ hụt khối, năng lượng liên kết và độ bền vững.
\loigiai{
Cần nêu được: Phải tính độ hụt khối đúng dấu, đổi sang năng lượng liên kết và dùng năng lượng liên kết riêng khi so sánh độ bền. Khi vận dụng phải xác định đúng dữ kiện, điều kiện áp dụng, hệ đơn vị và kiểm tra tính hợp lí của kết quả. Nhận định “Chỉ cần so sánh năng lượng liên kết toàn phần để kết luận độ bền cho mọi hạt nhân.” sai vì trái với kiến thức trên.
}
\end{bt}

% ===== Câu 63 =====
% ID: L12C4B25-16-TLN
% Mức: VD
\begin{ex}
% ID: L12C4B25-16-TLN
% Mức: VD
Trong bốn phát biểu sau về tổng hợp độ hụt khối, năng lượng liên kết và độ bền vững, có bao nhiêu phát biểu đúng? Chỉ ghi một số nguyên. (1) Phải tính độ hụt khối đúng dấu, đổi sang năng lượng liên kết và dùng năng lượng liên kết riêng khi so sánh độ bền. (2) Cần đối chiếu kết quả với điều kiện bài toán. (3) Chỉ cần so sánh năng lượng liên kết toàn phần để kết luận độ bền cho mọi hạt nhân. (4) Mọi trường hợp đều cho cùng một kết quả.
\shortans{2}
\loigiai{
Đối chiếu từng phát biểu với kiến thức cốt lõi: Phải tính độ hụt khối đúng dấu, đổi sang năng lượng liên kết và dùng năng lượng liên kết riêng khi so sánh độ bền. Có 2 phát biểu đúng.
}
\end{ex}

% ===== Câu 64 =====
% ID: L12C4B25-16-TN
% Mức: TH
\dangbt{Tổng hợp cấu trúc hạt nhân và đồng vị}
\begin{ex}
% ID: L12C4B25-16-TN
% Mức: TH
Kết luận nào của học sinh là đúng về tổng hợp cấu trúc hạt nhân và đồng vị?
\choice
{Các đồng vị của cùng nguyên tố có số proton khác nhau.}
{\True Bài tổng hợp cấu trúc cần xác định đúng $Z,N,A$, quan hệ đồng vị và trung bình có trọng số khi biết tỉ lệ đồng vị.}
{Có thể bỏ qua điều kiện áp dụng và đơn vị trong mọi trường hợp.}
{Kết quả của bài toán không phụ thuộc dữ kiện đã cho.}
\loigiai{
Bài tổng hợp cấu trúc cần xác định đúng $Z,N,A$, quan hệ đồng vị và trung bình có trọng số khi biết tỉ lệ đồng vị. Vì vậy phương án $B$ đúng; các phương án còn lại trái với kiến thức hoặc điều kiện áp dụng.
}
\end{ex}

% ===== Câu 65 =====
% ID: L12C4B25-17-DS
% Mức: TH
\dangbt{Tổng hợp phương trình và năng lượng phản ứng hạt nhân}
\begin{ex}
% ID: L12C4B25-17-DS
% Mức: TH
Xét tổng hợp phương trình và năng lượng phản ứng hạt nhân (tình huống 2). Hãy xác định tính đúng, sai của bốn phát biểu sau.
\choiceTF
{Phương trình hạt nhân chỉ cần cân bằng số hạt ở hai vế.}
{\True Cần bảo toàn $A,Z$, xác định hạt còn thiếu rồi tính $Q$ từ chênh lệch khối lượng trước và sau phản ứng.}
{Có thể bỏ qua dữ kiện, đơn vị và ý nghĩa vật lí của kết quả.}
{\True Khi giải cần đọc đúng dữ kiện, dùng hệ đơn vị thống nhất và kiểm tra điều kiện áp dụng.}
\loigiai{
\begin{itemchoice}
\item (Sai) Phương trình hạt nhân chỉ cần cân bằng số hạt ở hai vế. (Vì): Phương trình hạt nhân chỉ cần cân bằng số hạt ở hai vế. b) Đúng: Cần bảo toàn $A,Z$, xác định hạt còn thiếu rồi tính $Q$ từ chênh lệch khối lượng trước và sau phản ứng. c) Sai: Có thể bỏ qua dữ kiện, đơn vị và ý nghĩa vật lí của kết quả. d) Đúng: Khi giải cần đọc đúng dữ kiện, dùng hệ đơn vị thống nhất và kiểm tra điều kiện áp dụng. Kiến thức cốt lõi: Cần bảo toàn $A,Z$, xác định hạt còn thiếu rồi tính $Q$ từ chênh lệch khối lượng trước và sau phản ứng
\item (Đúng) Cần bảo toàn $A,Z$, xác định hạt còn thiếu rồi tính $Q$ từ chênh lệch khối lượng trước và sau phản ứng.
\item (Sai) Có thể bỏ qua dữ kiện, đơn vị và ý nghĩa vật lí của kết quả.
\item (Đúng) Khi giải cần đọc đúng dữ kiện, dùng hệ đơn vị thống nhất và kiểm tra điều kiện áp dụng.
\end{itemchoice}
}
\end{ex}

% ===== Câu 66 =====
% ID: L12C4B25-17-TL
% Mức: VD
\dangbt{Tổng hợp độ hụt khối, năng lượng liên kết và độ bền vững}
\begin{bt}
% ID: L12C4B25-17-TL
% Mức: VD
Nêu một tình huống thực tiễn liên quan đến tổng hợp độ hụt khối, năng lượng liên kết và độ bền vững và giải thích bằng kiến thức Vật lí.
\loigiai{
Cần nêu được: Phải tính độ hụt khối đúng dấu, đổi sang năng lượng liên kết và dùng năng lượng liên kết riêng khi so sánh độ bền. Khi vận dụng phải xác định đúng dữ kiện, điều kiện áp dụng, hệ đơn vị và kiểm tra tính hợp lí của kết quả. Nhận định “Chỉ cần so sánh năng lượng liên kết toàn phần để kết luận độ bền cho mọi hạt nhân.” sai vì trái với kiến thức trên.
}
\end{bt}

% ===== Câu 67 =====
% ID: L12C4B25-17-TLN
% Mức: VD
\begin{ex}
% ID: L12C4B25-17-TLN
% Mức: VD
Trong bốn phát biểu sau về tổng hợp độ hụt khối, năng lượng liên kết và độ bền vững, có bao nhiêu phát biểu đúng? Chỉ ghi một số nguyên. (1) Chỉ cần so sánh năng lượng liên kết toàn phần để kết luận độ bền cho mọi hạt nhân. (2) Cần đọc đúng dữ kiện và giả thiết. (3) Phải tính độ hụt khối đúng dấu, đổi sang năng lượng liên kết và dùng năng lượng liên kết riêng khi so sánh độ bền. (4) Có thể thay số trước khi chọn công thức.
\shortans{2}
\loigiai{
Đối chiếu từng phát biểu với kiến thức cốt lõi: Phải tính độ hụt khối đúng dấu, đổi sang năng lượng liên kết và dùng năng lượng liên kết riêng khi so sánh độ bền. Có 2 phát biểu đúng.
}
\end{ex}

% ===== Câu 68 =====
% ID: L12C4B25-17-TN
% Mức: TH
\dangbt{Tổng hợp cấu trúc hạt nhân và đồng vị}
\begin{ex}
% ID: L12C4B25-17-TN
% Mức: TH
Chọn phát biểu không trái với kiến thức về tổng hợp cấu trúc hạt nhân và đồng vị?
\choice
{Các đồng vị của cùng nguyên tố có số proton khác nhau.}
{Có thể bỏ qua điều kiện áp dụng và đơn vị trong mọi trường hợp.}
{\True Bài tổng hợp cấu trúc cần xác định đúng $Z,N,A$, quan hệ đồng vị và trung bình có trọng số khi biết tỉ lệ đồng vị.}
{Kết quả của bài toán không phụ thuộc dữ kiện đã cho.}
\loigiai{
Bài tổng hợp cấu trúc cần xác định đúng $Z,N,A$, quan hệ đồng vị và trung bình có trọng số khi biết tỉ lệ đồng vị. Vì vậy phương án $C$ đúng; các phương án còn lại trái với kiến thức hoặc điều kiện áp dụng.
}
\end{ex}

% ===== Câu 69 =====
% ID: L12C4B25-18-DS
% Mức: VD
\dangbt{Tổng hợp phương trình và năng lượng phản ứng hạt nhân}
\begin{ex}
% ID: L12C4B25-18-DS
% Mức: VD
Xét tổng hợp phương trình và năng lượng phản ứng hạt nhân (tình huống 3). Hãy xác định tính đúng, sai của bốn phát biểu sau.
\choiceTF
{\True Cần bảo toàn $A,Z$, xác định hạt còn thiếu rồi tính $Q$ từ chênh lệch khối lượng trước và sau phản ứng.}
{Phương trình hạt nhân chỉ cần cân bằng số hạt ở hai vế.}
{\True Khi giải cần đọc đúng dữ kiện, dùng hệ đơn vị thống nhất và kiểm tra điều kiện áp dụng.}
{Có thể bỏ qua dữ kiện, đơn vị và ý nghĩa vật lí của kết quả.}
\loigiai{
\begin{itemchoice}
\item (Đúng) Cần bảo toàn $A,Z$, xác định hạt còn thiếu rồi tính $Q$ từ chênh lệch khối lượng trước và sau phản ứng. (Vì): Cần bảo toàn $A,Z$, xác định hạt còn thiếu rồi tính $Q$ từ chênh lệch khối lượng trước và sau phản ứng. b) Sai: Phương trình hạt nhân chỉ cần cân bằng số hạt ở hai vế. c) Đúng: Khi giải cần đọc đúng dữ kiện, dùng hệ đơn vị thống nhất và kiểm tra điều kiện áp dụng. d) Sai: Có thể bỏ qua dữ kiện, đơn vị và ý nghĩa vật lí của kết quả. Kiến thức cốt lõi: Cần bảo toàn $A,Z$, xác định hạt còn thiếu rồi tính $Q$ từ chênh lệch khối lượng trước và sau phản ứng
\item (Sai) Phương trình hạt nhân chỉ cần cân bằng số hạt ở hai vế.
\item (Đúng) Khi giải cần đọc đúng dữ kiện, dùng hệ đơn vị thống nhất và kiểm tra điều kiện áp dụng.
\item (Sai) Có thể bỏ qua dữ kiện, đơn vị và ý nghĩa vật lí của kết quả.
\end{itemchoice}
}
\end{ex}

% ===== Câu 70 =====
% ID: L12C4B25-18-TL
% Mức: VDC
\dangbt{Tổng hợp độ hụt khối, năng lượng liên kết và độ bền vững}
\begin{bt}
% ID: L12C4B25-18-TL
% Mức: VDC
So sánh nhận định đúng “Phải tính độ hụt khối đúng dấu, đổi sang năng lượng liên kết và dùng năng lượng liên kết riêng khi so sánh độ bền.” với nhận định sai “Chỉ cần so sánh năng lượng liên kết toàn phần để kết luận độ bền cho mọi hạt nhân.”; chỉ rõ căn cứ Vật lí.
\loigiai{
Cần nêu được: Phải tính độ hụt khối đúng dấu, đổi sang năng lượng liên kết và dùng năng lượng liên kết riêng khi so sánh độ bền. Khi vận dụng phải xác định đúng dữ kiện, điều kiện áp dụng, hệ đơn vị và kiểm tra tính hợp lí của kết quả. Nhận định “Chỉ cần so sánh năng lượng liên kết toàn phần để kết luận độ bền cho mọi hạt nhân.” sai vì trái với kiến thức trên.
}
\end{bt}

% ===== Câu 71 =====
% ID: L12C4B25-18-TLN
% Mức: VDC
\begin{ex}
% ID: L12C4B25-18-TLN
% Mức: VDC
Trong bốn phát biểu sau về tổng hợp độ hụt khối, năng lượng liên kết và độ bền vững, có bao nhiêu phát biểu đúng? Chỉ ghi một số nguyên. (1) Kết quả cần có ý nghĩa vật lí. (2) Chỉ cần so sánh năng lượng liên kết toàn phần để kết luận độ bền cho mọi hạt nhân. (3) Phải dùng đúng định luật cho tình huống. (4) Phải tính độ hụt khối đúng dấu, đổi sang năng lượng liên kết và dùng năng lượng liên kết riêng khi so sánh độ bền.
\shortans{3}
\loigiai{
Đối chiếu từng phát biểu với kiến thức cốt lõi: Phải tính độ hụt khối đúng dấu, đổi sang năng lượng liên kết và dùng năng lượng liên kết riêng khi so sánh độ bền. Có 3 phát biểu đúng.
}
\end{ex}

% ===== Câu 72 =====
% ID: L12C4B25-18-TN
% Mức: VD
\dangbt{Tổng hợp cấu trúc hạt nhân và đồng vị}
\begin{ex}
% ID: L12C4B25-18-TN
% Mức: VD
Cơ sở Vật lí đúng để giải bài là gì về tổng hợp cấu trúc hạt nhân và đồng vị?
\choice
{Các đồng vị của cùng nguyên tố có số proton khác nhau.}
{Có thể bỏ qua điều kiện áp dụng và đơn vị trong mọi trường hợp.}
{Kết quả của bài toán không phụ thuộc dữ kiện đã cho.}
{\True Bài tổng hợp cấu trúc cần xác định đúng $Z,N,A$, quan hệ đồng vị và trung bình có trọng số khi biết tỉ lệ đồng vị.}
\loigiai{
Bài tổng hợp cấu trúc cần xác định đúng $Z,N,A$, quan hệ đồng vị và trung bình có trọng số khi biết tỉ lệ đồng vị. Vì vậy phương án $D$ đúng; các phương án còn lại trái với kiến thức hoặc điều kiện áp dụng.
}
\end{ex}

% ===== Câu 73 =====
% ID: L12C4B25-19-DS
% Mức: VD
\dangbt{Tổng hợp phương trình và năng lượng phản ứng hạt nhân}
\begin{ex}
% ID: L12C4B25-19-DS
% Mức: VD
Xét tổng hợp phương trình và năng lượng phản ứng hạt nhân (tình huống 4). Hãy xác định tính đúng, sai của bốn phát biểu sau.
\choiceTF
{Phương trình hạt nhân chỉ cần cân bằng số hạt ở hai vế.}
{\True Cần bảo toàn $A,Z$, xác định hạt còn thiếu rồi tính $Q$ từ chênh lệch khối lượng trước và sau phản ứng.}
{Có thể bỏ qua dữ kiện, đơn vị và ý nghĩa vật lí của kết quả.}
{\True Khi giải cần đọc đúng dữ kiện, dùng hệ đơn vị thống nhất và kiểm tra điều kiện áp dụng.}
\loigiai{
\begin{itemchoice}
\item (Sai) Phương trình hạt nhân chỉ cần cân bằng số hạt ở hai vế. (Vì): Phương trình hạt nhân chỉ cần cân bằng số hạt ở hai vế. b) Đúng: Cần bảo toàn $A,Z$, xác định hạt còn thiếu rồi tính $Q$ từ chênh lệch khối lượng trước và sau phản ứng. c) Sai: Có thể bỏ qua dữ kiện, đơn vị và ý nghĩa vật lí của kết quả. d) Đúng: Khi giải cần đọc đúng dữ kiện, dùng hệ đơn vị thống nhất và kiểm tra điều kiện áp dụng. Kiến thức cốt lõi: Cần bảo toàn $A,Z$, xác định hạt còn thiếu rồi tính $Q$ từ chênh lệch khối lượng trước và sau phản ứng
\item (Đúng) Cần bảo toàn $A,Z$, xác định hạt còn thiếu rồi tính $Q$ từ chênh lệch khối lượng trước và sau phản ứng.
\item (Sai) Có thể bỏ qua dữ kiện, đơn vị và ý nghĩa vật lí của kết quả.
\item (Đúng) Khi giải cần đọc đúng dữ kiện, dùng hệ đơn vị thống nhất và kiểm tra điều kiện áp dụng.
\end{itemchoice}
}
\end{ex}

% ===== Câu 74 =====
% ID: L12C4B25-19-TL
% Mức: TH
\begin{bt}
% ID: L12C4B25-19-TL
% Mức: TH
Trình bày kiến thức cốt lõi của tổng hợp phương trình và năng lượng phản ứng hạt nhân và nêu điều kiện áp dụng.
\loigiai{
Cần nêu được: Cần bảo toàn $A,Z$, xác định hạt còn thiếu rồi tính $Q$ từ chênh lệch khối lượng trước và sau phản ứng. Khi vận dụng phải xác định đúng dữ kiện, điều kiện áp dụng, hệ đơn vị và kiểm tra tính hợp lí của kết quả. Nhận định “Phương trình hạt nhân chỉ cần cân bằng số hạt ở hai vế.” sai vì trái với kiến thức trên.
}
\end{bt}

% ===== Câu 75 =====
% ID: L12C4B25-19-TLN
% Mức: TH
\begin{ex}
% ID: L12C4B25-19-TLN
% Mức: TH
Trong bốn phát biểu sau về tổng hợp phương trình và năng lượng phản ứng hạt nhân, có bao nhiêu phát biểu đúng? Chỉ ghi một số nguyên. (1) Cần bảo toàn $A,Z$, xác định hạt còn thiếu rồi tính $Q$ từ chênh lệch khối lượng trước và sau phản ứng. (2) Phương trình hạt nhân chỉ cần cân bằng số hạt ở hai vế. (3) Cần kiểm tra điều kiện áp dụng trước khi tính toán. (4) Có thể bỏ qua đơn vị của kết quả.
\shortans{2}
\loigiai{
Đối chiếu từng phát biểu với kiến thức cốt lõi: Cần bảo toàn $A,Z$, xác định hạt còn thiếu rồi tính $Q$ từ chênh lệch khối lượng trước và sau phản ứng. Có 2 phát biểu đúng.
}
\end{ex}

% ===== Câu 76 =====
% ID: L12C4B25-19-TN
% Mức: VD
\dangbt{Tổng hợp cấu trúc hạt nhân và đồng vị}
\begin{ex}
% ID: L12C4B25-19-TN
% Mức: VD
Trong một bài thực hành, nhận xét nào đúng về tổng hợp cấu trúc hạt nhân và đồng vị?
\choice
{\True Bài tổng hợp cấu trúc cần xác định đúng $Z,N,A$, quan hệ đồng vị và trung bình có trọng số khi biết tỉ lệ đồng vị.}
{Các đồng vị của cùng nguyên tố có số proton khác nhau.}
{Có thể bỏ qua điều kiện áp dụng và đơn vị trong mọi trường hợp.}
{Kết quả của bài toán không phụ thuộc dữ kiện đã cho.}
\loigiai{
Bài tổng hợp cấu trúc cần xác định đúng $Z,N,A$, quan hệ đồng vị và trung bình có trọng số khi biết tỉ lệ đồng vị. Vì vậy phương án $A$ đúng; các phương án còn lại trái với kiến thức hoặc điều kiện áp dụng.
}
\end{ex}

% ===== Câu 77 =====
% ID: L12C4B25-20-DS
% Mức: VD
\dangbt{Tổng hợp phương trình và năng lượng phản ứng hạt nhân}
\begin{ex}
% ID: L12C4B25-20-DS
% Mức: VD
Xét tổng hợp phương trình và năng lượng phản ứng hạt nhân (tình huống 5). Hãy xác định tính đúng, sai của bốn phát biểu sau.
\choiceTF
{\True Cần bảo toàn $A,Z$, xác định hạt còn thiếu rồi tính $Q$ từ chênh lệch khối lượng trước và sau phản ứng.}
{Có thể bỏ qua dữ kiện, đơn vị và ý nghĩa vật lí của kết quả.}
{Phương trình hạt nhân chỉ cần cân bằng số hạt ở hai vế.}
{Có thể bỏ qua dữ kiện, đơn vị và ý nghĩa vật lí của kết quả.}
\loigiai{
\begin{itemchoice}
\item (Đúng) Cần bảo toàn $A,Z$, xác định hạt còn thiếu rồi tính $Q$ từ chênh lệch khối lượng trước và sau phản ứng. (Vì): Cần bảo toàn $A,Z$, xác định hạt còn thiếu rồi tính $Q$ từ chênh lệch khối lượng trước và sau phản ứng. b) Sai: Có thể bỏ qua dữ kiện, đơn vị và ý nghĩa vật lí của kết quả. c) Sai: Phương trình hạt nhân chỉ cần cân bằng số hạt ở hai vế. d) Sai: Có thể bỏ qua dữ kiện, đơn vị và ý nghĩa vật lí của kết quả. Kiến thức cốt lõi: Cần bảo toàn $A,Z$, xác định hạt còn thiếu rồi tính $Q$ từ chênh lệch khối lượng trước và sau phản ứng
\item (Sai) Có thể bỏ qua dữ kiện, đơn vị và ý nghĩa vật lí của kết quả.
\item (Sai) Phương trình hạt nhân chỉ cần cân bằng số hạt ở hai vế.
\item (Sai) Có thể bỏ qua dữ kiện, đơn vị và ý nghĩa vật lí của kết quả.
\end{itemchoice}
}
\end{ex}

% ===== Câu 78 =====
% ID: L12C4B25-20-TL
% Mức: TH
\begin{bt}
% ID: L12C4B25-20-TL
% Mức: TH
Giải thích vì sao nhận định sau đúng: “Cần bảo toàn $A,Z$, xác định hạt còn thiếu rồi tính $Q$ từ chênh lệch khối lượng trước và sau phản ứng.”
\loigiai{
Cần nêu được: Cần bảo toàn $A,Z$, xác định hạt còn thiếu rồi tính $Q$ từ chênh lệch khối lượng trước và sau phản ứng. Khi vận dụng phải xác định đúng dữ kiện, điều kiện áp dụng, hệ đơn vị và kiểm tra tính hợp lí của kết quả. Nhận định “Phương trình hạt nhân chỉ cần cân bằng số hạt ở hai vế.” sai vì trái với kiến thức trên.
}
\end{bt}

% ===== Câu 79 =====
% ID: L12C4B25-20-TLN
% Mức: TH
\begin{ex}
% ID: L12C4B25-20-TLN
% Mức: TH
Trong bốn phát biểu sau về tổng hợp phương trình và năng lượng phản ứng hạt nhân, có bao nhiêu phát biểu đúng? Chỉ ghi một số nguyên. (1) Phương trình hạt nhân chỉ cần cân bằng số hạt ở hai vế. (2) Cần bảo toàn $A,Z$, xác định hạt còn thiếu rồi tính $Q$ từ chênh lệch khối lượng trước và sau phản ứng. (3) Kết quả phải phù hợp ý nghĩa vật lí. (4) Mọi đại lượng đều có cùng bản chất.
\shortans{2}
\loigiai{
Đối chiếu từng phát biểu với kiến thức cốt lõi: Cần bảo toàn $A,Z$, xác định hạt còn thiếu rồi tính $Q$ từ chênh lệch khối lượng trước và sau phản ứng. Có 2 phát biểu đúng.
}
\end{ex}

% ===== Câu 80 =====
% ID: L12C4B25-20-TN
% Mức: VD
\dangbt{Tổng hợp cấu trúc hạt nhân và đồng vị}
\begin{ex}
% ID: L12C4B25-20-TN
% Mức: VD
Kết luận đúng khi vận dụng là về tổng hợp cấu trúc hạt nhân và đồng vị?
\choice
{Các đồng vị của cùng nguyên tố có số proton khác nhau.}
{\True Bài tổng hợp cấu trúc cần xác định đúng $Z,N,A$, quan hệ đồng vị và trung bình có trọng số khi biết tỉ lệ đồng vị.}
{Có thể bỏ qua điều kiện áp dụng và đơn vị trong mọi trường hợp.}
{Kết quả của bài toán không phụ thuộc dữ kiện đã cho.}
\loigiai{
Bài tổng hợp cấu trúc cần xác định đúng $Z,N,A$, quan hệ đồng vị và trung bình có trọng số khi biết tỉ lệ đồng vị. Vì vậy phương án $B$ đúng; các phương án còn lại trái với kiến thức hoặc điều kiện áp dụng.
}
\end{ex}

% ===== Câu 81 =====
% ID: L12C4B25-21-DS
% Mức: TH
\dangbt{Tổng hợp định luật phóng xạ, độ phóng xạ và đồ thị}
\begin{ex}
% ID: L12C4B25-21-DS
% Mức: TH
Xét tổng hợp định luật phóng xạ, độ phóng xạ và đồ thị (tình huống 1). Hãy xác định tính đúng, sai của bốn phát biểu sau.
\choiceTF
{\True Các đại lượng $N,m,H$ giảm theo hàm mũ; phải phân biệt lượng còn lại với lượng đã phân rã khi đọc đồ thị.}
{Lượng chất đã phân rã cũng giảm theo cùng biểu thức với lượng chất còn lại.}
{\True Khi giải cần đọc đúng dữ kiện, dùng hệ đơn vị thống nhất và kiểm tra điều kiện áp dụng.}
{Có thể bỏ qua dữ kiện, đơn vị và ý nghĩa vật lí của kết quả.}
\loigiai{
\begin{itemchoice}
\item (Đúng) Các đại lượng $N,m,H$ giảm theo hàm mũ; phải phân biệt lượng còn lại với lượng đã phân rã khi đọc đồ thị. (Vì): Các đại lượng $N,m,H$ giảm theo hàm mũ; phải phân biệt lượng còn lại với lượng đã phân rã khi đọc đồ thị. b) Sai: Lượng chất đã phân rã cũng giảm theo cùng biểu thức với lượng chất còn lại. c) Đúng: Khi giải cần đọc đúng dữ kiện, dùng hệ đơn vị thống nhất và kiểm tra điều kiện áp dụng. d) Sai: Có thể bỏ qua dữ kiện, đơn vị và ý nghĩa vật lí của kết quả. Kiến thức cốt lõi: Các đại lượng $N,m,H$ giảm theo hàm mũ; phải phân biệt lượng còn lại với lượng đã phân rã khi đọc đồ thị
\item (Sai) Lượng chất đã phân rã cũng giảm theo cùng biểu thức với lượng chất còn lại.
\item (Đúng) Khi giải cần đọc đúng dữ kiện, dùng hệ đơn vị thống nhất và kiểm tra điều kiện áp dụng.
\item (Sai) Có thể bỏ qua dữ kiện, đơn vị và ý nghĩa vật lí của kết quả.
\end{itemchoice}
}
\end{ex}

% ===== Câu 82 =====
% ID: L12C4B25-21-TL
% Mức: VD
\dangbt{Tổng hợp phương trình và năng lượng phản ứng hạt nhân}
\begin{bt}
% ID: L12C4B25-21-TL
% Mức: VD
Phân tích sai lầm trong nhận định: “Phương trình hạt nhân chỉ cần cân bằng số hạt ở hai vế.”
\loigiai{
Cần nêu được: Cần bảo toàn $A,Z$, xác định hạt còn thiếu rồi tính $Q$ từ chênh lệch khối lượng trước và sau phản ứng. Khi vận dụng phải xác định đúng dữ kiện, điều kiện áp dụng, hệ đơn vị và kiểm tra tính hợp lí của kết quả. Nhận định “Phương trình hạt nhân chỉ cần cân bằng số hạt ở hai vế.” sai vì trái với kiến thức trên.
}
\end{bt}

% ===== Câu 83 =====
% ID: L12C4B25-21-TLN
% Mức: TH
\begin{ex}
% ID: L12C4B25-21-TLN
% Mức: TH
Trong bốn phát biểu sau về tổng hợp phương trình và năng lượng phản ứng hạt nhân, có bao nhiêu phát biểu đúng? Chỉ ghi một số nguyên. (1) Cần đổi các đại lượng về hệ đơn vị thống nhất. (2) Cần bảo toàn $A,Z$, xác định hạt còn thiếu rồi tính $Q$ từ chênh lệch khối lượng trước và sau phản ứng. (3) Phương trình hạt nhân chỉ cần cân bằng số hạt ở hai vế. (4) Không cần kiểm tra dữ kiện.
\shortans{2}
\loigiai{
Đối chiếu từng phát biểu với kiến thức cốt lõi: Cần bảo toàn $A,Z$, xác định hạt còn thiếu rồi tính $Q$ từ chênh lệch khối lượng trước và sau phản ứng. Có 2 phát biểu đúng.
}
\end{ex}

% ===== Câu 84 =====
% ID: L12C4B25-21-TN
% Mức: NB
\dangbt{Tổng hợp độ hụt khối, năng lượng liên kết và độ bền vững}
\begin{ex}
% ID: L12C4B25-21-TN
% Mức: NB
Phát biểu nào sau đây đúng về tổng hợp độ hụt khối, năng lượng liên kết và độ bền vững?
\choice
{\True Phải tính độ hụt khối đúng dấu, đổi sang năng lượng liên kết và dùng năng lượng liên kết riêng khi so sánh độ bền.}
{Chỉ cần so sánh năng lượng liên kết toàn phần để kết luận độ bền cho mọi hạt nhân.}
{Có thể bỏ qua điều kiện áp dụng và đơn vị trong mọi trường hợp.}
{Kết quả của bài toán không phụ thuộc dữ kiện đã cho.}
\loigiai{
Phải tính độ hụt khối đúng dấu, đổi sang năng lượng liên kết và dùng năng lượng liên kết riêng khi so sánh độ bền. Vì vậy phương án $A$ đúng; các phương án còn lại trái với kiến thức hoặc điều kiện áp dụng.
}
\end{ex}

% ===== Câu 85 =====
% ID: L12C4B25-22-DS
% Mức: TH
\dangbt{Tổng hợp định luật phóng xạ, độ phóng xạ và đồ thị}
\begin{ex}
% ID: L12C4B25-22-DS
% Mức: TH
Xét tổng hợp định luật phóng xạ, độ phóng xạ và đồ thị (tình huống 2). Hãy xác định tính đúng, sai của bốn phát biểu sau.
\choiceTF
{Lượng chất đã phân rã cũng giảm theo cùng biểu thức với lượng chất còn lại.}
{\True Các đại lượng $N,m,H$ giảm theo hàm mũ; phải phân biệt lượng còn lại với lượng đã phân rã khi đọc đồ thị.}
{Có thể bỏ qua dữ kiện, đơn vị và ý nghĩa vật lí của kết quả.}
{\True Khi giải cần đọc đúng dữ kiện, dùng hệ đơn vị thống nhất và kiểm tra điều kiện áp dụng.}
\loigiai{
\begin{itemchoice}
\item (Sai) Lượng chất đã phân rã cũng giảm theo cùng biểu thức với lượng chất còn lại. (Vì): Lượng chất đã phân rã cũng giảm theo cùng biểu thức với lượng chất còn lại. b) Đúng: Các đại lượng $N,m,H$ giảm theo hàm mũ; phải phân biệt lượng còn lại với lượng đã phân rã khi đọc đồ thị. c) Sai: Có thể bỏ qua dữ kiện, đơn vị và ý nghĩa vật lí của kết quả. d) Đúng: Khi giải cần đọc đúng dữ kiện, dùng hệ đơn vị thống nhất và kiểm tra điều kiện áp dụng. Kiến thức cốt lõi: Các đại lượng $N,m,H$ giảm theo hàm mũ; phải phân biệt lượng còn lại với lượng đã phân rã khi đọc đồ thị
\item (Đúng) Các đại lượng $N,m,H$ giảm theo hàm mũ; phải phân biệt lượng còn lại với lượng đã phân rã khi đọc đồ thị.
\item (Sai) Có thể bỏ qua dữ kiện, đơn vị và ý nghĩa vật lí của kết quả.
\item (Đúng) Khi giải cần đọc đúng dữ kiện, dùng hệ đơn vị thống nhất và kiểm tra điều kiện áp dụng.
\end{itemchoice}
}
\end{ex}

% ===== Câu 86 =====
% ID: L12C4B25-22-TL
% Mức: VD
\dangbt{Tổng hợp phương trình và năng lượng phản ứng hạt nhân}
\begin{bt}
% ID: L12C4B25-22-TL
% Mức: VD
Đề xuất các bước giải một bài toán thuộc dạng tổng hợp phương trình và năng lượng phản ứng hạt nhân.
\loigiai{
Cần nêu được: Cần bảo toàn $A,Z$, xác định hạt còn thiếu rồi tính $Q$ từ chênh lệch khối lượng trước và sau phản ứng. Khi vận dụng phải xác định đúng dữ kiện, điều kiện áp dụng, hệ đơn vị và kiểm tra tính hợp lí của kết quả. Nhận định “Phương trình hạt nhân chỉ cần cân bằng số hạt ở hai vế.” sai vì trái với kiến thức trên.
}
\end{bt}

% ===== Câu 87 =====
% ID: L12C4B25-22-TLN
% Mức: VD
\begin{ex}
% ID: L12C4B25-22-TLN
% Mức: VD
Trong bốn phát biểu sau về tổng hợp phương trình và năng lượng phản ứng hạt nhân, có bao nhiêu phát biểu đúng? Chỉ ghi một số nguyên. (1) Cần bảo toàn $A,Z$, xác định hạt còn thiếu rồi tính $Q$ từ chênh lệch khối lượng trước và sau phản ứng. (2) Cần đối chiếu kết quả với điều kiện bài toán. (3) Phương trình hạt nhân chỉ cần cân bằng số hạt ở hai vế. (4) Mọi trường hợp đều cho cùng một kết quả.
\shortans{2}
\loigiai{
Đối chiếu từng phát biểu với kiến thức cốt lõi: Cần bảo toàn $A,Z$, xác định hạt còn thiếu rồi tính $Q$ từ chênh lệch khối lượng trước và sau phản ứng. Có 2 phát biểu đúng.
}
\end{ex}

% ===== Câu 88 =====
% ID: L12C4B25-22-TN
% Mức: NB
\dangbt{Tổng hợp độ hụt khối, năng lượng liên kết và độ bền vững}
\begin{ex}
% ID: L12C4B25-22-TN
% Mức: NB
Nhận định nào phù hợp với kiến thức Vật lí về tổng hợp độ hụt khối, năng lượng liên kết và độ bền vững?
\choice
{Chỉ cần so sánh năng lượng liên kết toàn phần để kết luận độ bền cho mọi hạt nhân.}
{\True Phải tính độ hụt khối đúng dấu, đổi sang năng lượng liên kết và dùng năng lượng liên kết riêng khi so sánh độ bền.}
{Có thể bỏ qua điều kiện áp dụng và đơn vị trong mọi trường hợp.}
{Kết quả của bài toán không phụ thuộc dữ kiện đã cho.}
\loigiai{
Phải tính độ hụt khối đúng dấu, đổi sang năng lượng liên kết và dùng năng lượng liên kết riêng khi so sánh độ bền. Vì vậy phương án $B$ đúng; các phương án còn lại trái với kiến thức hoặc điều kiện áp dụng.
}
\end{ex}

% ===== Câu 89 =====
% ID: L12C4B25-23-DS
% Mức: VD
\dangbt{Tổng hợp định luật phóng xạ, độ phóng xạ và đồ thị}
\begin{ex}
% ID: L12C4B25-23-DS
% Mức: VD
Xét tổng hợp định luật phóng xạ, độ phóng xạ và đồ thị (tình huống 3). Hãy xác định tính đúng, sai của bốn phát biểu sau.
\choiceTF
{\True Các đại lượng $N,m,H$ giảm theo hàm mũ; phải phân biệt lượng còn lại với lượng đã phân rã khi đọc đồ thị.}
{Lượng chất đã phân rã cũng giảm theo cùng biểu thức với lượng chất còn lại.}
{\True Khi giải cần đọc đúng dữ kiện, dùng hệ đơn vị thống nhất và kiểm tra điều kiện áp dụng.}
{Có thể bỏ qua dữ kiện, đơn vị và ý nghĩa vật lí của kết quả.}
\loigiai{
\begin{itemchoice}
\item (Đúng) Các đại lượng $N,m,H$ giảm theo hàm mũ; phải phân biệt lượng còn lại với lượng đã phân rã khi đọc đồ thị. (Vì): Các đại lượng $N,m,H$ giảm theo hàm mũ; phải phân biệt lượng còn lại với lượng đã phân rã khi đọc đồ thị. b) Sai: Lượng chất đã phân rã cũng giảm theo cùng biểu thức với lượng chất còn lại. c) Đúng: Khi giải cần đọc đúng dữ kiện, dùng hệ đơn vị thống nhất và kiểm tra điều kiện áp dụng. d) Sai: Có thể bỏ qua dữ kiện, đơn vị và ý nghĩa vật lí của kết quả. Kiến thức cốt lõi: Các đại lượng $N,m,H$ giảm theo hàm mũ; phải phân biệt lượng còn lại với lượng đã phân rã khi đọc đồ thị
\item (Sai) Lượng chất đã phân rã cũng giảm theo cùng biểu thức với lượng chất còn lại.
\item (Đúng) Khi giải cần đọc đúng dữ kiện, dùng hệ đơn vị thống nhất và kiểm tra điều kiện áp dụng.
\item (Sai) Có thể bỏ qua dữ kiện, đơn vị và ý nghĩa vật lí của kết quả.
\end{itemchoice}
}
\end{ex}

% ===== Câu 90 =====
% ID: L12C4B25-23-TL
% Mức: VD
\dangbt{Tổng hợp phương trình và năng lượng phản ứng hạt nhân}
\begin{bt}
% ID: L12C4B25-23-TL
% Mức: VD
Nêu một tình huống thực tiễn liên quan đến tổng hợp phương trình và năng lượng phản ứng hạt nhân và giải thích bằng kiến thức Vật lí.
\loigiai{
Cần nêu được: Cần bảo toàn $A,Z$, xác định hạt còn thiếu rồi tính $Q$ từ chênh lệch khối lượng trước và sau phản ứng. Khi vận dụng phải xác định đúng dữ kiện, điều kiện áp dụng, hệ đơn vị và kiểm tra tính hợp lí của kết quả. Nhận định “Phương trình hạt nhân chỉ cần cân bằng số hạt ở hai vế.” sai vì trái với kiến thức trên.
}
\end{bt}

% ===== Câu 91 =====
% ID: L12C4B25-23-TLN
% Mức: VD
\begin{ex}
% ID: L12C4B25-23-TLN
% Mức: VD
Trong bốn phát biểu sau về tổng hợp phương trình và năng lượng phản ứng hạt nhân, có bao nhiêu phát biểu đúng? Chỉ ghi một số nguyên. (1) Phương trình hạt nhân chỉ cần cân bằng số hạt ở hai vế. (2) Cần đọc đúng dữ kiện và giả thiết. (3) Cần bảo toàn $A,Z$, xác định hạt còn thiếu rồi tính $Q$ từ chênh lệch khối lượng trước và sau phản ứng. (4) Có thể thay số trước khi chọn công thức.
\shortans{2}
\loigiai{
Đối chiếu từng phát biểu với kiến thức cốt lõi: Cần bảo toàn $A,Z$, xác định hạt còn thiếu rồi tính $Q$ từ chênh lệch khối lượng trước và sau phản ứng. Có 2 phát biểu đúng.
}
\end{ex}

% ===== Câu 92 =====
% ID: L12C4B25-23-TN
% Mức: NB
\dangbt{Tổng hợp độ hụt khối, năng lượng liên kết và độ bền vững}
\begin{ex}
% ID: L12C4B25-23-TN
% Mức: NB
Chọn kết luận chính xác về tổng hợp độ hụt khối, năng lượng liên kết và độ bền vững?
\choice
{Chỉ cần so sánh năng lượng liên kết toàn phần để kết luận độ bền cho mọi hạt nhân.}
{Có thể bỏ qua điều kiện áp dụng và đơn vị trong mọi trường hợp.}
{\True Phải tính độ hụt khối đúng dấu, đổi sang năng lượng liên kết và dùng năng lượng liên kết riêng khi so sánh độ bền.}
{Kết quả của bài toán không phụ thuộc dữ kiện đã cho.}
\loigiai{
Phải tính độ hụt khối đúng dấu, đổi sang năng lượng liên kết và dùng năng lượng liên kết riêng khi so sánh độ bền. Vì vậy phương án $C$ đúng; các phương án còn lại trái với kiến thức hoặc điều kiện áp dụng.
}
\end{ex}

% ===== Câu 93 =====
% ID: L12C4B25-24-DS
% Mức: VD
\dangbt{Tổng hợp định luật phóng xạ, độ phóng xạ và đồ thị}
\begin{ex}
% ID: L12C4B25-24-DS
% Mức: VD
Xét tổng hợp định luật phóng xạ, độ phóng xạ và đồ thị (tình huống 4). Hãy xác định tính đúng, sai của bốn phát biểu sau.
\choiceTF
{Lượng chất đã phân rã cũng giảm theo cùng biểu thức với lượng chất còn lại.}
{\True Các đại lượng $N,m,H$ giảm theo hàm mũ; phải phân biệt lượng còn lại với lượng đã phân rã khi đọc đồ thị.}
{Có thể bỏ qua dữ kiện, đơn vị và ý nghĩa vật lí của kết quả.}
{\True Khi giải cần đọc đúng dữ kiện, dùng hệ đơn vị thống nhất và kiểm tra điều kiện áp dụng.}
\loigiai{
\begin{itemchoice}
\item (Sai) Lượng chất đã phân rã cũng giảm theo cùng biểu thức với lượng chất còn lại. (Vì): Lượng chất đã phân rã cũng giảm theo cùng biểu thức với lượng chất còn lại. b) Đúng: Các đại lượng $N,m,H$ giảm theo hàm mũ; phải phân biệt lượng còn lại với lượng đã phân rã khi đọc đồ thị. c) Sai: Có thể bỏ qua dữ kiện, đơn vị và ý nghĩa vật lí của kết quả. d) Đúng: Khi giải cần đọc đúng dữ kiện, dùng hệ đơn vị thống nhất và kiểm tra điều kiện áp dụng. Kiến thức cốt lõi: Các đại lượng $N,m,H$ giảm theo hàm mũ; phải phân biệt lượng còn lại với lượng đã phân rã khi đọc đồ thị
\item (Đúng) Các đại lượng $N,m,H$ giảm theo hàm mũ; phải phân biệt lượng còn lại với lượng đã phân rã khi đọc đồ thị.
\item (Sai) Có thể bỏ qua dữ kiện, đơn vị và ý nghĩa vật lí của kết quả.
\item (Đúng) Khi giải cần đọc đúng dữ kiện, dùng hệ đơn vị thống nhất và kiểm tra điều kiện áp dụng.
\end{itemchoice}
}
\end{ex}

% ===== Câu 94 =====
% ID: L12C4B25-24-TL
% Mức: VDC
\dangbt{Tổng hợp phương trình và năng lượng phản ứng hạt nhân}
\begin{bt}
% ID: L12C4B25-24-TL
% Mức: VDC
So sánh nhận định đúng “Cần bảo toàn $A,Z$, xác định hạt còn thiếu rồi tính $Q$ từ chênh lệch khối lượng trước và sau phản ứng.” với nhận định sai “Phương trình hạt nhân chỉ cần cân bằng số hạt ở hai vế.”; chỉ rõ căn cứ Vật lí.
\loigiai{
Cần nêu được: Cần bảo toàn $A,Z$, xác định hạt còn thiếu rồi tính $Q$ từ chênh lệch khối lượng trước và sau phản ứng. Khi vận dụng phải xác định đúng dữ kiện, điều kiện áp dụng, hệ đơn vị và kiểm tra tính hợp lí của kết quả. Nhận định “Phương trình hạt nhân chỉ cần cân bằng số hạt ở hai vế.” sai vì trái với kiến thức trên.
}
\end{bt}

% ===== Câu 95 =====
% ID: L12C4B25-24-TLN
% Mức: VDC
\begin{ex}
% ID: L12C4B25-24-TLN
% Mức: VDC
Trong bốn phát biểu sau về tổng hợp phương trình và năng lượng phản ứng hạt nhân, có bao nhiêu phát biểu đúng? Chỉ ghi một số nguyên. (1) Kết quả cần có ý nghĩa vật lí. (2) Phương trình hạt nhân chỉ cần cân bằng số hạt ở hai vế. (3) Phải dùng đúng định luật cho tình huống. (4) Cần bảo toàn $A,Z$, xác định hạt còn thiếu rồi tính $Q$ từ chênh lệch khối lượng trước và sau phản ứng.
\shortans{3}
\loigiai{
Đối chiếu từng phát biểu với kiến thức cốt lõi: Cần bảo toàn $A,Z$, xác định hạt còn thiếu rồi tính $Q$ từ chênh lệch khối lượng trước và sau phản ứng. Có 3 phát biểu đúng.
}
\end{ex}

% ===== Câu 96 =====
% ID: L12C4B25-24-TN
% Mức: TH
\dangbt{Tổng hợp độ hụt khối, năng lượng liên kết và độ bền vững}
\begin{ex}
% ID: L12C4B25-24-TN
% Mức: TH
Nội dung nào mô tả đúng nhất về tổng hợp độ hụt khối, năng lượng liên kết và độ bền vững?
\choice
{Chỉ cần so sánh năng lượng liên kết toàn phần để kết luận độ bền cho mọi hạt nhân.}
{Có thể bỏ qua điều kiện áp dụng và đơn vị trong mọi trường hợp.}
{Kết quả của bài toán không phụ thuộc dữ kiện đã cho.}
{\True Phải tính độ hụt khối đúng dấu, đổi sang năng lượng liên kết và dùng năng lượng liên kết riêng khi so sánh độ bền.}
\loigiai{
Phải tính độ hụt khối đúng dấu, đổi sang năng lượng liên kết và dùng năng lượng liên kết riêng khi so sánh độ bền. Vì vậy phương án $D$ đúng; các phương án còn lại trái với kiến thức hoặc điều kiện áp dụng.
}
\end{ex}

% ===== Câu 97 =====
% ID: L12C4B25-25-DS
% Mức: VD
\dangbt{Tổng hợp định luật phóng xạ, độ phóng xạ và đồ thị}
\begin{ex}
% ID: L12C4B25-25-DS
% Mức: VD
Xét tổng hợp định luật phóng xạ, độ phóng xạ và đồ thị (tình huống 5). Hãy xác định tính đúng, sai của bốn phát biểu sau.
\choiceTF
{\True Các đại lượng $N,m,H$ giảm theo hàm mũ; phải phân biệt lượng còn lại với lượng đã phân rã khi đọc đồ thị.}
{Có thể bỏ qua dữ kiện, đơn vị và ý nghĩa vật lí của kết quả.}
{Lượng chất đã phân rã cũng giảm theo cùng biểu thức với lượng chất còn lại.}
{Có thể bỏ qua dữ kiện, đơn vị và ý nghĩa vật lí của kết quả.}
\loigiai{
\begin{itemchoice}
\item (Đúng) Các đại lượng $N,m,H$ giảm theo hàm mũ; phải phân biệt lượng còn lại với lượng đã phân rã khi đọc đồ thị. (Vì): Các đại lượng $N,m,H$ giảm theo hàm mũ; phải phân biệt lượng còn lại với lượng đã phân rã khi đọc đồ thị. b) Sai: Có thể bỏ qua dữ kiện, đơn vị và ý nghĩa vật lí của kết quả. c) Sai: Lượng chất đã phân rã cũng giảm theo cùng biểu thức với lượng chất còn lại. d) Sai: Có thể bỏ qua dữ kiện, đơn vị và ý nghĩa vật lí của kết quả. Kiến thức cốt lõi: Các đại lượng $N,m,H$ giảm theo hàm mũ; phải phân biệt lượng còn lại với lượng đã phân rã khi đọc đồ thị
\item (Sai) Có thể bỏ qua dữ kiện, đơn vị và ý nghĩa vật lí của kết quả.
\item (Sai) Lượng chất đã phân rã cũng giảm theo cùng biểu thức với lượng chất còn lại.
\item (Sai) Có thể bỏ qua dữ kiện, đơn vị và ý nghĩa vật lí của kết quả.
\end{itemchoice}
}
\end{ex}

% ===== Câu 98 =====
% ID: L12C4B25-25-TL
% Mức: TH
\begin{bt}
% ID: L12C4B25-25-TL
% Mức: TH
Trình bày kiến thức cốt lõi của tổng hợp định luật phóng xạ, độ phóng xạ và đồ thị và nêu điều kiện áp dụng.
\loigiai{
Cần nêu được: Các đại lượng $N,m,H$ giảm theo hàm mũ; phải phân biệt lượng còn lại với lượng đã phân rã khi đọc đồ thị. Khi vận dụng phải xác định đúng dữ kiện, điều kiện áp dụng, hệ đơn vị và kiểm tra tính hợp lí của kết quả. Nhận định “Lượng chất đã phân rã cũng giảm theo cùng biểu thức với lượng chất còn lại.” sai vì trái với kiến thức trên.
}
\end{bt}

% ===== Câu 99 =====
% ID: L12C4B25-25-TLN
% Mức: TH
\begin{ex}
% ID: L12C4B25-25-TLN
% Mức: TH
Trong bốn phát biểu sau về tổng hợp định luật phóng xạ, độ phóng xạ và đồ thị, có bao nhiêu phát biểu đúng? Chỉ ghi một số nguyên. (1) Các đại lượng $N,m,H$ giảm theo hàm mũ; phải phân biệt lượng còn lại với lượng đã phân rã khi đọc đồ thị. (2) Lượng chất đã phân rã cũng giảm theo cùng biểu thức với lượng chất còn lại. (3) Cần kiểm tra điều kiện áp dụng trước khi tính toán. (4) Có thể bỏ qua đơn vị của kết quả.
\shortans{2}
\loigiai{
Đối chiếu từng phát biểu với kiến thức cốt lõi: Các đại lượng $N,m,H$ giảm theo hàm mũ; phải phân biệt lượng còn lại với lượng đã phân rã khi đọc đồ thị. Có 2 phát biểu đúng.
}
\end{ex}

% ===== Câu 100 =====
% ID: L12C4B25-25-TN
% Mức: TH
\dangbt{Tổng hợp độ hụt khối, năng lượng liên kết và độ bền vững}
\begin{ex}
% ID: L12C4B25-25-TN
% Mức: TH
Khi phân tích, phát biểu nào đúng về tổng hợp độ hụt khối, năng lượng liên kết và độ bền vững?
\choice
{\True Phải tính độ hụt khối đúng dấu, đổi sang năng lượng liên kết và dùng năng lượng liên kết riêng khi so sánh độ bền.}
{Chỉ cần so sánh năng lượng liên kết toàn phần để kết luận độ bền cho mọi hạt nhân.}
{Có thể bỏ qua điều kiện áp dụng và đơn vị trong mọi trường hợp.}
{Kết quả của bài toán không phụ thuộc dữ kiện đã cho.}
\loigiai{
Phải tính độ hụt khối đúng dấu, đổi sang năng lượng liên kết và dùng năng lượng liên kết riêng khi so sánh độ bền. Vì vậy phương án $A$ đúng; các phương án còn lại trái với kiến thức hoặc điều kiện áp dụng.
}
\end{ex}

% ===== Câu 101 =====
% ID: L12C4B25-26-DS
% Mức: TH
\dangbt{Tổng hợp phân hạch, nhiệt hạch và công nghiệp hạt nhân}
\begin{ex}
% ID: L12C4B25-26-DS
% Mức: TH
Xét tổng hợp phân hạch, nhiệt hạch và công nghiệp hạt nhân (tình huống 1). Hãy xác định tính đúng, sai của bốn phát biểu sau.
\choiceTF
{\True Bài toán năng lượng cần xác định số phản ứng, năng lượng mỗi phản ứng và hiệu suất; phần công nghiệp phải xét kiểm soát neutron và an toàn.}
{Hiệu suất không ảnh hưởng điện năng thu được từ nhà máy.}
{\True Khi giải cần đọc đúng dữ kiện, dùng hệ đơn vị thống nhất và kiểm tra điều kiện áp dụng.}
{Có thể bỏ qua dữ kiện, đơn vị và ý nghĩa vật lí của kết quả.}
\loigiai{
\begin{itemchoice}
\item (Đúng) Bài toán năng lượng cần xác định số phản ứng, năng lượng mỗi phản ứng và hiệu suất; phần công nghiệp phải xét kiểm soát neutron và an toàn. (Vì): Bài toán năng lượng cần xác định số phản ứng, năng lượng mỗi phản ứng và hiệu suất; phần công nghiệp phải xét kiểm soát neutron và an toàn. b) Sai: Hiệu suất không ảnh hưởng điện năng thu được từ nhà máy. c) Đúng: Khi giải cần đọc đúng dữ kiện, dùng hệ đơn vị thống nhất và kiểm tra điều kiện áp dụng. d) Sai: Có thể bỏ qua dữ kiện, đơn vị và ý nghĩa vật lí của kết quả. Kiến thức cốt lõi: Bài toán năng lượng cần xác định số phản ứng, năng lượng mỗi phản ứng và hiệu suất; phần công nghiệp phải xét kiểm soát neutron và an toàn
\item (Sai) Hiệu suất không ảnh hưởng điện năng thu được từ nhà máy.
\item (Đúng) Khi giải cần đọc đúng dữ kiện, dùng hệ đơn vị thống nhất và kiểm tra điều kiện áp dụng.
\item (Sai) Có thể bỏ qua dữ kiện, đơn vị và ý nghĩa vật lí của kết quả.
\end{itemchoice}
}
\end{ex}

% ===== Câu 102 =====
% ID: L12C4B25-26-TL
% Mức: TH
\dangbt{Tổng hợp định luật phóng xạ, độ phóng xạ và đồ thị}
\begin{bt}
% ID: L12C4B25-26-TL
% Mức: TH
Giải thích vì sao nhận định sau đúng: “Các đại lượng $N,m,H$ giảm theo hàm mũ; phải phân biệt lượng còn lại với lượng đã phân rã khi đọc đồ thị.”
\loigiai{
Cần nêu được: Các đại lượng $N,m,H$ giảm theo hàm mũ; phải phân biệt lượng còn lại với lượng đã phân rã khi đọc đồ thị. Khi vận dụng phải xác định đúng dữ kiện, điều kiện áp dụng, hệ đơn vị và kiểm tra tính hợp lí của kết quả. Nhận định “Lượng chất đã phân rã cũng giảm theo cùng biểu thức với lượng chất còn lại.” sai vì trái với kiến thức trên.
}
\end{bt}

% ===== Câu 103 =====
% ID: L12C4B25-26-TLN
% Mức: TH
\begin{ex}
% ID: L12C4B25-26-TLN
% Mức: TH
Trong bốn phát biểu sau về tổng hợp định luật phóng xạ, độ phóng xạ và đồ thị, có bao nhiêu phát biểu đúng? Chỉ ghi một số nguyên. (1) Lượng chất đã phân rã cũng giảm theo cùng biểu thức với lượng chất còn lại. (2) Các đại lượng $N,m,H$ giảm theo hàm mũ; phải phân biệt lượng còn lại với lượng đã phân rã khi đọc đồ thị. (3) Kết quả phải phù hợp ý nghĩa vật lí. (4) Mọi đại lượng đều có cùng bản chất.
\shortans{2}
\loigiai{
Đối chiếu từng phát biểu với kiến thức cốt lõi: Các đại lượng $N,m,H$ giảm theo hàm mũ; phải phân biệt lượng còn lại với lượng đã phân rã khi đọc đồ thị. Có 2 phát biểu đúng.
}
\end{ex}

% ===== Câu 104 =====
% ID: L12C4B25-26-TN
% Mức: TH
\dangbt{Tổng hợp độ hụt khối, năng lượng liên kết và độ bền vững}
\begin{ex}
% ID: L12C4B25-26-TN
% Mức: TH
Kết luận nào của học sinh là đúng về tổng hợp độ hụt khối, năng lượng liên kết và độ bền vững?
\choice
{Chỉ cần so sánh năng lượng liên kết toàn phần để kết luận độ bền cho mọi hạt nhân.}
{\True Phải tính độ hụt khối đúng dấu, đổi sang năng lượng liên kết và dùng năng lượng liên kết riêng khi so sánh độ bền.}
{Có thể bỏ qua điều kiện áp dụng và đơn vị trong mọi trường hợp.}
{Kết quả của bài toán không phụ thuộc dữ kiện đã cho.}
\loigiai{
Phải tính độ hụt khối đúng dấu, đổi sang năng lượng liên kết và dùng năng lượng liên kết riêng khi so sánh độ bền. Vì vậy phương án $B$ đúng; các phương án còn lại trái với kiến thức hoặc điều kiện áp dụng.
}
\end{ex}

% ===== Câu 105 =====
% ID: L12C4B25-27-DS
% Mức: TH
\dangbt{Tổng hợp phân hạch, nhiệt hạch và công nghiệp hạt nhân}
\begin{ex}
% ID: L12C4B25-27-DS
% Mức: TH
Xét tổng hợp phân hạch, nhiệt hạch và công nghiệp hạt nhân (tình huống 2). Hãy xác định tính đúng, sai của bốn phát biểu sau.
\choiceTF
{Hiệu suất không ảnh hưởng điện năng thu được từ nhà máy.}
{\True Bài toán năng lượng cần xác định số phản ứng, năng lượng mỗi phản ứng và hiệu suất; phần công nghiệp phải xét kiểm soát neutron và an toàn.}
{Có thể bỏ qua dữ kiện, đơn vị và ý nghĩa vật lí của kết quả.}
{\True Khi giải cần đọc đúng dữ kiện, dùng hệ đơn vị thống nhất và kiểm tra điều kiện áp dụng.}
\loigiai{
\begin{itemchoice}
\item (Sai) Hiệu suất không ảnh hưởng điện năng thu được từ nhà máy. (Vì): Hiệu suất không ảnh hưởng điện năng thu được từ nhà máy. b) Đúng: Bài toán năng lượng cần xác định số phản ứng, năng lượng mỗi phản ứng và hiệu suất; phần công nghiệp phải xét kiểm soát neutron và an toàn. c) Sai: Có thể bỏ qua dữ kiện, đơn vị và ý nghĩa vật lí của kết quả. d) Đúng: Khi giải cần đọc đúng dữ kiện, dùng hệ đơn vị thống nhất và kiểm tra điều kiện áp dụng. Kiến thức cốt lõi: Bài toán năng lượng cần xác định số phản ứng, năng lượng mỗi phản ứng và hiệu suất; phần công nghiệp phải xét kiểm soát neutron và an toàn
\item (Đúng) Bài toán năng lượng cần xác định số phản ứng, năng lượng mỗi phản ứng và hiệu suất; phần công nghiệp phải xét kiểm soát neutron và an toàn.
\item (Sai) Có thể bỏ qua dữ kiện, đơn vị và ý nghĩa vật lí của kết quả.
\item (Đúng) Khi giải cần đọc đúng dữ kiện, dùng hệ đơn vị thống nhất và kiểm tra điều kiện áp dụng.
\end{itemchoice}
}
\end{ex}

% ===== Câu 106 =====
% ID: L12C4B25-27-TL
% Mức: VD
\dangbt{Tổng hợp định luật phóng xạ, độ phóng xạ và đồ thị}
\begin{bt}
% ID: L12C4B25-27-TL
% Mức: VD
Phân tích sai lầm trong nhận định: “Lượng chất đã phân rã cũng giảm theo cùng biểu thức với lượng chất còn lại.”
\loigiai{
Cần nêu được: Các đại lượng $N,m,H$ giảm theo hàm mũ; phải phân biệt lượng còn lại với lượng đã phân rã khi đọc đồ thị. Khi vận dụng phải xác định đúng dữ kiện, điều kiện áp dụng, hệ đơn vị và kiểm tra tính hợp lí của kết quả. Nhận định “Lượng chất đã phân rã cũng giảm theo cùng biểu thức với lượng chất còn lại.” sai vì trái với kiến thức trên.
}
\end{bt}

% ===== Câu 107 =====
% ID: L12C4B25-27-TLN
% Mức: TH
\begin{ex}
% ID: L12C4B25-27-TLN
% Mức: TH
Trong bốn phát biểu sau về tổng hợp định luật phóng xạ, độ phóng xạ và đồ thị, có bao nhiêu phát biểu đúng? Chỉ ghi một số nguyên. (1) Cần đổi các đại lượng về hệ đơn vị thống nhất. (2) Các đại lượng $N,m,H$ giảm theo hàm mũ; phải phân biệt lượng còn lại với lượng đã phân rã khi đọc đồ thị. (3) Lượng chất đã phân rã cũng giảm theo cùng biểu thức với lượng chất còn lại. (4) Không cần kiểm tra dữ kiện.
\shortans{2}
\loigiai{
Đối chiếu từng phát biểu với kiến thức cốt lõi: Các đại lượng $N,m,H$ giảm theo hàm mũ; phải phân biệt lượng còn lại với lượng đã phân rã khi đọc đồ thị. Có 2 phát biểu đúng.
}
\end{ex}

% ===== Câu 108 =====
% ID: L12C4B25-27-TN
% Mức: TH
\dangbt{Tổng hợp độ hụt khối, năng lượng liên kết và độ bền vững}
\begin{ex}
% ID: L12C4B25-27-TN
% Mức: TH
Chọn phát biểu không trái với kiến thức về tổng hợp độ hụt khối, năng lượng liên kết và độ bền vững?
\choice
{Chỉ cần so sánh năng lượng liên kết toàn phần để kết luận độ bền cho mọi hạt nhân.}
{Có thể bỏ qua điều kiện áp dụng và đơn vị trong mọi trường hợp.}
{\True Phải tính độ hụt khối đúng dấu, đổi sang năng lượng liên kết và dùng năng lượng liên kết riêng khi so sánh độ bền.}
{Kết quả của bài toán không phụ thuộc dữ kiện đã cho.}
\loigiai{
Phải tính độ hụt khối đúng dấu, đổi sang năng lượng liên kết và dùng năng lượng liên kết riêng khi so sánh độ bền. Vì vậy phương án $C$ đúng; các phương án còn lại trái với kiến thức hoặc điều kiện áp dụng.
}
\end{ex}

% ===== Câu 109 =====
% ID: L12C4B25-28-DS
% Mức: VD
\dangbt{Tổng hợp phân hạch, nhiệt hạch và công nghiệp hạt nhân}
\begin{ex}
% ID: L12C4B25-28-DS
% Mức: VD
Xét tổng hợp phân hạch, nhiệt hạch và công nghiệp hạt nhân (tình huống 3). Hãy xác định tính đúng, sai của bốn phát biểu sau.
\choiceTF
{\True Bài toán năng lượng cần xác định số phản ứng, năng lượng mỗi phản ứng và hiệu suất; phần công nghiệp phải xét kiểm soát neutron và an toàn.}
{Hiệu suất không ảnh hưởng điện năng thu được từ nhà máy.}
{\True Khi giải cần đọc đúng dữ kiện, dùng hệ đơn vị thống nhất và kiểm tra điều kiện áp dụng.}
{Có thể bỏ qua dữ kiện, đơn vị và ý nghĩa vật lí của kết quả.}
\loigiai{
\begin{itemchoice}
\item (Đúng) Bài toán năng lượng cần xác định số phản ứng, năng lượng mỗi phản ứng và hiệu suất; phần công nghiệp phải xét kiểm soát neutron và an toàn. (Vì): Bài toán năng lượng cần xác định số phản ứng, năng lượng mỗi phản ứng và hiệu suất; phần công nghiệp phải xét kiểm soát neutron và an toàn. b) Sai: Hiệu suất không ảnh hưởng điện năng thu được từ nhà máy. c) Đúng: Khi giải cần đọc đúng dữ kiện, dùng hệ đơn vị thống nhất và kiểm tra điều kiện áp dụng. d) Sai: Có thể bỏ qua dữ kiện, đơn vị và ý nghĩa vật lí của kết quả. Kiến thức cốt lõi: Bài toán năng lượng cần xác định số phản ứng, năng lượng mỗi phản ứng và hiệu suất; phần công nghiệp phải xét kiểm soát neutron và an toàn
\item (Sai) Hiệu suất không ảnh hưởng điện năng thu được từ nhà máy.
\item (Đúng) Khi giải cần đọc đúng dữ kiện, dùng hệ đơn vị thống nhất và kiểm tra điều kiện áp dụng.
\item (Sai) Có thể bỏ qua dữ kiện, đơn vị và ý nghĩa vật lí của kết quả.
\end{itemchoice}
}
\end{ex}

% ===== Câu 110 =====
% ID: L12C4B25-28-TL
% Mức: VD
\dangbt{Tổng hợp định luật phóng xạ, độ phóng xạ và đồ thị}
\begin{bt}
% ID: L12C4B25-28-TL
% Mức: VD
Đề xuất các bước giải một bài toán thuộc dạng tổng hợp định luật phóng xạ, độ phóng xạ và đồ thị.
\loigiai{
Cần nêu được: Các đại lượng $N,m,H$ giảm theo hàm mũ; phải phân biệt lượng còn lại với lượng đã phân rã khi đọc đồ thị. Khi vận dụng phải xác định đúng dữ kiện, điều kiện áp dụng, hệ đơn vị và kiểm tra tính hợp lí của kết quả. Nhận định “Lượng chất đã phân rã cũng giảm theo cùng biểu thức với lượng chất còn lại.” sai vì trái với kiến thức trên.
}
\end{bt}

% ===== Câu 111 =====
% ID: L12C4B25-28-TLN
% Mức: VD
\begin{ex}
% ID: L12C4B25-28-TLN
% Mức: VD
Trong bốn phát biểu sau về tổng hợp định luật phóng xạ, độ phóng xạ và đồ thị, có bao nhiêu phát biểu đúng? Chỉ ghi một số nguyên. (1) Các đại lượng $N,m,H$ giảm theo hàm mũ; phải phân biệt lượng còn lại với lượng đã phân rã khi đọc đồ thị. (2) Cần đối chiếu kết quả với điều kiện bài toán. (3) Lượng chất đã phân rã cũng giảm theo cùng biểu thức với lượng chất còn lại. (4) Mọi trường hợp đều cho cùng một kết quả.
\shortans{2}
\loigiai{
Đối chiếu từng phát biểu với kiến thức cốt lõi: Các đại lượng $N,m,H$ giảm theo hàm mũ; phải phân biệt lượng còn lại với lượng đã phân rã khi đọc đồ thị. Có 2 phát biểu đúng.
}
\end{ex}

% ===== Câu 112 =====
% ID: L12C4B25-28-TN
% Mức: VD
\dangbt{Tổng hợp độ hụt khối, năng lượng liên kết và độ bền vững}
\begin{ex}
% ID: L12C4B25-28-TN
% Mức: VD
Cơ sở Vật lí đúng để giải bài là gì về tổng hợp độ hụt khối, năng lượng liên kết và độ bền vững?
\choice
{Chỉ cần so sánh năng lượng liên kết toàn phần để kết luận độ bền cho mọi hạt nhân.}
{Có thể bỏ qua điều kiện áp dụng và đơn vị trong mọi trường hợp.}
{Kết quả của bài toán không phụ thuộc dữ kiện đã cho.}
{\True Phải tính độ hụt khối đúng dấu, đổi sang năng lượng liên kết và dùng năng lượng liên kết riêng khi so sánh độ bền.}
\loigiai{
Phải tính độ hụt khối đúng dấu, đổi sang năng lượng liên kết và dùng năng lượng liên kết riêng khi so sánh độ bền. Vì vậy phương án $D$ đúng; các phương án còn lại trái với kiến thức hoặc điều kiện áp dụng.
}
\end{ex}

% ===== Câu 113 =====
% ID: L12C4B25-29-DS
% Mức: VD
\dangbt{Tổng hợp phân hạch, nhiệt hạch và công nghiệp hạt nhân}
\begin{ex}
% ID: L12C4B25-29-DS
% Mức: VD
Xét tổng hợp phân hạch, nhiệt hạch và công nghiệp hạt nhân (tình huống 4). Hãy xác định tính đúng, sai của bốn phát biểu sau.
\choiceTF
{Hiệu suất không ảnh hưởng điện năng thu được từ nhà máy.}
{\True Bài toán năng lượng cần xác định số phản ứng, năng lượng mỗi phản ứng và hiệu suất; phần công nghiệp phải xét kiểm soát neutron và an toàn.}
{Có thể bỏ qua dữ kiện, đơn vị và ý nghĩa vật lí của kết quả.}
{\True Khi giải cần đọc đúng dữ kiện, dùng hệ đơn vị thống nhất và kiểm tra điều kiện áp dụng.}
\loigiai{
\begin{itemchoice}
\item (Sai) Hiệu suất không ảnh hưởng điện năng thu được từ nhà máy. (Vì): Hiệu suất không ảnh hưởng điện năng thu được từ nhà máy. b) Đúng: Bài toán năng lượng cần xác định số phản ứng, năng lượng mỗi phản ứng và hiệu suất; phần công nghiệp phải xét kiểm soát neutron và an toàn. c) Sai: Có thể bỏ qua dữ kiện, đơn vị và ý nghĩa vật lí của kết quả. d) Đúng: Khi giải cần đọc đúng dữ kiện, dùng hệ đơn vị thống nhất và kiểm tra điều kiện áp dụng. Kiến thức cốt lõi: Bài toán năng lượng cần xác định số phản ứng, năng lượng mỗi phản ứng và hiệu suất; phần công nghiệp phải xét kiểm soát neutron và an toàn
\item (Đúng) Bài toán năng lượng cần xác định số phản ứng, năng lượng mỗi phản ứng và hiệu suất; phần công nghiệp phải xét kiểm soát neutron và an toàn.
\item (Sai) Có thể bỏ qua dữ kiện, đơn vị và ý nghĩa vật lí của kết quả.
\item (Đúng) Khi giải cần đọc đúng dữ kiện, dùng hệ đơn vị thống nhất và kiểm tra điều kiện áp dụng.
\end{itemchoice}
}
\end{ex}

% ===== Câu 114 =====
% ID: L12C4B25-29-TL
% Mức: VD
\dangbt{Tổng hợp định luật phóng xạ, độ phóng xạ và đồ thị}
\begin{bt}
% ID: L12C4B25-29-TL
% Mức: VD
Nêu một tình huống thực tiễn liên quan đến tổng hợp định luật phóng xạ, độ phóng xạ và đồ thị và giải thích bằng kiến thức Vật lí.
\loigiai{
Cần nêu được: Các đại lượng $N,m,H$ giảm theo hàm mũ; phải phân biệt lượng còn lại với lượng đã phân rã khi đọc đồ thị. Khi vận dụng phải xác định đúng dữ kiện, điều kiện áp dụng, hệ đơn vị và kiểm tra tính hợp lí của kết quả. Nhận định “Lượng chất đã phân rã cũng giảm theo cùng biểu thức với lượng chất còn lại.” sai vì trái với kiến thức trên.
}
\end{bt}

% ===== Câu 115 =====
% ID: L12C4B25-29-TLN
% Mức: VD
\begin{ex}
% ID: L12C4B25-29-TLN
% Mức: VD
Trong bốn phát biểu sau về tổng hợp định luật phóng xạ, độ phóng xạ và đồ thị, có bao nhiêu phát biểu đúng? Chỉ ghi một số nguyên. (1) Lượng chất đã phân rã cũng giảm theo cùng biểu thức với lượng chất còn lại. (2) Cần đọc đúng dữ kiện và giả thiết. (3) Các đại lượng $N,m,H$ giảm theo hàm mũ; phải phân biệt lượng còn lại với lượng đã phân rã khi đọc đồ thị. (4) Có thể thay số trước khi chọn công thức.
\shortans{2}
\loigiai{
Đối chiếu từng phát biểu với kiến thức cốt lõi: Các đại lượng $N,m,H$ giảm theo hàm mũ; phải phân biệt lượng còn lại với lượng đã phân rã khi đọc đồ thị. Có 2 phát biểu đúng.
}
\end{ex}

% ===== Câu 116 =====
% ID: L12C4B25-29-TN
% Mức: VD
\dangbt{Tổng hợp độ hụt khối, năng lượng liên kết và độ bền vững}
\begin{ex}
% ID: L12C4B25-29-TN
% Mức: VD
Trong một bài thực hành, nhận xét nào đúng về tổng hợp độ hụt khối, năng lượng liên kết và độ bền vững?
\choice
{\True Phải tính độ hụt khối đúng dấu, đổi sang năng lượng liên kết và dùng năng lượng liên kết riêng khi so sánh độ bền.}
{Chỉ cần so sánh năng lượng liên kết toàn phần để kết luận độ bền cho mọi hạt nhân.}
{Có thể bỏ qua điều kiện áp dụng và đơn vị trong mọi trường hợp.}
{Kết quả của bài toán không phụ thuộc dữ kiện đã cho.}
\loigiai{
Phải tính độ hụt khối đúng dấu, đổi sang năng lượng liên kết và dùng năng lượng liên kết riêng khi so sánh độ bền. Vì vậy phương án $A$ đúng; các phương án còn lại trái với kiến thức hoặc điều kiện áp dụng.
}
\end{ex}

% ===== Câu 117 =====
% ID: L12C4B25-30-DS
% Mức: VD
\dangbt{Tổng hợp phân hạch, nhiệt hạch và công nghiệp hạt nhân}
\begin{ex}
% ID: L12C4B25-30-DS
% Mức: VD
Xét tổng hợp phân hạch, nhiệt hạch và công nghiệp hạt nhân (tình huống 5). Hãy xác định tính đúng, sai của bốn phát biểu sau.
\choiceTF
{\True Bài toán năng lượng cần xác định số phản ứng, năng lượng mỗi phản ứng và hiệu suất; phần công nghiệp phải xét kiểm soát neutron và an toàn.}
{Có thể bỏ qua dữ kiện, đơn vị và ý nghĩa vật lí của kết quả.}
{Hiệu suất không ảnh hưởng điện năng thu được từ nhà máy.}
{Có thể bỏ qua dữ kiện, đơn vị và ý nghĩa vật lí của kết quả.}
\loigiai{
\begin{itemchoice}
\item (Đúng) Bài toán năng lượng cần xác định số phản ứng, năng lượng mỗi phản ứng và hiệu suất; phần công nghiệp phải xét kiểm soát neutron và an toàn. (Vì): Bài toán năng lượng cần xác định số phản ứng, năng lượng mỗi phản ứng và hiệu suất; phần công nghiệp phải xét kiểm soát neutron và an toàn. b) Sai: Có thể bỏ qua dữ kiện, đơn vị và ý nghĩa vật lí của kết quả. c) Sai: Hiệu suất không ảnh hưởng điện năng thu được từ nhà máy. d) Sai: Có thể bỏ qua dữ kiện, đơn vị và ý nghĩa vật lí của kết quả. Kiến thức cốt lõi: Bài toán năng lượng cần xác định số phản ứng, năng lượng mỗi phản ứng và hiệu suất; phần công nghiệp phải xét kiểm soát neutron và an toàn
\item (Sai) Có thể bỏ qua dữ kiện, đơn vị và ý nghĩa vật lí của kết quả.
\item (Sai) Hiệu suất không ảnh hưởng điện năng thu được từ nhà máy.
\item (Sai) Có thể bỏ qua dữ kiện, đơn vị và ý nghĩa vật lí của kết quả.
\end{itemchoice}
}
\end{ex}

% ===== Câu 118 =====
% ID: L12C4B25-30-TL
% Mức: VDC
\dangbt{Tổng hợp định luật phóng xạ, độ phóng xạ và đồ thị}
\begin{bt}
% ID: L12C4B25-30-TL
% Mức: VDC
So sánh nhận định đúng “Các đại lượng $N,m,H$ giảm theo hàm mũ; phải phân biệt lượng còn lại với lượng đã phân rã khi đọc đồ thị.” với nhận định sai “Lượng chất đã phân rã cũng giảm theo cùng biểu thức với lượng chất còn lại.”; chỉ rõ căn cứ Vật lí.
\loigiai{
Cần nêu được: Các đại lượng $N,m,H$ giảm theo hàm mũ; phải phân biệt lượng còn lại với lượng đã phân rã khi đọc đồ thị. Khi vận dụng phải xác định đúng dữ kiện, điều kiện áp dụng, hệ đơn vị và kiểm tra tính hợp lí của kết quả. Nhận định “Lượng chất đã phân rã cũng giảm theo cùng biểu thức với lượng chất còn lại.” sai vì trái với kiến thức trên.
}
\end{bt}

% ===== Câu 119 =====
% ID: L12C4B25-30-TLN
% Mức: VDC
\begin{ex}
% ID: L12C4B25-30-TLN
% Mức: VDC
Trong bốn phát biểu sau về tổng hợp định luật phóng xạ, độ phóng xạ và đồ thị, có bao nhiêu phát biểu đúng? Chỉ ghi một số nguyên. (1) Kết quả cần có ý nghĩa vật lí. (2) Lượng chất đã phân rã cũng giảm theo cùng biểu thức với lượng chất còn lại. (3) Phải dùng đúng định luật cho tình huống. (4) Các đại lượng $N,m,H$ giảm theo hàm mũ; phải phân biệt lượng còn lại với lượng đã phân rã khi đọc đồ thị.
\shortans{3}
\loigiai{
Đối chiếu từng phát biểu với kiến thức cốt lõi: Các đại lượng $N,m,H$ giảm theo hàm mũ; phải phân biệt lượng còn lại với lượng đã phân rã khi đọc đồ thị. Có 3 phát biểu đúng.
}
\end{ex}

% ===== Câu 120 =====
% ID: L12C4B25-30-TN
% Mức: VD
\dangbt{Tổng hợp độ hụt khối, năng lượng liên kết và độ bền vững}
\begin{ex}
% ID: L12C4B25-30-TN
% Mức: VD
Kết luận đúng khi vận dụng là về tổng hợp độ hụt khối, năng lượng liên kết và độ bền vững?
\choice
{Chỉ cần so sánh năng lượng liên kết toàn phần để kết luận độ bền cho mọi hạt nhân.}
{\True Phải tính độ hụt khối đúng dấu, đổi sang năng lượng liên kết và dùng năng lượng liên kết riêng khi so sánh độ bền.}
{Có thể bỏ qua điều kiện áp dụng và đơn vị trong mọi trường hợp.}
{Kết quả của bài toán không phụ thuộc dữ kiện đã cho.}
\loigiai{
Phải tính độ hụt khối đúng dấu, đổi sang năng lượng liên kết và dùng năng lượng liên kết riêng khi so sánh độ bền. Vì vậy phương án $B$ đúng; các phương án còn lại trái với kiến thức hoặc điều kiện áp dụng.
}
\end{ex}

% ===== Câu 121 =====
% ID: L12C4B25-31-DS
% Mức: TH
\dangbt{Bài toán vận dụng cao kết hợp nhiều nội dung của Chương IV}
\begin{ex}
% ID: L12C4B25-31-DS
% Mức: TH
Xét bài toán vận dụng cao kết hợp nhiều nội dung của chương iv (tình huống 1). Hãy xác định tính đúng, sai của bốn phát biểu sau.
\choiceTF
{\True Bài vận dụng cao cần tách giai đoạn, dùng đúng bảo toàn, liên hệ khối lượng–năng lượng và định luật phóng xạ, đồng thời kiểm tra đơn vị và ý nghĩa kết quả.}
{Có thể dùng một công thức duy nhất cho mọi bài toán hạt nhân mà không xét điều kiện.}
{\True Khi giải cần đọc đúng dữ kiện, dùng hệ đơn vị thống nhất và kiểm tra điều kiện áp dụng.}
{Có thể bỏ qua dữ kiện, đơn vị và ý nghĩa vật lí của kết quả.}
\loigiai{
\begin{itemchoice}
\item (Đúng) Bài vận dụng cao cần tách giai đoạn, dùng đúng bảo toàn, liên hệ khối lượng–năng lượng và định luật phóng xạ, đồng thời kiểm tra đơn vị và ý nghĩa kết quả. (Vì): Bài vận dụng cao cần tách giai đoạn, dùng đúng bảo toàn, liên hệ khối lượng–năng lượng và định luật phóng xạ, đồng thời kiểm tra đơn vị và ý nghĩa kết quả. b) Sai: Có thể dùng một công thức duy nhất cho mọi bài toán hạt nhân mà không xét điều kiện. c) Đúng: Khi giải cần đọc đúng dữ kiện, dùng hệ đơn vị thống nhất và kiểm tra điều kiện áp dụng. d) Sai: Có thể bỏ qua dữ kiện, đơn vị và ý nghĩa vật lí của kết quả. Kiến thức cốt lõi: Bài vận dụng cao cần tách giai đoạn, dùng đúng bảo toàn, liên hệ khối lượng–năng lượng và định luật phóng xạ, đồng thời kiểm tra đơn vị và ý nghĩa kết quả
\item (Sai) Có thể dùng một công thức duy nhất cho mọi bài toán hạt nhân mà không xét điều kiện.
\item (Đúng) Khi giải cần đọc đúng dữ kiện, dùng hệ đơn vị thống nhất và kiểm tra điều kiện áp dụng.
\item (Sai) Có thể bỏ qua dữ kiện, đơn vị và ý nghĩa vật lí của kết quả.
\end{itemchoice}
}
\end{ex}

% ===== Câu 122 =====
% ID: L12C4B25-31-TL
% Mức: TH
\dangbt{Tổng hợp phân hạch, nhiệt hạch và công nghiệp hạt nhân}
\begin{bt}
% ID: L12C4B25-31-TL
% Mức: TH
Trình bày kiến thức cốt lõi của tổng hợp phân hạch, nhiệt hạch và công nghiệp hạt nhân và nêu điều kiện áp dụng.
\loigiai{
Cần nêu được: Bài toán năng lượng cần xác định số phản ứng, năng lượng mỗi phản ứng và hiệu suất; phần công nghiệp phải xét kiểm soát neutron và an toàn. Khi vận dụng phải xác định đúng dữ kiện, điều kiện áp dụng, hệ đơn vị và kiểm tra tính hợp lí của kết quả. Nhận định “Hiệu suất không ảnh hưởng điện năng thu được từ nhà máy.” sai vì trái với kiến thức trên.
}
\end{bt}

% ===== Câu 123 =====
% ID: L12C4B25-31-TLN
% Mức: TH
\begin{ex}
% ID: L12C4B25-31-TLN
% Mức: TH
Trong bốn phát biểu sau về tổng hợp phân hạch, nhiệt hạch và công nghiệp hạt nhân, có bao nhiêu phát biểu đúng? Chỉ ghi một số nguyên. (1) Bài toán năng lượng cần xác định số phản ứng, năng lượng mỗi phản ứng và hiệu suất; phần công nghiệp phải xét kiểm soát neutron và an toàn. (2) Hiệu suất không ảnh hưởng điện năng thu được từ nhà máy. (3) Cần kiểm tra điều kiện áp dụng trước khi tính toán. (4) Có thể bỏ qua đơn vị của kết quả.
\shortans{2}
\loigiai{
Đối chiếu từng phát biểu với kiến thức cốt lõi: Bài toán năng lượng cần xác định số phản ứng, năng lượng mỗi phản ứng và hiệu suất; phần công nghiệp phải xét kiểm soát neutron và an toàn. Có 2 phát biểu đúng.
}
\end{ex}

% ===== Câu 124 =====
% ID: L12C4B25-31-TN
% Mức: NB
\dangbt{Tổng hợp phương trình và năng lượng phản ứng hạt nhân}
\begin{ex}
% ID: L12C4B25-31-TN
% Mức: NB
Phát biểu nào sau đây đúng về tổng hợp phương trình và năng lượng phản ứng hạt nhân?
\choice
{\True Cần bảo toàn $A,Z$, xác định hạt còn thiếu rồi tính $Q$ từ chênh lệch khối lượng trước và sau phản ứng.}
{Phương trình hạt nhân chỉ cần cân bằng số hạt ở hai vế.}
{Có thể bỏ qua điều kiện áp dụng và đơn vị trong mọi trường hợp.}
{Kết quả của bài toán không phụ thuộc dữ kiện đã cho.}
\loigiai{
Cần bảo toàn $A,Z$, xác định hạt còn thiếu rồi tính $Q$ từ chênh lệch khối lượng trước và sau phản ứng. Vì vậy phương án $A$ đúng; các phương án còn lại trái với kiến thức hoặc điều kiện áp dụng.
}
\end{ex}

% ===== Câu 125 =====
% ID: L12C4B25-32-DS
% Mức: TH
\dangbt{Bài toán vận dụng cao kết hợp nhiều nội dung của Chương IV}
\begin{ex}
% ID: L12C4B25-32-DS
% Mức: TH
Xét bài toán vận dụng cao kết hợp nhiều nội dung của chương iv (tình huống 2). Hãy xác định tính đúng, sai của bốn phát biểu sau.
\choiceTF
{Có thể dùng một công thức duy nhất cho mọi bài toán hạt nhân mà không xét điều kiện.}
{\True Bài vận dụng cao cần tách giai đoạn, dùng đúng bảo toàn, liên hệ khối lượng–năng lượng và định luật phóng xạ, đồng thời kiểm tra đơn vị và ý nghĩa kết quả.}
{Có thể bỏ qua dữ kiện, đơn vị và ý nghĩa vật lí của kết quả.}
{\True Khi giải cần đọc đúng dữ kiện, dùng hệ đơn vị thống nhất và kiểm tra điều kiện áp dụng.}
\loigiai{
\begin{itemchoice}
\item (Sai) Có thể dùng một công thức duy nhất cho mọi bài toán hạt nhân mà không xét điều kiện. (Vì): Có thể dùng một công thức duy nhất cho mọi bài toán hạt nhân mà không xét điều kiện. b) Đúng: Bài vận dụng cao cần tách giai đoạn, dùng đúng bảo toàn, liên hệ khối lượng–năng lượng và định luật phóng xạ, đồng thời kiểm tra đơn vị và ý nghĩa kết quả. c) Sai: Có thể bỏ qua dữ kiện, đơn vị và ý nghĩa vật lí của kết quả. d) Đúng: Khi giải cần đọc đúng dữ kiện, dùng hệ đơn vị thống nhất và kiểm tra điều kiện áp dụng. Kiến thức cốt lõi: Bài vận dụng cao cần tách giai đoạn, dùng đúng bảo toàn, liên hệ khối lượng–năng lượng và định luật phóng xạ, đồng thời kiểm tra đơn vị và ý nghĩa kết quả
\item (Đúng) Bài vận dụng cao cần tách giai đoạn, dùng đúng bảo toàn, liên hệ khối lượng–năng lượng và định luật phóng xạ, đồng thời kiểm tra đơn vị và ý nghĩa kết quả.
\item (Sai) Có thể bỏ qua dữ kiện, đơn vị và ý nghĩa vật lí của kết quả.
\item (Đúng) Khi giải cần đọc đúng dữ kiện, dùng hệ đơn vị thống nhất và kiểm tra điều kiện áp dụng.
\end{itemchoice}
}
\end{ex}

% ===== Câu 126 =====
% ID: L12C4B25-32-TL
% Mức: TH
\dangbt{Tổng hợp phân hạch, nhiệt hạch và công nghiệp hạt nhân}
\begin{bt}
% ID: L12C4B25-32-TL
% Mức: TH
Giải thích vì sao nhận định sau đúng: “Bài toán năng lượng cần xác định số phản ứng, năng lượng mỗi phản ứng và hiệu suất; phần công nghiệp phải xét kiểm soát neutron và an toàn.”
\loigiai{
Cần nêu được: Bài toán năng lượng cần xác định số phản ứng, năng lượng mỗi phản ứng và hiệu suất; phần công nghiệp phải xét kiểm soát neutron và an toàn. Khi vận dụng phải xác định đúng dữ kiện, điều kiện áp dụng, hệ đơn vị và kiểm tra tính hợp lí của kết quả. Nhận định “Hiệu suất không ảnh hưởng điện năng thu được từ nhà máy.” sai vì trái với kiến thức trên.
}
\end{bt}

% ===== Câu 127 =====
% ID: L12C4B25-32-TLN
% Mức: TH
\begin{ex}
% ID: L12C4B25-32-TLN
% Mức: TH
Trong bốn phát biểu sau về tổng hợp phân hạch, nhiệt hạch và công nghiệp hạt nhân, có bao nhiêu phát biểu đúng? Chỉ ghi một số nguyên. (1) Hiệu suất không ảnh hưởng điện năng thu được từ nhà máy. (2) Bài toán năng lượng cần xác định số phản ứng, năng lượng mỗi phản ứng và hiệu suất; phần công nghiệp phải xét kiểm soát neutron và an toàn. (3) Kết quả phải phù hợp ý nghĩa vật lí. (4) Mọi đại lượng đều có cùng bản chất.
\shortans{2}
\loigiai{
Đối chiếu từng phát biểu với kiến thức cốt lõi: Bài toán năng lượng cần xác định số phản ứng, năng lượng mỗi phản ứng và hiệu suất; phần công nghiệp phải xét kiểm soát neutron và an toàn. Có 2 phát biểu đúng.
}
\end{ex}

% ===== Câu 128 =====
% ID: L12C4B25-32-TN
% Mức: NB
\dangbt{Tổng hợp phương trình và năng lượng phản ứng hạt nhân}
\begin{ex}
% ID: L12C4B25-32-TN
% Mức: NB
Nhận định nào phù hợp với kiến thức Vật lí về tổng hợp phương trình và năng lượng phản ứng hạt nhân?
\choice
{Phương trình hạt nhân chỉ cần cân bằng số hạt ở hai vế.}
{\True Cần bảo toàn $A,Z$, xác định hạt còn thiếu rồi tính $Q$ từ chênh lệch khối lượng trước và sau phản ứng.}
{Có thể bỏ qua điều kiện áp dụng và đơn vị trong mọi trường hợp.}
{Kết quả của bài toán không phụ thuộc dữ kiện đã cho.}
\loigiai{
Cần bảo toàn $A,Z$, xác định hạt còn thiếu rồi tính $Q$ từ chênh lệch khối lượng trước và sau phản ứng. Vì vậy phương án $B$ đúng; các phương án còn lại trái với kiến thức hoặc điều kiện áp dụng.
}
\end{ex}

% ===== Câu 129 =====
% ID: L12C4B25-33-DS
% Mức: VD
\dangbt{Bài toán vận dụng cao kết hợp nhiều nội dung của Chương IV}
\begin{ex}
% ID: L12C4B25-33-DS
% Mức: VD
Xét bài toán vận dụng cao kết hợp nhiều nội dung của chương iv (tình huống 3). Hãy xác định tính đúng, sai của bốn phát biểu sau.
\choiceTF
{\True Bài vận dụng cao cần tách giai đoạn, dùng đúng bảo toàn, liên hệ khối lượng–năng lượng và định luật phóng xạ, đồng thời kiểm tra đơn vị và ý nghĩa kết quả.}
{Có thể dùng một công thức duy nhất cho mọi bài toán hạt nhân mà không xét điều kiện.}
{\True Khi giải cần đọc đúng dữ kiện, dùng hệ đơn vị thống nhất và kiểm tra điều kiện áp dụng.}
{Có thể bỏ qua dữ kiện, đơn vị và ý nghĩa vật lí của kết quả.}
\loigiai{
\begin{itemchoice}
\item (Đúng) Bài vận dụng cao cần tách giai đoạn, dùng đúng bảo toàn, liên hệ khối lượng–năng lượng và định luật phóng xạ, đồng thời kiểm tra đơn vị và ý nghĩa kết quả. (Vì): Bài vận dụng cao cần tách giai đoạn, dùng đúng bảo toàn, liên hệ khối lượng–năng lượng và định luật phóng xạ, đồng thời kiểm tra đơn vị và ý nghĩa kết quả. b) Sai: Có thể dùng một công thức duy nhất cho mọi bài toán hạt nhân mà không xét điều kiện. c) Đúng: Khi giải cần đọc đúng dữ kiện, dùng hệ đơn vị thống nhất và kiểm tra điều kiện áp dụng. d) Sai: Có thể bỏ qua dữ kiện, đơn vị và ý nghĩa vật lí của kết quả. Kiến thức cốt lõi: Bài vận dụng cao cần tách giai đoạn, dùng đúng bảo toàn, liên hệ khối lượng–năng lượng và định luật phóng xạ, đồng thời kiểm tra đơn vị và ý nghĩa kết quả
\item (Sai) Có thể dùng một công thức duy nhất cho mọi bài toán hạt nhân mà không xét điều kiện.
\item (Đúng) Khi giải cần đọc đúng dữ kiện, dùng hệ đơn vị thống nhất và kiểm tra điều kiện áp dụng.
\item (Sai) Có thể bỏ qua dữ kiện, đơn vị và ý nghĩa vật lí của kết quả.
\end{itemchoice}
}
\end{ex}

% ===== Câu 130 =====
% ID: L12C4B25-33-TL
% Mức: VD
\dangbt{Tổng hợp phân hạch, nhiệt hạch và công nghiệp hạt nhân}
\begin{bt}
% ID: L12C4B25-33-TL
% Mức: VD
Phân tích sai lầm trong nhận định: “Hiệu suất không ảnh hưởng điện năng thu được từ nhà máy.”
\loigiai{
Cần nêu được: Bài toán năng lượng cần xác định số phản ứng, năng lượng mỗi phản ứng và hiệu suất; phần công nghiệp phải xét kiểm soát neutron và an toàn. Khi vận dụng phải xác định đúng dữ kiện, điều kiện áp dụng, hệ đơn vị và kiểm tra tính hợp lí của kết quả. Nhận định “Hiệu suất không ảnh hưởng điện năng thu được từ nhà máy.” sai vì trái với kiến thức trên.
}
\end{bt}

% ===== Câu 131 =====
% ID: L12C4B25-33-TLN
% Mức: TH
\begin{ex}
% ID: L12C4B25-33-TLN
% Mức: TH
Trong bốn phát biểu sau về tổng hợp phân hạch, nhiệt hạch và công nghiệp hạt nhân, có bao nhiêu phát biểu đúng? Chỉ ghi một số nguyên. (1) Cần đổi các đại lượng về hệ đơn vị thống nhất. (2) Bài toán năng lượng cần xác định số phản ứng, năng lượng mỗi phản ứng và hiệu suất; phần công nghiệp phải xét kiểm soát neutron và an toàn. (3) Hiệu suất không ảnh hưởng điện năng thu được từ nhà máy. (4) Không cần kiểm tra dữ kiện.
\shortans{2}
\loigiai{
Đối chiếu từng phát biểu với kiến thức cốt lõi: Bài toán năng lượng cần xác định số phản ứng, năng lượng mỗi phản ứng và hiệu suất; phần công nghiệp phải xét kiểm soát neutron và an toàn. Có 2 phát biểu đúng.
}
\end{ex}

% ===== Câu 132 =====
% ID: L12C4B25-33-TN
% Mức: NB
\dangbt{Tổng hợp phương trình và năng lượng phản ứng hạt nhân}
\begin{ex}
% ID: L12C4B25-33-TN
% Mức: NB
Chọn kết luận chính xác về tổng hợp phương trình và năng lượng phản ứng hạt nhân?
\choice
{Phương trình hạt nhân chỉ cần cân bằng số hạt ở hai vế.}
{Có thể bỏ qua điều kiện áp dụng và đơn vị trong mọi trường hợp.}
{\True Cần bảo toàn $A,Z$, xác định hạt còn thiếu rồi tính $Q$ từ chênh lệch khối lượng trước và sau phản ứng.}
{Kết quả của bài toán không phụ thuộc dữ kiện đã cho.}
\loigiai{
Cần bảo toàn $A,Z$, xác định hạt còn thiếu rồi tính $Q$ từ chênh lệch khối lượng trước và sau phản ứng. Vì vậy phương án $C$ đúng; các phương án còn lại trái với kiến thức hoặc điều kiện áp dụng.
}
\end{ex}

% ===== Câu 133 =====
% ID: L12C4B25-34-DS
% Mức: VD
\dangbt{Bài toán vận dụng cao kết hợp nhiều nội dung của Chương IV}
\begin{ex}
% ID: L12C4B25-34-DS
% Mức: VD
Xét bài toán vận dụng cao kết hợp nhiều nội dung của chương iv (tình huống 4). Hãy xác định tính đúng, sai của bốn phát biểu sau.
\choiceTF
{Có thể dùng một công thức duy nhất cho mọi bài toán hạt nhân mà không xét điều kiện.}
{\True Bài vận dụng cao cần tách giai đoạn, dùng đúng bảo toàn, liên hệ khối lượng–năng lượng và định luật phóng xạ, đồng thời kiểm tra đơn vị và ý nghĩa kết quả.}
{Có thể bỏ qua dữ kiện, đơn vị và ý nghĩa vật lí của kết quả.}
{\True Khi giải cần đọc đúng dữ kiện, dùng hệ đơn vị thống nhất và kiểm tra điều kiện áp dụng.}
\loigiai{
\begin{itemchoice}
\item (Sai) Có thể dùng một công thức duy nhất cho mọi bài toán hạt nhân mà không xét điều kiện. (Vì): Có thể dùng một công thức duy nhất cho mọi bài toán hạt nhân mà không xét điều kiện. b) Đúng: Bài vận dụng cao cần tách giai đoạn, dùng đúng bảo toàn, liên hệ khối lượng–năng lượng và định luật phóng xạ, đồng thời kiểm tra đơn vị và ý nghĩa kết quả. c) Sai: Có thể bỏ qua dữ kiện, đơn vị và ý nghĩa vật lí của kết quả. d) Đúng: Khi giải cần đọc đúng dữ kiện, dùng hệ đơn vị thống nhất và kiểm tra điều kiện áp dụng. Kiến thức cốt lõi: Bài vận dụng cao cần tách giai đoạn, dùng đúng bảo toàn, liên hệ khối lượng–năng lượng và định luật phóng xạ, đồng thời kiểm tra đơn vị và ý nghĩa kết quả
\item (Đúng) Bài vận dụng cao cần tách giai đoạn, dùng đúng bảo toàn, liên hệ khối lượng–năng lượng và định luật phóng xạ, đồng thời kiểm tra đơn vị và ý nghĩa kết quả.
\item (Sai) Có thể bỏ qua dữ kiện, đơn vị và ý nghĩa vật lí của kết quả.
\item (Đúng) Khi giải cần đọc đúng dữ kiện, dùng hệ đơn vị thống nhất và kiểm tra điều kiện áp dụng.
\end{itemchoice}
}
\end{ex}

% ===== Câu 134 =====
% ID: L12C4B25-34-TL
% Mức: VD
\dangbt{Tổng hợp phân hạch, nhiệt hạch và công nghiệp hạt nhân}
\begin{bt}
% ID: L12C4B25-34-TL
% Mức: VD
Đề xuất các bước giải một bài toán thuộc dạng tổng hợp phân hạch, nhiệt hạch và công nghiệp hạt nhân.
\loigiai{
Cần nêu được: Bài toán năng lượng cần xác định số phản ứng, năng lượng mỗi phản ứng và hiệu suất; phần công nghiệp phải xét kiểm soát neutron và an toàn. Khi vận dụng phải xác định đúng dữ kiện, điều kiện áp dụng, hệ đơn vị và kiểm tra tính hợp lí của kết quả. Nhận định “Hiệu suất không ảnh hưởng điện năng thu được từ nhà máy.” sai vì trái với kiến thức trên.
}
\end{bt}

% ===== Câu 135 =====
% ID: L12C4B25-34-TLN
% Mức: VD
\begin{ex}
% ID: L12C4B25-34-TLN
% Mức: VD
Trong bốn phát biểu sau về tổng hợp phân hạch, nhiệt hạch và công nghiệp hạt nhân, có bao nhiêu phát biểu đúng? Chỉ ghi một số nguyên. (1) Bài toán năng lượng cần xác định số phản ứng, năng lượng mỗi phản ứng và hiệu suất; phần công nghiệp phải xét kiểm soát neutron và an toàn. (2) Cần đối chiếu kết quả với điều kiện bài toán. (3) Hiệu suất không ảnh hưởng điện năng thu được từ nhà máy. (4) Mọi trường hợp đều cho cùng một kết quả.
\shortans{2}
\loigiai{
Đối chiếu từng phát biểu với kiến thức cốt lõi: Bài toán năng lượng cần xác định số phản ứng, năng lượng mỗi phản ứng và hiệu suất; phần công nghiệp phải xét kiểm soát neutron và an toàn. Có 2 phát biểu đúng.
}
\end{ex}

% ===== Câu 136 =====
% ID: L12C4B25-34-TN
% Mức: TH
\dangbt{Tổng hợp phương trình và năng lượng phản ứng hạt nhân}
\begin{ex}
% ID: L12C4B25-34-TN
% Mức: TH
Nội dung nào mô tả đúng nhất về tổng hợp phương trình và năng lượng phản ứng hạt nhân?
\choice
{Phương trình hạt nhân chỉ cần cân bằng số hạt ở hai vế.}
{Có thể bỏ qua điều kiện áp dụng và đơn vị trong mọi trường hợp.}
{Kết quả của bài toán không phụ thuộc dữ kiện đã cho.}
{\True Cần bảo toàn $A,Z$, xác định hạt còn thiếu rồi tính $Q$ từ chênh lệch khối lượng trước và sau phản ứng.}
\loigiai{
Cần bảo toàn $A,Z$, xác định hạt còn thiếu rồi tính $Q$ từ chênh lệch khối lượng trước và sau phản ứng. Vì vậy phương án $D$ đúng; các phương án còn lại trái với kiến thức hoặc điều kiện áp dụng.
}
\end{ex}

% ===== Câu 137 =====
% ID: L12C4B25-35-DS
% Mức: VD
\dangbt{Bài toán vận dụng cao kết hợp nhiều nội dung của Chương IV}
\begin{ex}
% ID: L12C4B25-35-DS
% Mức: VD
Xét bài toán vận dụng cao kết hợp nhiều nội dung của chương iv (tình huống 5). Hãy xác định tính đúng, sai của bốn phát biểu sau.
\choiceTF
{\True Bài vận dụng cao cần tách giai đoạn, dùng đúng bảo toàn, liên hệ khối lượng–năng lượng và định luật phóng xạ, đồng thời kiểm tra đơn vị và ý nghĩa kết quả.}
{Có thể bỏ qua dữ kiện, đơn vị và ý nghĩa vật lí của kết quả.}
{Có thể dùng một công thức duy nhất cho mọi bài toán hạt nhân mà không xét điều kiện.}
{Có thể bỏ qua dữ kiện, đơn vị và ý nghĩa vật lí của kết quả.}
\loigiai{
\begin{itemchoice}
\item (Đúng) Bài vận dụng cao cần tách giai đoạn, dùng đúng bảo toàn, liên hệ khối lượng–năng lượng và định luật phóng xạ, đồng thời kiểm tra đơn vị và ý nghĩa kết quả. (Vì): Bài vận dụng cao cần tách giai đoạn, dùng đúng bảo toàn, liên hệ khối lượng–năng lượng và định luật phóng xạ, đồng thời kiểm tra đơn vị và ý nghĩa kết quả. b) Sai: Có thể bỏ qua dữ kiện, đơn vị và ý nghĩa vật lí của kết quả. c) Sai: Có thể dùng một công thức duy nhất cho mọi bài toán hạt nhân mà không xét điều kiện. d) Sai: Có thể bỏ qua dữ kiện, đơn vị và ý nghĩa vật lí của kết quả. Kiến thức cốt lõi: Bài vận dụng cao cần tách giai đoạn, dùng đúng bảo toàn, liên hệ khối lượng–năng lượng và định luật phóng xạ, đồng thời kiểm tra đơn vị và ý nghĩa kết quả
\item (Sai) Có thể bỏ qua dữ kiện, đơn vị và ý nghĩa vật lí của kết quả.
\item (Sai) Có thể dùng một công thức duy nhất cho mọi bài toán hạt nhân mà không xét điều kiện.
\item (Sai) Có thể bỏ qua dữ kiện, đơn vị và ý nghĩa vật lí của kết quả.
\end{itemchoice}
}
\end{ex}

% ===== Câu 138 =====
% ID: L12C4B25-35-TL
% Mức: VD
\dangbt{Tổng hợp phân hạch, nhiệt hạch và công nghiệp hạt nhân}
\begin{bt}
% ID: L12C4B25-35-TL
% Mức: VD
Nêu một tình huống thực tiễn liên quan đến tổng hợp phân hạch, nhiệt hạch và công nghiệp hạt nhân và giải thích bằng kiến thức Vật lí.
\loigiai{
Cần nêu được: Bài toán năng lượng cần xác định số phản ứng, năng lượng mỗi phản ứng và hiệu suất; phần công nghiệp phải xét kiểm soát neutron và an toàn. Khi vận dụng phải xác định đúng dữ kiện, điều kiện áp dụng, hệ đơn vị và kiểm tra tính hợp lí của kết quả. Nhận định “Hiệu suất không ảnh hưởng điện năng thu được từ nhà máy.” sai vì trái với kiến thức trên.
}
\end{bt}

% ===== Câu 139 =====
% ID: L12C4B25-35-TLN
% Mức: VD
\begin{ex}
% ID: L12C4B25-35-TLN
% Mức: VD
Trong bốn phát biểu sau về tổng hợp phân hạch, nhiệt hạch và công nghiệp hạt nhân, có bao nhiêu phát biểu đúng? Chỉ ghi một số nguyên. (1) Hiệu suất không ảnh hưởng điện năng thu được từ nhà máy. (2) Cần đọc đúng dữ kiện và giả thiết. (3) Bài toán năng lượng cần xác định số phản ứng, năng lượng mỗi phản ứng và hiệu suất; phần công nghiệp phải xét kiểm soát neutron và an toàn. (4) Có thể thay số trước khi chọn công thức.
\shortans{2}
\loigiai{
Đối chiếu từng phát biểu với kiến thức cốt lõi: Bài toán năng lượng cần xác định số phản ứng, năng lượng mỗi phản ứng và hiệu suất; phần công nghiệp phải xét kiểm soát neutron và an toàn. Có 2 phát biểu đúng.
}
\end{ex}

% ===== Câu 140 =====
% ID: L12C4B25-35-TN
% Mức: TH
\dangbt{Tổng hợp phương trình và năng lượng phản ứng hạt nhân}
\begin{ex}
% ID: L12C4B25-35-TN
% Mức: TH
Khi phân tích, phát biểu nào đúng về tổng hợp phương trình và năng lượng phản ứng hạt nhân?
\choice
{\True Cần bảo toàn $A,Z$, xác định hạt còn thiếu rồi tính $Q$ từ chênh lệch khối lượng trước và sau phản ứng.}
{Phương trình hạt nhân chỉ cần cân bằng số hạt ở hai vế.}
{Có thể bỏ qua điều kiện áp dụng và đơn vị trong mọi trường hợp.}
{Kết quả của bài toán không phụ thuộc dữ kiện đã cho.}
\loigiai{
Cần bảo toàn $A,Z$, xác định hạt còn thiếu rồi tính $Q$ từ chênh lệch khối lượng trước và sau phản ứng. Vì vậy phương án $A$ đúng; các phương án còn lại trái với kiến thức hoặc điều kiện áp dụng.
}
\end{ex}

% ===== Câu 141 =====
% ID: L12C4B25-36-TL
% Mức: VDC
\dangbt{Tổng hợp phân hạch, nhiệt hạch và công nghiệp hạt nhân}
\begin{bt}
% ID: L12C4B25-36-TL
% Mức: VDC
So sánh nhận định đúng “Bài toán năng lượng cần xác định số phản ứng, năng lượng mỗi phản ứng và hiệu suất; phần công nghiệp phải xét kiểm soát neutron và an toàn.” với nhận định sai “Hiệu suất không ảnh hưởng điện năng thu được từ nhà máy.”; chỉ rõ căn cứ Vật lí.
\loigiai{
Cần nêu được: Bài toán năng lượng cần xác định số phản ứng, năng lượng mỗi phản ứng và hiệu suất; phần công nghiệp phải xét kiểm soát neutron và an toàn. Khi vận dụng phải xác định đúng dữ kiện, điều kiện áp dụng, hệ đơn vị và kiểm tra tính hợp lí của kết quả. Nhận định “Hiệu suất không ảnh hưởng điện năng thu được từ nhà máy.” sai vì trái với kiến thức trên.
}
\end{bt}

% ===== Câu 142 =====
% ID: L12C4B25-36-TLN
% Mức: VDC
\begin{ex}
% ID: L12C4B25-36-TLN
% Mức: VDC
Trong bốn phát biểu sau về tổng hợp phân hạch, nhiệt hạch và công nghiệp hạt nhân, có bao nhiêu phát biểu đúng? Chỉ ghi một số nguyên. (1) Kết quả cần có ý nghĩa vật lí. (2) Hiệu suất không ảnh hưởng điện năng thu được từ nhà máy. (3) Phải dùng đúng định luật cho tình huống. (4) Bài toán năng lượng cần xác định số phản ứng, năng lượng mỗi phản ứng và hiệu suất; phần công nghiệp phải xét kiểm soát neutron và an toàn.
\shortans{3}
\loigiai{
Đối chiếu từng phát biểu với kiến thức cốt lõi: Bài toán năng lượng cần xác định số phản ứng, năng lượng mỗi phản ứng và hiệu suất; phần công nghiệp phải xét kiểm soát neutron và an toàn. Có 3 phát biểu đúng.
}
\end{ex}

% ===== Câu 143 =====
% ID: L12C4B25-36-TN
% Mức: TH
\dangbt{Tổng hợp phương trình và năng lượng phản ứng hạt nhân}
\begin{ex}
% ID: L12C4B25-36-TN
% Mức: TH
Kết luận nào của học sinh là đúng về tổng hợp phương trình và năng lượng phản ứng hạt nhân?
\choice
{Phương trình hạt nhân chỉ cần cân bằng số hạt ở hai vế.}
{\True Cần bảo toàn $A,Z$, xác định hạt còn thiếu rồi tính $Q$ từ chênh lệch khối lượng trước và sau phản ứng.}
{Có thể bỏ qua điều kiện áp dụng và đơn vị trong mọi trường hợp.}
{Kết quả của bài toán không phụ thuộc dữ kiện đã cho.}
\loigiai{
Cần bảo toàn $A,Z$, xác định hạt còn thiếu rồi tính $Q$ từ chênh lệch khối lượng trước và sau phản ứng. Vì vậy phương án $B$ đúng; các phương án còn lại trái với kiến thức hoặc điều kiện áp dụng.
}
\end{ex}

% ===== Câu 144 =====
% ID: L12C4B25-37-TL
% Mức: TH
\dangbt{Bài toán vận dụng cao kết hợp nhiều nội dung của Chương IV}
\begin{bt}
% ID: L12C4B25-37-TL
% Mức: TH
Trình bày kiến thức cốt lõi của bài toán vận dụng cao kết hợp nhiều nội dung của chương iv và nêu điều kiện áp dụng.
\loigiai{
Cần nêu được: Bài vận dụng cao cần tách giai đoạn, dùng đúng bảo toàn, liên hệ khối lượng–năng lượng và định luật phóng xạ, đồng thời kiểm tra đơn vị và ý nghĩa kết quả. Khi vận dụng phải xác định đúng dữ kiện, điều kiện áp dụng, hệ đơn vị và kiểm tra tính hợp lí của kết quả. Nhận định “Có thể dùng một công thức duy nhất cho mọi bài toán hạt nhân mà không xét điều kiện.” sai vì trái với kiến thức trên.
}
\end{bt}

% ===== Câu 145 =====
% ID: L12C4B25-37-TLN
% Mức: TH
\begin{ex}
% ID: L12C4B25-37-TLN
% Mức: TH
Trong bốn phát biểu sau về bài toán vận dụng cao kết hợp nhiều nội dung của chương iv, có bao nhiêu phát biểu đúng? Chỉ ghi một số nguyên. (1) Bài vận dụng cao cần tách giai đoạn, dùng đúng bảo toàn, liên hệ khối lượng–năng lượng và định luật phóng xạ, đồng thời kiểm tra đơn vị và ý nghĩa kết quả. (2) Có thể dùng một công thức duy nhất cho mọi bài toán hạt nhân mà không xét điều kiện. (3) Cần kiểm tra điều kiện áp dụng trước khi tính toán. (4) Có thể bỏ qua đơn vị của kết quả.
\shortans{2}
\loigiai{
Đối chiếu từng phát biểu với kiến thức cốt lõi: Bài vận dụng cao cần tách giai đoạn, dùng đúng bảo toàn, liên hệ khối lượng–năng lượng và định luật phóng xạ, đồng thời kiểm tra đơn vị và ý nghĩa kết quả. Có 2 phát biểu đúng.
}
\end{ex}

% ===== Câu 146 =====
% ID: L12C4B25-37-TN
% Mức: TH
\dangbt{Tổng hợp phương trình và năng lượng phản ứng hạt nhân}
\begin{ex}
% ID: L12C4B25-37-TN
% Mức: TH
Chọn phát biểu không trái với kiến thức về tổng hợp phương trình và năng lượng phản ứng hạt nhân?
\choice
{Phương trình hạt nhân chỉ cần cân bằng số hạt ở hai vế.}
{Có thể bỏ qua điều kiện áp dụng và đơn vị trong mọi trường hợp.}
{\True Cần bảo toàn $A,Z$, xác định hạt còn thiếu rồi tính $Q$ từ chênh lệch khối lượng trước và sau phản ứng.}
{Kết quả của bài toán không phụ thuộc dữ kiện đã cho.}
\loigiai{
Cần bảo toàn $A,Z$, xác định hạt còn thiếu rồi tính $Q$ từ chênh lệch khối lượng trước và sau phản ứng. Vì vậy phương án $C$ đúng; các phương án còn lại trái với kiến thức hoặc điều kiện áp dụng.
}
\end{ex}

% ===== Câu 147 =====
% ID: L12C4B25-38-TL
% Mức: TH
\dangbt{Bài toán vận dụng cao kết hợp nhiều nội dung của Chương IV}
\begin{bt}
% ID: L12C4B25-38-TL
% Mức: TH
Giải thích vì sao nhận định sau đúng: “Bài vận dụng cao cần tách giai đoạn, dùng đúng bảo toàn, liên hệ khối lượng–năng lượng và định luật phóng xạ, đồng thời kiểm tra đơn vị và ý nghĩa kết quả.”
\loigiai{
Cần nêu được: Bài vận dụng cao cần tách giai đoạn, dùng đúng bảo toàn, liên hệ khối lượng–năng lượng và định luật phóng xạ, đồng thời kiểm tra đơn vị và ý nghĩa kết quả. Khi vận dụng phải xác định đúng dữ kiện, điều kiện áp dụng, hệ đơn vị và kiểm tra tính hợp lí của kết quả. Nhận định “Có thể dùng một công thức duy nhất cho mọi bài toán hạt nhân mà không xét điều kiện.” sai vì trái với kiến thức trên.
}
\end{bt}

% ===== Câu 148 =====
% ID: L12C4B25-38-TLN
% Mức: TH
\begin{ex}
% ID: L12C4B25-38-TLN
% Mức: TH
Trong bốn phát biểu sau về bài toán vận dụng cao kết hợp nhiều nội dung của chương iv, có bao nhiêu phát biểu đúng? Chỉ ghi một số nguyên. (1) Có thể dùng một công thức duy nhất cho mọi bài toán hạt nhân mà không xét điều kiện. (2) Bài vận dụng cao cần tách giai đoạn, dùng đúng bảo toàn, liên hệ khối lượng–năng lượng và định luật phóng xạ, đồng thời kiểm tra đơn vị và ý nghĩa kết quả. (3) Kết quả phải phù hợp ý nghĩa vật lí. (4) Mọi đại lượng đều có cùng bản chất.
\shortans{2}
\loigiai{
Đối chiếu từng phát biểu với kiến thức cốt lõi: Bài vận dụng cao cần tách giai đoạn, dùng đúng bảo toàn, liên hệ khối lượng–năng lượng và định luật phóng xạ, đồng thời kiểm tra đơn vị và ý nghĩa kết quả. Có 2 phát biểu đúng.
}
\end{ex}

% ===== Câu 149 =====
% ID: L12C4B25-38-TN
% Mức: VD
\dangbt{Tổng hợp phương trình và năng lượng phản ứng hạt nhân}
\begin{ex}
% ID: L12C4B25-38-TN
% Mức: VD
Cơ sở Vật lí đúng để giải bài là gì về tổng hợp phương trình và năng lượng phản ứng hạt nhân?
\choice
{Phương trình hạt nhân chỉ cần cân bằng số hạt ở hai vế.}
{Có thể bỏ qua điều kiện áp dụng và đơn vị trong mọi trường hợp.}
{Kết quả của bài toán không phụ thuộc dữ kiện đã cho.}
{\True Cần bảo toàn $A,Z$, xác định hạt còn thiếu rồi tính $Q$ từ chênh lệch khối lượng trước và sau phản ứng.}
\loigiai{
Cần bảo toàn $A,Z$, xác định hạt còn thiếu rồi tính $Q$ từ chênh lệch khối lượng trước và sau phản ứng. Vì vậy phương án $D$ đúng; các phương án còn lại trái với kiến thức hoặc điều kiện áp dụng.
}
\end{ex}

% ===== Câu 150 =====
% ID: L12C4B25-39-TL
% Mức: VD
\dangbt{Bài toán vận dụng cao kết hợp nhiều nội dung của Chương IV}
\begin{bt}
% ID: L12C4B25-39-TL
% Mức: VD
Phân tích sai lầm trong nhận định: “Có thể dùng một công thức duy nhất cho mọi bài toán hạt nhân mà không xét điều kiện.”
\loigiai{
Cần nêu được: Bài vận dụng cao cần tách giai đoạn, dùng đúng bảo toàn, liên hệ khối lượng–năng lượng và định luật phóng xạ, đồng thời kiểm tra đơn vị và ý nghĩa kết quả. Khi vận dụng phải xác định đúng dữ kiện, điều kiện áp dụng, hệ đơn vị và kiểm tra tính hợp lí của kết quả. Nhận định “Có thể dùng một công thức duy nhất cho mọi bài toán hạt nhân mà không xét điều kiện.” sai vì trái với kiến thức trên.
}
\end{bt}

% ===== Câu 151 =====
% ID: L12C4B25-39-TLN
% Mức: TH
\begin{ex}
% ID: L12C4B25-39-TLN
% Mức: TH
Trong bốn phát biểu sau về bài toán vận dụng cao kết hợp nhiều nội dung của chương iv, có bao nhiêu phát biểu đúng? Chỉ ghi một số nguyên. (1) Cần đổi các đại lượng về hệ đơn vị thống nhất. (2) Bài vận dụng cao cần tách giai đoạn, dùng đúng bảo toàn, liên hệ khối lượng–năng lượng và định luật phóng xạ, đồng thời kiểm tra đơn vị và ý nghĩa kết quả. (3) Có thể dùng một công thức duy nhất cho mọi bài toán hạt nhân mà không xét điều kiện. (4) Không cần kiểm tra dữ kiện.
\shortans{2}
\loigiai{
Đối chiếu từng phát biểu với kiến thức cốt lõi: Bài vận dụng cao cần tách giai đoạn, dùng đúng bảo toàn, liên hệ khối lượng–năng lượng và định luật phóng xạ, đồng thời kiểm tra đơn vị và ý nghĩa kết quả. Có 2 phát biểu đúng.
}
\end{ex}

% ===== Câu 152 =====
% ID: L12C4B25-39-TN
% Mức: VD
\dangbt{Tổng hợp phương trình và năng lượng phản ứng hạt nhân}
\begin{ex}
% ID: L12C4B25-39-TN
% Mức: VD
Trong một bài thực hành, nhận xét nào đúng về tổng hợp phương trình và năng lượng phản ứng hạt nhân?
\choice
{\True Cần bảo toàn $A,Z$, xác định hạt còn thiếu rồi tính $Q$ từ chênh lệch khối lượng trước và sau phản ứng.}
{Phương trình hạt nhân chỉ cần cân bằng số hạt ở hai vế.}
{Có thể bỏ qua điều kiện áp dụng và đơn vị trong mọi trường hợp.}
{Kết quả của bài toán không phụ thuộc dữ kiện đã cho.}
\loigiai{
Cần bảo toàn $A,Z$, xác định hạt còn thiếu rồi tính $Q$ từ chênh lệch khối lượng trước và sau phản ứng. Vì vậy phương án $A$ đúng; các phương án còn lại trái với kiến thức hoặc điều kiện áp dụng.
}
\end{ex}

% ===== Câu 153 =====
% ID: L12C4B25-40-TL
% Mức: VD
\dangbt{Bài toán vận dụng cao kết hợp nhiều nội dung của Chương IV}
\begin{bt}
% ID: L12C4B25-40-TL
% Mức: VD
Đề xuất các bước giải một bài toán thuộc dạng bài toán vận dụng cao kết hợp nhiều nội dung của chương iv.
\loigiai{
Cần nêu được: Bài vận dụng cao cần tách giai đoạn, dùng đúng bảo toàn, liên hệ khối lượng–năng lượng và định luật phóng xạ, đồng thời kiểm tra đơn vị và ý nghĩa kết quả. Khi vận dụng phải xác định đúng dữ kiện, điều kiện áp dụng, hệ đơn vị và kiểm tra tính hợp lí của kết quả. Nhận định “Có thể dùng một công thức duy nhất cho mọi bài toán hạt nhân mà không xét điều kiện.” sai vì trái với kiến thức trên.
}
\end{bt}

% ===== Câu 154 =====
% ID: L12C4B25-40-TLN
% Mức: VD
\begin{ex}
% ID: L12C4B25-40-TLN
% Mức: VD
Trong bốn phát biểu sau về bài toán vận dụng cao kết hợp nhiều nội dung của chương iv, có bao nhiêu phát biểu đúng? Chỉ ghi một số nguyên. (1) Bài vận dụng cao cần tách giai đoạn, dùng đúng bảo toàn, liên hệ khối lượng–năng lượng và định luật phóng xạ, đồng thời kiểm tra đơn vị và ý nghĩa kết quả. (2) Cần đối chiếu kết quả với điều kiện bài toán. (3) Có thể dùng một công thức duy nhất cho mọi bài toán hạt nhân mà không xét điều kiện. (4) Mọi trường hợp đều cho cùng một kết quả.
\shortans{2}
\loigiai{
Đối chiếu từng phát biểu với kiến thức cốt lõi: Bài vận dụng cao cần tách giai đoạn, dùng đúng bảo toàn, liên hệ khối lượng–năng lượng và định luật phóng xạ, đồng thời kiểm tra đơn vị và ý nghĩa kết quả. Có 2 phát biểu đúng.
}
\end{ex}

% ===== Câu 155 =====
% ID: L12C4B25-40-TN
% Mức: VD
\dangbt{Tổng hợp phương trình và năng lượng phản ứng hạt nhân}
\begin{ex}
% ID: L12C4B25-40-TN
% Mức: VD
Kết luận đúng khi vận dụng là về tổng hợp phương trình và năng lượng phản ứng hạt nhân?
\choice
{Phương trình hạt nhân chỉ cần cân bằng số hạt ở hai vế.}
{\True Cần bảo toàn $A,Z$, xác định hạt còn thiếu rồi tính $Q$ từ chênh lệch khối lượng trước và sau phản ứng.}
{Có thể bỏ qua điều kiện áp dụng và đơn vị trong mọi trường hợp.}
{Kết quả của bài toán không phụ thuộc dữ kiện đã cho.}
\loigiai{
Cần bảo toàn $A,Z$, xác định hạt còn thiếu rồi tính $Q$ từ chênh lệch khối lượng trước và sau phản ứng. Vì vậy phương án $B$ đúng; các phương án còn lại trái với kiến thức hoặc điều kiện áp dụng.
}
\end{ex}

% ===== Câu 156 =====
% ID: L12C4B25-41-TL
% Mức: VD
\dangbt{Bài toán vận dụng cao kết hợp nhiều nội dung của Chương IV}
\begin{bt}
% ID: L12C4B25-41-TL
% Mức: VD
Nêu một tình huống thực tiễn liên quan đến bài toán vận dụng cao kết hợp nhiều nội dung của chương iv và giải thích bằng kiến thức Vật lí.
\loigiai{
Cần nêu được: Bài vận dụng cao cần tách giai đoạn, dùng đúng bảo toàn, liên hệ khối lượng–năng lượng và định luật phóng xạ, đồng thời kiểm tra đơn vị và ý nghĩa kết quả. Khi vận dụng phải xác định đúng dữ kiện, điều kiện áp dụng, hệ đơn vị và kiểm tra tính hợp lí của kết quả. Nhận định “Có thể dùng một công thức duy nhất cho mọi bài toán hạt nhân mà không xét điều kiện.” sai vì trái với kiến thức trên.
}
\end{bt}

% ===== Câu 157 =====
% ID: L12C4B25-41-TLN
% Mức: VD
\begin{ex}
% ID: L12C4B25-41-TLN
% Mức: VD
Trong bốn phát biểu sau về bài toán vận dụng cao kết hợp nhiều nội dung của chương iv, có bao nhiêu phát biểu đúng? Chỉ ghi một số nguyên. (1) Có thể dùng một công thức duy nhất cho mọi bài toán hạt nhân mà không xét điều kiện. (2) Cần đọc đúng dữ kiện và giả thiết. (3) Bài vận dụng cao cần tách giai đoạn, dùng đúng bảo toàn, liên hệ khối lượng–năng lượng và định luật phóng xạ, đồng thời kiểm tra đơn vị và ý nghĩa kết quả. (4) Có thể thay số trước khi chọn công thức.
\shortans{2}
\loigiai{
Đối chiếu từng phát biểu với kiến thức cốt lõi: Bài vận dụng cao cần tách giai đoạn, dùng đúng bảo toàn, liên hệ khối lượng–năng lượng và định luật phóng xạ, đồng thời kiểm tra đơn vị và ý nghĩa kết quả. Có 2 phát biểu đúng.
}
\end{ex}

% ===== Câu 158 =====
% ID: L12C4B25-41-TN
% Mức: NB
\dangbt{Tổng hợp định luật phóng xạ, độ phóng xạ và đồ thị}
\begin{ex}
% ID: L12C4B25-41-TN
% Mức: NB
Phát biểu nào sau đây đúng về tổng hợp định luật phóng xạ, độ phóng xạ và đồ thị?
\choice
{\True Các đại lượng $N,m,H$ giảm theo hàm mũ; phải phân biệt lượng còn lại với lượng đã phân rã khi đọc đồ thị.}
{Lượng chất đã phân rã cũng giảm theo cùng biểu thức với lượng chất còn lại.}
{Có thể bỏ qua điều kiện áp dụng và đơn vị trong mọi trường hợp.}
{Kết quả của bài toán không phụ thuộc dữ kiện đã cho.}
\loigiai{
Các đại lượng $N,m,H$ giảm theo hàm mũ; phải phân biệt lượng còn lại với lượng đã phân rã khi đọc đồ thị. Vì vậy phương án $A$ đúng; các phương án còn lại trái với kiến thức hoặc điều kiện áp dụng.
}
\end{ex}

% ===== Câu 159 =====
% ID: L12C4B25-42-TL
% Mức: VDC
\dangbt{Bài toán vận dụng cao kết hợp nhiều nội dung của Chương IV}
\begin{bt}
% ID: L12C4B25-42-TL
% Mức: VDC
So sánh nhận định đúng “Bài vận dụng cao cần tách giai đoạn, dùng đúng bảo toàn, liên hệ khối lượng–năng lượng và định luật phóng xạ, đồng thời kiểm tra đơn vị và ý nghĩa kết quả.” với nhận định sai “Có thể dùng một công thức duy nhất cho mọi bài toán hạt nhân mà không xét điều kiện.”; chỉ rõ căn cứ Vật lí.
\loigiai{
Cần nêu được: Bài vận dụng cao cần tách giai đoạn, dùng đúng bảo toàn, liên hệ khối lượng–năng lượng và định luật phóng xạ, đồng thời kiểm tra đơn vị và ý nghĩa kết quả. Khi vận dụng phải xác định đúng dữ kiện, điều kiện áp dụng, hệ đơn vị và kiểm tra tính hợp lí của kết quả. Nhận định “Có thể dùng một công thức duy nhất cho mọi bài toán hạt nhân mà không xét điều kiện.” sai vì trái với kiến thức trên.
}
\end{bt}

% ===== Câu 160 =====
% ID: L12C4B25-42-TLN
% Mức: VDC
\begin{ex}
% ID: L12C4B25-42-TLN
% Mức: VDC
Trong bốn phát biểu sau về bài toán vận dụng cao kết hợp nhiều nội dung của chương iv, có bao nhiêu phát biểu đúng? Chỉ ghi một số nguyên. (1) Kết quả cần có ý nghĩa vật lí. (2) Có thể dùng một công thức duy nhất cho mọi bài toán hạt nhân mà không xét điều kiện. (3) Phải dùng đúng định luật cho tình huống. (4) Bài vận dụng cao cần tách giai đoạn, dùng đúng bảo toàn, liên hệ khối lượng–năng lượng và định luật phóng xạ, đồng thời kiểm tra đơn vị và ý nghĩa kết quả.
\shortans{3}
\loigiai{
Đối chiếu từng phát biểu với kiến thức cốt lõi: Bài vận dụng cao cần tách giai đoạn, dùng đúng bảo toàn, liên hệ khối lượng–năng lượng và định luật phóng xạ, đồng thời kiểm tra đơn vị và ý nghĩa kết quả. Có 3 phát biểu đúng.
}
\end{ex}

% ===== Câu 161 =====
% ID: L12C4B25-42-TN
% Mức: NB
\dangbt{Tổng hợp định luật phóng xạ, độ phóng xạ và đồ thị}
\begin{ex}
% ID: L12C4B25-42-TN
% Mức: NB
Nhận định nào phù hợp với kiến thức Vật lí về tổng hợp định luật phóng xạ, độ phóng xạ và đồ thị?
\choice
{Lượng chất đã phân rã cũng giảm theo cùng biểu thức với lượng chất còn lại.}
{\True Các đại lượng $N,m,H$ giảm theo hàm mũ; phải phân biệt lượng còn lại với lượng đã phân rã khi đọc đồ thị.}
{Có thể bỏ qua điều kiện áp dụng và đơn vị trong mọi trường hợp.}
{Kết quả của bài toán không phụ thuộc dữ kiện đã cho.}
\loigiai{
Các đại lượng $N,m,H$ giảm theo hàm mũ; phải phân biệt lượng còn lại với lượng đã phân rã khi đọc đồ thị. Vì vậy phương án $B$ đúng; các phương án còn lại trái với kiến thức hoặc điều kiện áp dụng.
}
\end{ex}

% ===== Câu 162 =====
% ID: L12C4B25-43-TN
% Mức: NB
\begin{ex}
% ID: L12C4B25-43-TN
% Mức: NB
Chọn kết luận chính xác về tổng hợp định luật phóng xạ, độ phóng xạ và đồ thị?
\choice
{Lượng chất đã phân rã cũng giảm theo cùng biểu thức với lượng chất còn lại.}
{Có thể bỏ qua điều kiện áp dụng và đơn vị trong mọi trường hợp.}
{\True Các đại lượng $N,m,H$ giảm theo hàm mũ; phải phân biệt lượng còn lại với lượng đã phân rã khi đọc đồ thị.}
{Kết quả của bài toán không phụ thuộc dữ kiện đã cho.}
\loigiai{
Các đại lượng $N,m,H$ giảm theo hàm mũ; phải phân biệt lượng còn lại với lượng đã phân rã khi đọc đồ thị. Vì vậy phương án $C$ đúng; các phương án còn lại trái với kiến thức hoặc điều kiện áp dụng.
}
\end{ex}

% ===== Câu 163 =====
% ID: L12C4B25-44-TN
% Mức: TH
\begin{ex}
% ID: L12C4B25-44-TN
% Mức: TH
Nội dung nào mô tả đúng nhất về tổng hợp định luật phóng xạ, độ phóng xạ và đồ thị?
\choice
{Lượng chất đã phân rã cũng giảm theo cùng biểu thức với lượng chất còn lại.}
{Có thể bỏ qua điều kiện áp dụng và đơn vị trong mọi trường hợp.}
{Kết quả của bài toán không phụ thuộc dữ kiện đã cho.}
{\True Các đại lượng $N,m,H$ giảm theo hàm mũ; phải phân biệt lượng còn lại với lượng đã phân rã khi đọc đồ thị.}
\loigiai{
Các đại lượng $N,m,H$ giảm theo hàm mũ; phải phân biệt lượng còn lại với lượng đã phân rã khi đọc đồ thị. Vì vậy phương án $D$ đúng; các phương án còn lại trái với kiến thức hoặc điều kiện áp dụng.
}
\end{ex}

% ===== Câu 164 =====
% ID: L12C4B25-45-TN
% Mức: TH
\begin{ex}
% ID: L12C4B25-45-TN
% Mức: TH
Khi phân tích, phát biểu nào đúng về tổng hợp định luật phóng xạ, độ phóng xạ và đồ thị?
\choice
{\True Các đại lượng $N,m,H$ giảm theo hàm mũ; phải phân biệt lượng còn lại với lượng đã phân rã khi đọc đồ thị.}
{Lượng chất đã phân rã cũng giảm theo cùng biểu thức với lượng chất còn lại.}
{Có thể bỏ qua điều kiện áp dụng và đơn vị trong mọi trường hợp.}
{Kết quả của bài toán không phụ thuộc dữ kiện đã cho.}
\loigiai{
Các đại lượng $N,m,H$ giảm theo hàm mũ; phải phân biệt lượng còn lại với lượng đã phân rã khi đọc đồ thị. Vì vậy phương án $A$ đúng; các phương án còn lại trái với kiến thức hoặc điều kiện áp dụng.
}
\end{ex}

% ===== Câu 165 =====
% ID: L12C4B25-46-TN
% Mức: TH
\begin{ex}
% ID: L12C4B25-46-TN
% Mức: TH
Kết luận nào của học sinh là đúng về tổng hợp định luật phóng xạ, độ phóng xạ và đồ thị?
\choice
{Lượng chất đã phân rã cũng giảm theo cùng biểu thức với lượng chất còn lại.}
{\True Các đại lượng $N,m,H$ giảm theo hàm mũ; phải phân biệt lượng còn lại với lượng đã phân rã khi đọc đồ thị.}
{Có thể bỏ qua điều kiện áp dụng và đơn vị trong mọi trường hợp.}
{Kết quả của bài toán không phụ thuộc dữ kiện đã cho.}
\loigiai{
Các đại lượng $N,m,H$ giảm theo hàm mũ; phải phân biệt lượng còn lại với lượng đã phân rã khi đọc đồ thị. Vì vậy phương án $B$ đúng; các phương án còn lại trái với kiến thức hoặc điều kiện áp dụng.
}
\end{ex}

% ===== Câu 166 =====
% ID: L12C4B25-47-TN
% Mức: TH
\begin{ex}
% ID: L12C4B25-47-TN
% Mức: TH
Chọn phát biểu không trái với kiến thức về tổng hợp định luật phóng xạ, độ phóng xạ và đồ thị?
\choice
{Lượng chất đã phân rã cũng giảm theo cùng biểu thức với lượng chất còn lại.}
{Có thể bỏ qua điều kiện áp dụng và đơn vị trong mọi trường hợp.}
{\True Các đại lượng $N,m,H$ giảm theo hàm mũ; phải phân biệt lượng còn lại với lượng đã phân rã khi đọc đồ thị.}
{Kết quả của bài toán không phụ thuộc dữ kiện đã cho.}
\loigiai{
Các đại lượng $N,m,H$ giảm theo hàm mũ; phải phân biệt lượng còn lại với lượng đã phân rã khi đọc đồ thị. Vì vậy phương án $C$ đúng; các phương án còn lại trái với kiến thức hoặc điều kiện áp dụng.
}
\end{ex}

% ===== Câu 167 =====
% ID: L12C4B25-48-TN
% Mức: VD
\begin{ex}
% ID: L12C4B25-48-TN
% Mức: VD
Cơ sở Vật lí đúng để giải bài là gì về tổng hợp định luật phóng xạ, độ phóng xạ và đồ thị?
\choice
{Lượng chất đã phân rã cũng giảm theo cùng biểu thức với lượng chất còn lại.}
{Có thể bỏ qua điều kiện áp dụng và đơn vị trong mọi trường hợp.}
{Kết quả của bài toán không phụ thuộc dữ kiện đã cho.}
{\True Các đại lượng $N,m,H$ giảm theo hàm mũ; phải phân biệt lượng còn lại với lượng đã phân rã khi đọc đồ thị.}
\loigiai{
Các đại lượng $N,m,H$ giảm theo hàm mũ; phải phân biệt lượng còn lại với lượng đã phân rã khi đọc đồ thị. Vì vậy phương án $D$ đúng; các phương án còn lại trái với kiến thức hoặc điều kiện áp dụng.
}
\end{ex}

% ===== Câu 168 =====
% ID: L12C4B25-49-TN
% Mức: VD
\begin{ex}
% ID: L12C4B25-49-TN
% Mức: VD
Trong một bài thực hành, nhận xét nào đúng về tổng hợp định luật phóng xạ, độ phóng xạ và đồ thị?
\choice
{\True Các đại lượng $N,m,H$ giảm theo hàm mũ; phải phân biệt lượng còn lại với lượng đã phân rã khi đọc đồ thị.}
{Lượng chất đã phân rã cũng giảm theo cùng biểu thức với lượng chất còn lại.}
{Có thể bỏ qua điều kiện áp dụng và đơn vị trong mọi trường hợp.}
{Kết quả của bài toán không phụ thuộc dữ kiện đã cho.}
\loigiai{
Các đại lượng $N,m,H$ giảm theo hàm mũ; phải phân biệt lượng còn lại với lượng đã phân rã khi đọc đồ thị. Vì vậy phương án $A$ đúng; các phương án còn lại trái với kiến thức hoặc điều kiện áp dụng.
}
\end{ex}

% ===== Câu 169 =====
% ID: L12C4B25-50-TN
% Mức: VD
\begin{ex}
% ID: L12C4B25-50-TN
% Mức: VD
Kết luận đúng khi vận dụng là về tổng hợp định luật phóng xạ, độ phóng xạ và đồ thị?
\choice
{Lượng chất đã phân rã cũng giảm theo cùng biểu thức với lượng chất còn lại.}
{\True Các đại lượng $N,m,H$ giảm theo hàm mũ; phải phân biệt lượng còn lại với lượng đã phân rã khi đọc đồ thị.}
{Có thể bỏ qua điều kiện áp dụng và đơn vị trong mọi trường hợp.}
{Kết quả của bài toán không phụ thuộc dữ kiện đã cho.}
\loigiai{
Các đại lượng $N,m,H$ giảm theo hàm mũ; phải phân biệt lượng còn lại với lượng đã phân rã khi đọc đồ thị. Vì vậy phương án $B$ đúng; các phương án còn lại trái với kiến thức hoặc điều kiện áp dụng.
}
\end{ex}

% ===== Câu 170 =====
% ID: L12C4B25-51-TN
% Mức: NB
\dangbt{Tổng hợp phân hạch, nhiệt hạch và công nghiệp hạt nhân}
\begin{ex}
% ID: L12C4B25-51-TN
% Mức: NB
Phát biểu nào sau đây đúng về tổng hợp phân hạch, nhiệt hạch và công nghiệp hạt nhân?
\choice
{\True Bài toán năng lượng cần xác định số phản ứng, năng lượng mỗi phản ứng và hiệu suất; phần công nghiệp phải xét kiểm soát neutron và an toàn.}
{Hiệu suất không ảnh hưởng điện năng thu được từ nhà máy.}
{Có thể bỏ qua điều kiện áp dụng và đơn vị trong mọi trường hợp.}
{Kết quả của bài toán không phụ thuộc dữ kiện đã cho.}
\loigiai{
Bài toán năng lượng cần xác định số phản ứng, năng lượng mỗi phản ứng và hiệu suất; phần công nghiệp phải xét kiểm soát neutron và an toàn. Vì vậy phương án $A$ đúng; các phương án còn lại trái với kiến thức hoặc điều kiện áp dụng.
}
\end{ex}

% ===== Câu 171 =====
% ID: L12C4B25-52-TN
% Mức: NB
\begin{ex}
% ID: L12C4B25-52-TN
% Mức: NB
Nhận định nào phù hợp với kiến thức Vật lí về tổng hợp phân hạch, nhiệt hạch và công nghiệp hạt nhân?
\choice
{Hiệu suất không ảnh hưởng điện năng thu được từ nhà máy.}
{\True Bài toán năng lượng cần xác định số phản ứng, năng lượng mỗi phản ứng và hiệu suất; phần công nghiệp phải xét kiểm soát neutron và an toàn.}
{Có thể bỏ qua điều kiện áp dụng và đơn vị trong mọi trường hợp.}
{Kết quả của bài toán không phụ thuộc dữ kiện đã cho.}
\loigiai{
Bài toán năng lượng cần xác định số phản ứng, năng lượng mỗi phản ứng và hiệu suất; phần công nghiệp phải xét kiểm soát neutron và an toàn. Vì vậy phương án $B$ đúng; các phương án còn lại trái với kiến thức hoặc điều kiện áp dụng.
}
\end{ex}

% ===== Câu 172 =====
% ID: L12C4B25-53-TN
% Mức: NB
\begin{ex}
% ID: L12C4B25-53-TN
% Mức: NB
Chọn kết luận chính xác về tổng hợp phân hạch, nhiệt hạch và công nghiệp hạt nhân?
\choice
{Hiệu suất không ảnh hưởng điện năng thu được từ nhà máy.}
{Có thể bỏ qua điều kiện áp dụng và đơn vị trong mọi trường hợp.}
{\True Bài toán năng lượng cần xác định số phản ứng, năng lượng mỗi phản ứng và hiệu suất; phần công nghiệp phải xét kiểm soát neutron và an toàn.}
{Kết quả của bài toán không phụ thuộc dữ kiện đã cho.}
\loigiai{
Bài toán năng lượng cần xác định số phản ứng, năng lượng mỗi phản ứng và hiệu suất; phần công nghiệp phải xét kiểm soát neutron và an toàn. Vì vậy phương án $C$ đúng; các phương án còn lại trái với kiến thức hoặc điều kiện áp dụng.
}
\end{ex}

% ===== Câu 173 =====
% ID: L12C4B25-54-TN
% Mức: TH
\begin{ex}
% ID: L12C4B25-54-TN
% Mức: TH
Nội dung nào mô tả đúng nhất về tổng hợp phân hạch, nhiệt hạch và công nghiệp hạt nhân?
\choice
{Hiệu suất không ảnh hưởng điện năng thu được từ nhà máy.}
{Có thể bỏ qua điều kiện áp dụng và đơn vị trong mọi trường hợp.}
{Kết quả của bài toán không phụ thuộc dữ kiện đã cho.}
{\True Bài toán năng lượng cần xác định số phản ứng, năng lượng mỗi phản ứng và hiệu suất; phần công nghiệp phải xét kiểm soát neutron và an toàn.}
\loigiai{
Bài toán năng lượng cần xác định số phản ứng, năng lượng mỗi phản ứng và hiệu suất; phần công nghiệp phải xét kiểm soát neutron và an toàn. Vì vậy phương án $D$ đúng; các phương án còn lại trái với kiến thức hoặc điều kiện áp dụng.
}
\end{ex}

% ===== Câu 174 =====
% ID: L12C4B25-55-TN
% Mức: TH
\begin{ex}
% ID: L12C4B25-55-TN
% Mức: TH
Khi phân tích, phát biểu nào đúng về tổng hợp phân hạch, nhiệt hạch và công nghiệp hạt nhân?
\choice
{\True Bài toán năng lượng cần xác định số phản ứng, năng lượng mỗi phản ứng và hiệu suất; phần công nghiệp phải xét kiểm soát neutron và an toàn.}
{Hiệu suất không ảnh hưởng điện năng thu được từ nhà máy.}
{Có thể bỏ qua điều kiện áp dụng và đơn vị trong mọi trường hợp.}
{Kết quả của bài toán không phụ thuộc dữ kiện đã cho.}
\loigiai{
Bài toán năng lượng cần xác định số phản ứng, năng lượng mỗi phản ứng và hiệu suất; phần công nghiệp phải xét kiểm soát neutron và an toàn. Vì vậy phương án $A$ đúng; các phương án còn lại trái với kiến thức hoặc điều kiện áp dụng.
}
\end{ex}

% ===== Câu 175 =====
% ID: L12C4B25-56-TN
% Mức: TH
\begin{ex}
% ID: L12C4B25-56-TN
% Mức: TH
Kết luận nào của học sinh là đúng về tổng hợp phân hạch, nhiệt hạch và công nghiệp hạt nhân?
\choice
{Hiệu suất không ảnh hưởng điện năng thu được từ nhà máy.}
{\True Bài toán năng lượng cần xác định số phản ứng, năng lượng mỗi phản ứng và hiệu suất; phần công nghiệp phải xét kiểm soát neutron và an toàn.}
{Có thể bỏ qua điều kiện áp dụng và đơn vị trong mọi trường hợp.}
{Kết quả của bài toán không phụ thuộc dữ kiện đã cho.}
\loigiai{
Bài toán năng lượng cần xác định số phản ứng, năng lượng mỗi phản ứng và hiệu suất; phần công nghiệp phải xét kiểm soát neutron và an toàn. Vì vậy phương án $B$ đúng; các phương án còn lại trái với kiến thức hoặc điều kiện áp dụng.
}
\end{ex}

% ===== Câu 176 =====
% ID: L12C4B25-57-TN
% Mức: TH
\begin{ex}
% ID: L12C4B25-57-TN
% Mức: TH
Chọn phát biểu không trái với kiến thức về tổng hợp phân hạch, nhiệt hạch và công nghiệp hạt nhân?
\choice
{Hiệu suất không ảnh hưởng điện năng thu được từ nhà máy.}
{Có thể bỏ qua điều kiện áp dụng và đơn vị trong mọi trường hợp.}
{\True Bài toán năng lượng cần xác định số phản ứng, năng lượng mỗi phản ứng và hiệu suất; phần công nghiệp phải xét kiểm soát neutron và an toàn.}
{Kết quả của bài toán không phụ thuộc dữ kiện đã cho.}
\loigiai{
Bài toán năng lượng cần xác định số phản ứng, năng lượng mỗi phản ứng và hiệu suất; phần công nghiệp phải xét kiểm soát neutron và an toàn. Vì vậy phương án $C$ đúng; các phương án còn lại trái với kiến thức hoặc điều kiện áp dụng.
}
\end{ex}

% ===== Câu 177 =====
% ID: L12C4B25-58-TN
% Mức: VD
\begin{ex}
% ID: L12C4B25-58-TN
% Mức: VD
Cơ sở Vật lí đúng để giải bài là gì về tổng hợp phân hạch, nhiệt hạch và công nghiệp hạt nhân?
\choice
{Hiệu suất không ảnh hưởng điện năng thu được từ nhà máy.}
{Có thể bỏ qua điều kiện áp dụng và đơn vị trong mọi trường hợp.}
{Kết quả của bài toán không phụ thuộc dữ kiện đã cho.}
{\True Bài toán năng lượng cần xác định số phản ứng, năng lượng mỗi phản ứng và hiệu suất; phần công nghiệp phải xét kiểm soát neutron và an toàn.}
\loigiai{
Bài toán năng lượng cần xác định số phản ứng, năng lượng mỗi phản ứng và hiệu suất; phần công nghiệp phải xét kiểm soát neutron và an toàn. Vì vậy phương án $D$ đúng; các phương án còn lại trái với kiến thức hoặc điều kiện áp dụng.
}
\end{ex}

% ===== Câu 178 =====
% ID: L12C4B25-59-TN
% Mức: VD
\begin{ex}
% ID: L12C4B25-59-TN
% Mức: VD
Trong một bài thực hành, nhận xét nào đúng về tổng hợp phân hạch, nhiệt hạch và công nghiệp hạt nhân?
\choice
{\True Bài toán năng lượng cần xác định số phản ứng, năng lượng mỗi phản ứng và hiệu suất; phần công nghiệp phải xét kiểm soát neutron và an toàn.}
{Hiệu suất không ảnh hưởng điện năng thu được từ nhà máy.}
{Có thể bỏ qua điều kiện áp dụng và đơn vị trong mọi trường hợp.}
{Kết quả của bài toán không phụ thuộc dữ kiện đã cho.}
\loigiai{
Bài toán năng lượng cần xác định số phản ứng, năng lượng mỗi phản ứng và hiệu suất; phần công nghiệp phải xét kiểm soát neutron và an toàn. Vì vậy phương án $A$ đúng; các phương án còn lại trái với kiến thức hoặc điều kiện áp dụng.
}
\end{ex}

% ===== Câu 179 =====
% ID: L12C4B25-60-TN
% Mức: VD
\begin{ex}
% ID: L12C4B25-60-TN
% Mức: VD
Kết luận đúng khi vận dụng là về tổng hợp phân hạch, nhiệt hạch và công nghiệp hạt nhân?
\choice
{Hiệu suất không ảnh hưởng điện năng thu được từ nhà máy.}
{\True Bài toán năng lượng cần xác định số phản ứng, năng lượng mỗi phản ứng và hiệu suất; phần công nghiệp phải xét kiểm soát neutron và an toàn.}
{Có thể bỏ qua điều kiện áp dụng và đơn vị trong mọi trường hợp.}
{Kết quả của bài toán không phụ thuộc dữ kiện đã cho.}
\loigiai{
Bài toán năng lượng cần xác định số phản ứng, năng lượng mỗi phản ứng và hiệu suất; phần công nghiệp phải xét kiểm soát neutron và an toàn. Vì vậy phương án $B$ đúng; các phương án còn lại trái với kiến thức hoặc điều kiện áp dụng.
}
\end{ex}

% ===== Câu 180 =====
% ID: L12C4B25-61-TN
% Mức: NB
\dangbt{Bài toán vận dụng cao kết hợp nhiều nội dung của Chương IV}
\begin{ex}
% ID: L12C4B25-61-TN
% Mức: NB
Phát biểu nào sau đây đúng về bài toán vận dụng cao kết hợp nhiều nội dung của chương iv?
\choice
{\True Bài vận dụng cao cần tách giai đoạn, dùng đúng bảo toàn, liên hệ khối lượng–năng lượng và định luật phóng xạ, đồng thời kiểm tra đơn vị và ý nghĩa kết quả.}
{Có thể dùng một công thức duy nhất cho mọi bài toán hạt nhân mà không xét điều kiện.}
{Có thể bỏ qua điều kiện áp dụng và đơn vị trong mọi trường hợp.}
{Kết quả của bài toán không phụ thuộc dữ kiện đã cho.}
\loigiai{
Bài vận dụng cao cần tách giai đoạn, dùng đúng bảo toàn, liên hệ khối lượng–năng lượng và định luật phóng xạ, đồng thời kiểm tra đơn vị và ý nghĩa kết quả. Vì vậy phương án $A$ đúng; các phương án còn lại trái với kiến thức hoặc điều kiện áp dụng.
}
\end{ex}

% ===== Câu 181 =====
% ID: L12C4B25-62-TN
% Mức: NB
\begin{ex}
% ID: L12C4B25-62-TN
% Mức: NB
Nhận định nào phù hợp với kiến thức Vật lí về bài toán vận dụng cao kết hợp nhiều nội dung của chương iv?
\choice
{Có thể dùng một công thức duy nhất cho mọi bài toán hạt nhân mà không xét điều kiện.}
{\True Bài vận dụng cao cần tách giai đoạn, dùng đúng bảo toàn, liên hệ khối lượng–năng lượng và định luật phóng xạ, đồng thời kiểm tra đơn vị và ý nghĩa kết quả.}
{Có thể bỏ qua điều kiện áp dụng và đơn vị trong mọi trường hợp.}
{Kết quả của bài toán không phụ thuộc dữ kiện đã cho.}
\loigiai{
Bài vận dụng cao cần tách giai đoạn, dùng đúng bảo toàn, liên hệ khối lượng–năng lượng và định luật phóng xạ, đồng thời kiểm tra đơn vị và ý nghĩa kết quả. Vì vậy phương án $B$ đúng; các phương án còn lại trái với kiến thức hoặc điều kiện áp dụng.
}
\end{ex}

% ===== Câu 182 =====
% ID: L12C4B25-63-TN
% Mức: NB
\begin{ex}
% ID: L12C4B25-63-TN
% Mức: NB
Chọn kết luận chính xác về bài toán vận dụng cao kết hợp nhiều nội dung của chương iv?
\choice
{Có thể dùng một công thức duy nhất cho mọi bài toán hạt nhân mà không xét điều kiện.}
{Có thể bỏ qua điều kiện áp dụng và đơn vị trong mọi trường hợp.}
{\True Bài vận dụng cao cần tách giai đoạn, dùng đúng bảo toàn, liên hệ khối lượng–năng lượng và định luật phóng xạ, đồng thời kiểm tra đơn vị và ý nghĩa kết quả.}
{Kết quả của bài toán không phụ thuộc dữ kiện đã cho.}
\loigiai{
Bài vận dụng cao cần tách giai đoạn, dùng đúng bảo toàn, liên hệ khối lượng–năng lượng và định luật phóng xạ, đồng thời kiểm tra đơn vị và ý nghĩa kết quả. Vì vậy phương án $C$ đúng; các phương án còn lại trái với kiến thức hoặc điều kiện áp dụng.
}
\end{ex}

% ===== Câu 183 =====
% ID: L12C4B25-64-TN
% Mức: TH
\begin{ex}
% ID: L12C4B25-64-TN
% Mức: TH
Nội dung nào mô tả đúng nhất về bài toán vận dụng cao kết hợp nhiều nội dung của chương iv?
\choice
{Có thể dùng một công thức duy nhất cho mọi bài toán hạt nhân mà không xét điều kiện.}
{Có thể bỏ qua điều kiện áp dụng và đơn vị trong mọi trường hợp.}
{Kết quả của bài toán không phụ thuộc dữ kiện đã cho.}
{\True Bài vận dụng cao cần tách giai đoạn, dùng đúng bảo toàn, liên hệ khối lượng–năng lượng và định luật phóng xạ, đồng thời kiểm tra đơn vị và ý nghĩa kết quả.}
\loigiai{
Bài vận dụng cao cần tách giai đoạn, dùng đúng bảo toàn, liên hệ khối lượng–năng lượng và định luật phóng xạ, đồng thời kiểm tra đơn vị và ý nghĩa kết quả. Vì vậy phương án $D$ đúng; các phương án còn lại trái với kiến thức hoặc điều kiện áp dụng.
}
\end{ex}

% ===== Câu 184 =====
% ID: L12C4B25-65-TN
% Mức: TH
\begin{ex}
% ID: L12C4B25-65-TN
% Mức: TH
Khi phân tích, phát biểu nào đúng về bài toán vận dụng cao kết hợp nhiều nội dung của chương iv?
\choice
{\True Bài vận dụng cao cần tách giai đoạn, dùng đúng bảo toàn, liên hệ khối lượng–năng lượng và định luật phóng xạ, đồng thời kiểm tra đơn vị và ý nghĩa kết quả.}
{Có thể dùng một công thức duy nhất cho mọi bài toán hạt nhân mà không xét điều kiện.}
{Có thể bỏ qua điều kiện áp dụng và đơn vị trong mọi trường hợp.}
{Kết quả của bài toán không phụ thuộc dữ kiện đã cho.}
\loigiai{
Bài vận dụng cao cần tách giai đoạn, dùng đúng bảo toàn, liên hệ khối lượng–năng lượng và định luật phóng xạ, đồng thời kiểm tra đơn vị và ý nghĩa kết quả. Vì vậy phương án $A$ đúng; các phương án còn lại trái với kiến thức hoặc điều kiện áp dụng.
}
\end{ex}

% ===== Câu 185 =====
% ID: L12C4B25-66-TN
% Mức: TH
\begin{ex}
% ID: L12C4B25-66-TN
% Mức: TH
Kết luận nào của học sinh là đúng về bài toán vận dụng cao kết hợp nhiều nội dung của chương iv?
\choice
{Có thể dùng một công thức duy nhất cho mọi bài toán hạt nhân mà không xét điều kiện.}
{\True Bài vận dụng cao cần tách giai đoạn, dùng đúng bảo toàn, liên hệ khối lượng–năng lượng và định luật phóng xạ, đồng thời kiểm tra đơn vị và ý nghĩa kết quả.}
{Có thể bỏ qua điều kiện áp dụng và đơn vị trong mọi trường hợp.}
{Kết quả của bài toán không phụ thuộc dữ kiện đã cho.}
\loigiai{
Bài vận dụng cao cần tách giai đoạn, dùng đúng bảo toàn, liên hệ khối lượng–năng lượng và định luật phóng xạ, đồng thời kiểm tra đơn vị và ý nghĩa kết quả. Vì vậy phương án $B$ đúng; các phương án còn lại trái với kiến thức hoặc điều kiện áp dụng.
}
\end{ex}

% ===== Câu 186 =====
% ID: L12C4B25-67-TN
% Mức: TH
\begin{ex}
% ID: L12C4B25-67-TN
% Mức: TH
Chọn phát biểu không trái với kiến thức về bài toán vận dụng cao kết hợp nhiều nội dung của chương iv?
\choice
{Có thể dùng một công thức duy nhất cho mọi bài toán hạt nhân mà không xét điều kiện.}
{Có thể bỏ qua điều kiện áp dụng và đơn vị trong mọi trường hợp.}
{\True Bài vận dụng cao cần tách giai đoạn, dùng đúng bảo toàn, liên hệ khối lượng–năng lượng và định luật phóng xạ, đồng thời kiểm tra đơn vị và ý nghĩa kết quả.}
{Kết quả của bài toán không phụ thuộc dữ kiện đã cho.}
\loigiai{
Bài vận dụng cao cần tách giai đoạn, dùng đúng bảo toàn, liên hệ khối lượng–năng lượng và định luật phóng xạ, đồng thời kiểm tra đơn vị và ý nghĩa kết quả. Vì vậy phương án $C$ đúng; các phương án còn lại trái với kiến thức hoặc điều kiện áp dụng.
}
\end{ex}

% ===== Câu 187 =====
% ID: L12C4B25-68-TN
% Mức: VD
\begin{ex}
% ID: L12C4B25-68-TN
% Mức: VD
Cơ sở Vật lí đúng để giải bài là gì về bài toán vận dụng cao kết hợp nhiều nội dung của chương iv?
\choice
{Có thể dùng một công thức duy nhất cho mọi bài toán hạt nhân mà không xét điều kiện.}
{Có thể bỏ qua điều kiện áp dụng và đơn vị trong mọi trường hợp.}
{Kết quả của bài toán không phụ thuộc dữ kiện đã cho.}
{\True Bài vận dụng cao cần tách giai đoạn, dùng đúng bảo toàn, liên hệ khối lượng–năng lượng và định luật phóng xạ, đồng thời kiểm tra đơn vị và ý nghĩa kết quả.}
\loigiai{
Bài vận dụng cao cần tách giai đoạn, dùng đúng bảo toàn, liên hệ khối lượng–năng lượng và định luật phóng xạ, đồng thời kiểm tra đơn vị và ý nghĩa kết quả. Vì vậy phương án $D$ đúng; các phương án còn lại trái với kiến thức hoặc điều kiện áp dụng.
}
\end{ex}

% ===== Câu 188 =====
% ID: L12C4B25-69-TN
% Mức: VD
\begin{ex}
% ID: L12C4B25-69-TN
% Mức: VD
Trong một bài thực hành, nhận xét nào đúng về bài toán vận dụng cao kết hợp nhiều nội dung của chương iv?
\choice
{\True Bài vận dụng cao cần tách giai đoạn, dùng đúng bảo toàn, liên hệ khối lượng–năng lượng và định luật phóng xạ, đồng thời kiểm tra đơn vị và ý nghĩa kết quả.}
{Có thể dùng một công thức duy nhất cho mọi bài toán hạt nhân mà không xét điều kiện.}
{Có thể bỏ qua điều kiện áp dụng và đơn vị trong mọi trường hợp.}
{Kết quả của bài toán không phụ thuộc dữ kiện đã cho.}
\loigiai{
Bài vận dụng cao cần tách giai đoạn, dùng đúng bảo toàn, liên hệ khối lượng–năng lượng và định luật phóng xạ, đồng thời kiểm tra đơn vị và ý nghĩa kết quả. Vì vậy phương án $A$ đúng; các phương án còn lại trái với kiến thức hoặc điều kiện áp dụng.
}
\end{ex}

% ===== Câu 189 =====
% ID: L12C4B25-70-TN
% Mức: VD
\begin{ex}
% ID: L12C4B25-70-TN
% Mức: VD
Kết luận đúng khi vận dụng là về bài toán vận dụng cao kết hợp nhiều nội dung của chương iv?
\choice
{Có thể dùng một công thức duy nhất cho mọi bài toán hạt nhân mà không xét điều kiện.}
{\True Bài vận dụng cao cần tách giai đoạn, dùng đúng bảo toàn, liên hệ khối lượng–năng lượng và định luật phóng xạ, đồng thời kiểm tra đơn vị và ý nghĩa kết quả.}
{Có thể bỏ qua điều kiện áp dụng và đơn vị trong mọi trường hợp.}
{Kết quả của bài toán không phụ thuộc dữ kiện đã cho.}
\loigiai{
Bài vận dụng cao cần tách giai đoạn, dùng đúng bảo toàn, liên hệ khối lượng–năng lượng và định luật phóng xạ, đồng thời kiểm tra đơn vị và ý nghĩa kết quả. Vì vậy phương án $B$ đúng; các phương án còn lại trái với kiến thức hoặc điều kiện áp dụng.
}
\end{ex}
